\documentclass[10pt]{article}
\usepackage[utf8]{inputenc}
\usepackage[T1]{fontenc}
\usepackage{graphicx}
\usepackage[export]{adjustbox}
\graphicspath{ {./images/} }
\usepackage{hyperref}
\hypersetup{colorlinks=true, linkcolor=blue, filecolor=magenta, urlcolor=cyan,}
\urlstyle{same}
\usepackage{amsmath}
\usepackage{amsfonts}
\usepackage{amssymb}
\usepackage[version=4]{mhchem}
\usepackage{stmaryrd}
\usepackage{bbold}
\usepackage{mathrsfs}
\usepackage{caption}
\usepackage{multirow}

%New command to display footnote whose markers will always be hidden
\let\svthefootnote\thefootnote
\newcommand\blfootnotetext[1]{%
  \let\thefootnote\relax\footnote{#1}%
  \addtocounter{footnote}{-1}%
  \let\thefootnote\svthefootnote%
}
%Overriding the \footnotetext command to hide the marker if its value is `0`
\let\svfootnotetext\footnotetext
\renewcommand\footnotetext[2][?]{%
  \if\relax#1\relax%
    \ifnum\value{footnote}=0\blfootnotetext{#2}\else\svfootnotetext{#2}\fi%
  \else%
    \if?#1\ifnum\value{footnote}=0\blfootnotetext{#2}\else\svfootnotetext{#2}\fi%
    \else\svfootnotetext[#1]{#2}\fi%
  \fi
}
\DeclareUnicodeCharacter{22C5}{\ifmmode\cdot\else{$\cdot$}\fi}
\DeclareUnicodeCharacter{0394}{\ifmmode\Delta\else{$\Delta$}\fi}
\DeclareUnicodeCharacter{22EE}{\ifmmode\vdots\else{$\vdots$}\fi}
\title{Chapter 10: Unobserved Effects Panel Data Models}

\author{}
\date{}

\begin{document}
\maketitle

\includegraphics[max width=\textwidth, alt={}]{2e17339c-7635-4ef8-aded-2191d8d853ff-0312_111_980_116_283}
\end{center}}
In Chapter 7 we covered a class of linear panel data models where, at a minimum, the error in each time period was assumed to be uncorrelated with the explanatory variables in the same time period. For certain panel data applications this assumption is too strong. In fact, a primary motivation for using panel data is to solve the omitted variables problem.

In this chapter we study population models that explicitly contain a time-constant, unobserved effect. The treatment in this chapter is "modern" in the sense that unobserved effects are treated as random variables, drawn from the population along with the observed explained and explanatory variables, as opposed to parameters to be estimated. In this framework, the key issue is whether the unobserved effect is uncorrelated with the explanatory variables.

\subsection*{10.1 Motivation: Omitted Variables Problem}
It is easy to see how panel data can be used, at least under certain assumptions, to obtain consistent estimators in the presence of omitted variables. Let $y$ and $\mathbf{x} \equiv \left(x_{1}, x_{2}, \ldots, x_{K}\right)$ be observable random variables, and let $c$ be an unobservable random variable; the vector $\left(y, x_{1}, x_{2}, \ldots, x_{K}, c\right)$ represents the population of interest. As is often the case in applied econometrics, we are interested in the partial effects of the observable explanatory variables $x_{j}$ in the population regression function\\
$\mathrm{E}\left(y \mid x_{1}, x_{2}, \ldots, x_{K}, c\right)$.

In words, we would like to hold $c$ constant when obtaining partial effects of the observable explanatory variables. We follow Chamberlain (1984) in using $c$ to denote the unobserved variable. Much of the panel data literature uses a Greek letter, such as $\alpha$ or $\phi$, but we want to emphasize that the unobservable is a random variable, not a parameter to be estimated. (We discuss this point further in Section 10.2.1.)

Assuming a linear model, with $c$ entering additively along with the $x_{j}$, we have\\
$\mathrm{E}(y \mid \mathbf{x}, c)=\beta_{0}+\mathbf{x} \boldsymbol{\beta}+c$,\\
where interest lies in the $K \times 1$ vector $\boldsymbol{\beta}$. On the one hand, if $c$ is uncorrelated with each $x_{j}$, then $c$ is just another unobserved factor affecting $y$ that is not systematically related to the observable explanatory variables whose effects are of interest. On the other hand, if $\operatorname{Cov}\left(x_{j}, c\right) \neq 0$ for some $j$, putting $c$ into the error term can cause serious problems. Without additional information we cannot consistently estimate $\boldsymbol{\beta}$, nor will we be able to determine whether there is a problem (except by introspection, or by concluding that the estimates of $\boldsymbol{\beta}$ are somehow "unreasonable").

Under additional assumptions there are ways to address the problem $\operatorname{Cov}(\mathbf{x}, c) \neq \mathbf{0}$. We have covered at least three possibilities in the context of cross section analysis: (1) we might be able to find a suitable proxy variable for $c$, in which case we can estimate an equation by OLS where the proxy is plugged in for $c$; (2) for the elements of $\mathbf{x}$ that are correlated with $c$, we may be able to find instrumental variables, in which case we can use an IV method such as 2SLS; or (3) we may be able to find indicators of $c$ that can then be used in multiple indicator instrumental variables procedure. These solutions are covered in Chapters 4 and 5.

If we have access to only a single cross section of observations, then the three remedies listed, or slight variants of them, largely exhaust the possibilities. However, if we can observe the same cross section units at different points in time-that is, if we can collect a panel data set-then other possibilties arise.

For illustration, suppose we can observe $y$ and $\mathbf{x}$ at two different time periods; call these $y_{t}, \mathbf{x}_{t}$ for $t=1,2$. The population now represents two time periods on the same unit. Also, suppose that the omitted variable $c$ is time constant. Then we are interested in the population regression function


\begin{equation*}
\mathrm{E}\left(y_{t} \mid \mathbf{x}_{t}, c\right)=\beta_{0}+\mathbf{x}_{t} \boldsymbol{\beta}+c, \quad t=1,2, \tag{10.3}
\end{equation*}


where $\mathbf{x}_{t} \boldsymbol{\beta}=\beta_{1} x_{t 1}+\cdots+\beta_{K} x_{t K}$ and $x_{t j}$ indicates variable $j$ at time $t$. Model (10.3) assumes that $c$ has the same effect on the mean response in each time period. Without loss of generality, we set the coefficient on $c$ equal to one. (Because $c$ is unobserved and virtually never has a natural unit of measurement, it would be meaningless to try to estimate its partial effect.)

The assumption that $c$ is constant over time, and has a constant partial effect over time, is crucial to the following analysis. An unobserved, time-constant variable is called an unobserved effect in panel data analysis. When $t$ represents different time periods for the same individual, the unobserved effect is often interpreted as capturing features of an individual, such as cognitive ability, motivation, or early family upbringing, that are given and do not change over time. Similarly, if the unit of observation is the firm, $c$ contains unobserved firm characteristics-such as managerial quality or structure-that can be viewed as being (roughly) constant over the period in question. We cover several specific examples of unobserved effects models in Section 10.2.

To discuss the additional assumptions sufficient to estimate $\boldsymbol{\beta}$, it is useful to write model (10.3) in error form as


\begin{equation*}
y_{t}=\beta_{0}+\mathbf{x}_{t} \boldsymbol{\beta}+c+u_{t}, \tag{10.4}
\end{equation*}


where, by definition,\\
$\mathrm{E}\left(u_{t} \mid \mathbf{x}_{t}, c\right)=0, \quad t=1,2$.

One implication of condition (10.5) is\\
$\mathrm{E}\left(\mathbf{x}_{t}^{\prime} u_{t}\right)=\mathbf{0}, \quad t=1,2$.

If we were to assume $\mathrm{E}\left(\mathbf{x}_{t}^{\prime} c\right)=\mathbf{0}$, we could apply pooled OLS, as we covered in Section 7.8. If $c$ is correlated with any element of $\mathbf{x}_{t}$, then pooled OLS is biased and inconsistent.

With two years of data we can difference equation (10.4) across the two time periods to eliminate the time-constant unobservable, c. Define $\Delta y=y_{2}-y_{1}, \Delta \mathbf{x}=\mathbf{x}_{2}-\mathbf{x}_{1}$, and $\Delta u=u_{2}-u_{1}$. Then, differencing equation (10.4) gives\\
$\Delta y=\Delta \mathbf{x} \boldsymbol{\beta}+\Delta u$,\\
which is just a standard linear model in the differences of all variables (although the intercept has dropped out). Importantly, the parameter vector of interest, $\boldsymbol{\beta}$, appears directly in equation (10.7), and its presence suggests estimating equation (10.7) by OLS. Given a panel data set with two time periods, equation (10.7) is just a standard cross section equation. Under what assumptions will the OLS estimator from equation (10.7) be consistent?

Because we assume a random sample from the population, we can apply the results in Chapter 4 directly to equation (10.7). The key conditions for OLS to consistently estimate $\boldsymbol{\beta}$ are the orthogonality condition (Assumption OLS.1)


\begin{equation*}
\mathrm{E}\left(\Delta \mathbf{x}^{\prime} \Delta u\right)=\mathbf{0} \tag{10.8}
\end{equation*}


and the rank condition (Assumption OLS.2)\\
$\operatorname{rank} \mathrm{E}\left(\Delta \mathbf{x}^{\prime} \Delta \mathbf{x}\right)=K$.

Consider condition (10.8) first. It is equivalent to $\mathrm{E}\left[\left(\mathbf{x}_{2}-\mathbf{x}_{1}\right)^{\prime}\left(u_{2}-u_{1}\right)\right]=\mathbf{0}$, or, after simple algebra,\\
$\mathrm{E}\left(\mathbf{x}_{2}^{\prime} u_{2}\right)+\mathrm{E}\left(\mathbf{x}_{1}^{\prime} u_{1}\right)-\mathrm{E}\left(\mathbf{x}_{1}^{\prime} u_{2}\right)-\mathrm{E}\left(\mathbf{x}_{2}^{\prime} u_{1}\right)=\mathbf{0}$.

The first two terms in equation (10.10) are zero by condition (10.6), which holds for $t=1,2$. But condition (10.5) does not guarantee that $\mathbf{x}_{1}$ and $u_{2}$ are uncorrelated or that $\mathbf{x}_{2}$ and $u_{1}$ are uncorrelated. It might be reasonable to assume that condition (10.8) holds, but we must recognize that it does not follow from condition (10.5). Assuming that the error $u_{t}$ is uncorrelated with $\mathbf{x}_{1}$ and $\mathbf{x}_{2}$ for $t=1,2$ is an example of a strict exogeneity assumption on the regressors in unobserved components panel data\\
models. We discuss strict exogeneity assumptions generally in Section 10.2. For now, we emphasize that assuming $\operatorname{Cov}\left(\mathbf{x}_{t}, u_{s}\right)=\mathbf{0}$ for all $t$ and $s$ puts no restrictions on the correlation between $\mathbf{x}_{t}$ and the unobserved effect, $c$.

The second assumption, condition (10.9), also deserves some attention now because the elements of $\mathbf{x}_{t}$ appearing in structural equation (10.3) have been differenced across time. If $\mathbf{x}_{t}$ contains a variable that is constant across time for every member of the population, then $\Delta \mathbf{x}$ contains an entry that is identically zero, and condition (10.9) fails. This outcome is not surprising: if $c$ is allowed to be arbitrarily correlated with the elements of $\mathbf{x}_{t}$, the effect of any variable that is constant across time cannot be distinguished from the effect of $c$. Therefore, we can consistently estimate $\beta_{j}$ only if there is some variation in $x_{t j}$ over time.

In the remainder of this chapter, we cover various ways of dealing with the presence of unobserved effects under different sets of assumptions. We assume we have repeated observations on a cross section of $N$ individuals, families, firms, school districts, cities, or some other economic unit. As in Chapter 7, we assume in this chapter that we have the same time periods, denoted $t=1,2, \ldots, T$, for each cross section observation. Such a data set is usually called a balanced panel because the same time periods are available for all cross section units. While the mechanics of the unbalanced case are similar to the balanced case, a careful treatment of the unbalanced case requires a formal description of why the panel may be unbalanced, and the sample selection issues can be somewhat subtle. Therefore, we hold off covering unbalanced panels until Chapter 21, where we discuss sample selection and attrition issues.

We still focus on asymptotic properties of estimators, where the time dimension, $T$, is fixed and the cross section dimension, $N$, grows without bound. With large- $N$ asymptotics it is convenient to view the cross section observations as independent, identically distributed draws from the population. For any cross section observation $i$-denoting a single individual, firm, city, and so on-we denote the observable variables for all $T$ time periods by $\left\{\left(y_{i t}, \mathbf{x}_{i t}\right): t=1,2, \ldots, T\right\}$. Because of the fixed $T$ assumption, the asymptotic analysis is valid for arbitrary time dependence and distributional heterogeneity across $t$.

When applying asymptotic analysis to panel data methods, it is important to remember that asymptotics are useful insofar as they provide a reasonable approximation to the finite sample properties of estimators and statistics. For example, a priori it is difficult to know whether $N \rightarrow \infty$ asymptotics works well with, say, $N=50$ states in the United States and $T=8$ years. But we can be pretty confident that $N \rightarrow \infty$ asymptotics are more appropriate than $T \rightarrow \infty$ asymptotics, even\\
though $N$ is practically fixed while $T$ can grow. With large geographical regions, the random sampling assumption in the cross section dimension is conceptually flawed. Nevertheless, if $N$ is sufficiently large relative to $T$, and if we can assume rough independence in the cross section, then our asymptotic analysis should provide suitable approximations.

If $T$ is of the same order as $N$-for example, $N=60$ countries and $T=55$ postWorld War II years-an asymptotic analysis that makes explicit assumptions about the nature of the time series dependence is needed. (In special cases, the conclusions about consistent estimation and approximate normality of $t$ statistics will be the same, but not generally.) This area is still relatively new. If $T$ is much larger than $N$, say $N=5$ companies and $T=40$ years, the framework becomes multiple time series analysis: $N$ can be held fixed while $T \rightarrow \infty$. We do not cover time series analysis.

\subsection*{10.2 Assumptions about the Unobserved Effects and Explanatory Variables}
Before analyzing panel data estimation methods in more detail, it is useful to generally discuss the nature of the unobserved effects and certain features of the observed explanatory variables.

\subsection*{10.2.1 Random or Fixed Effects?}
The basic unobserved effects model (UEM) can be written, for a randomly drawn cross section observation $i$, as\\
$y_{i t}=\mathbf{x}_{i t} \boldsymbol{\beta}+c_{i}+u_{i t}, \quad t=1,2, \ldots, T$,\\
where $\mathbf{x}_{i t}$ is $1 \times K$ and can contain observable variables that change across $t$ but not $i$, variables that change across $i$ but not $t$, and variables that change across $i$ and $t$. In addition to unobserved effect, there are many other names given to $c_{i}$ in applications: unobserved component, latent variable, and unobserved heterogeneity are common. If $i$ indexes individuals, then $c_{i}$ is sometimes called an individual effect or individual heterogeneity; analogous terms apply to families, firms, cities, and other cross-sectional units. The $u_{i t}$ are called the idiosyncratic errors or idiosyncratic disturbances because these change across $t$ as well as across $i$.

Especially in methodological papers, but also in applications, one often sees a discussion about whether $c_{i}$ will be treated as a random effect or a fixed effect. Originally, such discussions centered on whether $c_{i}$ is properly viewed as a random variable or as a parameter to be estimated. In the traditional approach to panel data models, $c_{i}$ is called a "random effect" when it is treated as a random variable and a "fixed\\
effect" when it is treated as a parameter to be estimated for each cross section observation $i$. Our view is that discussions about whether the $c_{i}$ should be treated as random variables or as parameters to be estimated are wrongheaded for microeconometric panel data applications. With a large number of random draws from the cross section, it almost always makes sense to treat the unobserved effects, $c_{i}$, as random draws from the population, along with $y_{i t}$ and $\mathbf{x}_{i t}$. This approach is certainly appropriate from an omitted variables or neglected heterogeneity perspective. As our discussion in Section 10.1 suggests, the key issue involving $c_{i}$ is whether or not it is uncorrelated with the observed explanatory variables $\mathbf{x}_{i t}, t=1,2, \ldots, T$. Mundlak (1978) made this argument many years ago, and it still is persuasive.

In modern econometric parlance, a random effects framework is synonymous with zero correlation between the observed explanatory variables and the unobserved effect: $\operatorname{Cov}\left(\mathbf{x}_{i t}, c_{i}\right)=\mathbf{0}, t=1,2, \ldots, T$. (Actually, a stronger conditional mean independence assumption, $\mathrm{E}\left(c_{i} \mid \mathbf{x}_{i 1}, \ldots, \mathbf{x}_{i T}\right)=\mathrm{E}\left(c_{i}\right)$, will be needed to fully justify statistical inference; more on this subject in Section 10.4.) In applied papers, when $c_{i}$ is referred to as, say, an "individual random effect," then $c_{i}$ is probably being assumed to be uncorrelated with the $\mathbf{x}_{i t}$.

In most microeconometric applications, a fixed effects framework does not actually mean that $c_{i}$ is being treated as nonrandom; rather, it means that one is allowing for arbitrary dependence between the unobserved effect $c_{i}$ and the observed explanatory variables $\mathbf{x}_{i t}$. So, if $c_{i}$ is called an "individual fixed effect" or a "firm fixed effect," then, for practical purposes, this terminology means that $c_{i}$ is allowed to be correlated arbitrarily with $\mathbf{x}_{i t}$.

More recently, another concept has cropped up for describing situations where we allow dependence between $c_{i}$ and $\left\{\mathbf{x}_{i t}: t=1, \ldots, T\right\}$, especially, but not only, for nonlinear models. (See, for example, Cameron and Trivedi (2005, pp. 719, 786).) If we model the dependence between $c_{i}$ and $\mathbf{x}_{i}=\left(\mathbf{x}_{i 1}, \mathbf{x}_{i 2}, \ldots, \mathbf{x}_{i T}\right)$-or, more generally, place substantive restrictions on the distribution of $c_{i}$ given $\mathbf{x}_{i}$-then we are using a correlated random effects (CRE) framework. (The name refers to allowing correlation between $c_{i}$ and $\mathbf{x}_{i}$.) Practically, the key difference between a fixed effects approach and a correlated random effects approach is that in the former case the relationship between $c_{i}$ and $\mathbf{x}_{i}$ is left entirely unspecified, while in the latter case we restrict this dependence in some way-sometimes rather severely, in the case of nonlinear models (as we will see in Part IV). In Section 10.7.2 we cover Mundlak's (1978) CRE framework in the context of the standard unobserved effects model. Mundlak's approach yields important insights for understanding the difference between the random and fixed effects frameworks, and it is very useful for testing whether $c_{i}$ is uncorrelated with the regressors (the critical assumption in a traditional random\\
effects analysis). Further, Mundlak's approach is indispensable for analyzing a broad class of nonlinear models with unobserved effects, a topic we cover in Part IV.

In this book, we avoid referring to $c_{i}$ as a random effect or a fixed effect because of the unwanted connotations of these terms concerning the nature of $c_{i}$. Instead, we will mostly refer to $c_{i}$ as an "unobserved effect" and sometimes as "unobserved heterogeneity." We will refer to "random effects assumptions," "fixed effects assumptions," and "correlated random effects assumptions." Also, we will label different estimation methods as random effects (RE) estimation, fixed effects (FE) estimation, and, less often for linear models, correlated random effects (CRE) estimation. The terminology for estimation methods is so ingrained in econometrics that it would be counterproductive to try to change it now.

\subsection*{10.2.2 Strict Exogeneity Assumptions on the Explanatory Variables}
Traditional unobserved components panel data models take the $\mathbf{x}_{i t}$ as nonrandom. We will never assume the $\mathbf{x}_{i t}$ are nonrandom because potential feedback from $y_{i t}$ to $\mathbf{x}_{i s}$ for $s>t$ needs to be addressed explicitly.

In Chapter 7 we discussed strict exogeneity assumptions in panel data models that did not explicitly contain unobserved effects. We now provide strict exogeneity assumptions for models with unobserved effects.

In Section 10.1 we stated the strict exogeneity assumption in terms of zero correlation. For inference and efficiency discussions, it is more convenient to state the strict exogeneity assumption in terms of conditional expectations. With an unobserved effect, the most revealing form of the strict exogeneity assumption is


\begin{equation*}
\mathrm{E}\left(y_{i t} \mid \mathbf{x}_{i 1}, \mathbf{x}_{i 2}, \ldots, \mathbf{x}_{i T}, c_{i}\right)=\mathrm{E}\left(y_{i t} \mid \mathbf{x}_{i t}, c_{i}\right)=\mathbf{x}_{i t} \boldsymbol{\beta}+c_{i} \tag{10.12}
\end{equation*}


for $t=1,2, \ldots, T$. The second equality is the functional form assumption on $\mathrm{E}\left(y_{i t} \mid \mathbf{x}_{i t}, c_{i}\right)$. It is the first equality that gives the strict exogeneity its interpretation. It means that, once $\mathbf{x}_{i t}$ and $c_{i}$ are controlled for, $\mathbf{x}_{i s}$ has no partial effect on $y_{i t}$ for $s \neq t$.

When assumption (10.12) holds, we say that the $\left\{\mathbf{x}_{i t}: t=1,2, \ldots, T\right\}$ are strictly exogenous conditional on the unobserved effect $c_{i}$. Assumption (10.12) and the corresponding terminology were introduced and used by Chamberlain (1982). We will explicitly cover Chamberlain's approach to estimating unobserved effects models in the next chapter, but his manner of stating assumptions is instructive even for traditional panel data analysis.

Assumption (10.12) restricts how the expected value of $y_{i t}$ can depend on explanatory variables in other time periods, but it is more reasonable than strict exogeneity without conditioning on the unobserved effect. Without conditioning on an unobserved effect, the strict exogeneity assumption is\\
$\mathrm{E}\left(y_{i t} \mid \mathbf{x}_{i 1}, \mathbf{x}_{i 2}, \ldots, \mathbf{x}_{i T}\right)=\mathrm{E}\left(y_{i t} \mid \mathbf{x}_{i t}\right)=\mathbf{x}_{i t} \boldsymbol{\beta}$,\\
$t=1, \ldots, T$. To see that assumption (10.13) is less likely to hold than assumption (10.12), first consider an example. Suppose that $y_{i t}$ is output of soybeans for farm $i$ during year $t$, and $\mathbf{x}_{i t}$ contains capital, labor, materials (such as fertilizer), rainfall, and other observable inputs. The unobserved effect, $c_{i}$, can capture average quality of land, managerial ability of the family running the farm, and other unobserved, timeconstant factors. A natural assumption is that, once current inputs have been controlled for along with $c_{i}$, inputs used in other years have no effect on output during the current year. However, since the optimal choice of inputs in every year generally depends on $c_{i}$, it is likely that some partial correlation between output in year $t$ and inputs in other years will exist if $c_{i}$ is not controlled for: assumption (10.12) is reasonable while assumption (10.13) is not.

More generally, it is easy to see that assumption (10.13) fails whenever assumption (10.12) holds and the expected value of $c_{i}$ depends on $\left(\mathbf{x}_{i 1}, \ldots, \mathbf{x}_{i T}\right)$. From the law of iterated expectations, if assumption (10.12) holds, then\\
$\mathrm{E}\left(y_{i t} \mid \mathbf{x}_{i 1}, \ldots, \mathbf{x}_{i T}\right)=\mathbf{x}_{i t} \boldsymbol{\beta}+\mathrm{E}\left(c_{i} \mid \mathbf{x}_{i 1}, \ldots, \mathbf{x}_{i T}\right)$,\\
and so assumption (10.13) fails if $\mathrm{E}\left(c_{i} \mid \mathbf{x}_{i 1}, \ldots, \mathbf{x}_{i T}\right) \neq \mathrm{E}\left(c_{i}\right)$. In particular, assumption (10.13) fails if $c_{i}$ is correlated with any of the $\mathbf{x}_{i t}$.

Given equation (10.11), the strict exogeneity assumption can be stated in terms of the idiosyncratic errors as\\
$\mathrm{E}\left(u_{i t} \mid \mathbf{x}_{i 1}, \ldots, \mathbf{x}_{i T}, c_{i}\right)=0, \quad t=1,2, \ldots, T$.

This assumption, in turn, implies that explanatory variables in each time period are uncorrelated with the idiosyncratic error in each time period:\\
$\mathrm{E}\left(\mathbf{x}_{i s}^{\prime} u_{i t}\right)=\mathbf{0}, \quad s, t=1, \ldots, T$.

This assumption is much stronger than assuming zero contemporaneous correlation: $\mathrm{E}\left(\mathbf{x}_{i t}^{\prime} u_{i t}\right)=\mathbf{0}, t=1, \ldots, T$. Nevertheless, assumption (10.15) does allow arbitary correlation between $c_{i}$ and $\mathbf{x}_{i t}$ for all $t$, something we ruled out in Section 7.8. Later, we will use the fact that assumption (10.14) implies that $u_{i t}$ and $c_{i}$ are uncorrelated.

For examining consistency of panel data estimators, the zero covariance assumption (10.15) generally suffices. Further, assumption (10.15) is often the easiest way to think about whether strict exogeneity is likely to hold in a particular application. But standard forms of statistical inference, as well as the efficiency properties of standard estimators, rely on the stronger conditional mean formulation in assumption (10.14). Therefore, we focus on assumption (10.14).

\subsection*{10.2.3 Some Examples of Unobserved Effects Panel Data Models}
Our discussions in Sections 10.2.1 and 10.2.2 emphasize that in any panel data application we should initially focus on two questions: (1) Is the unobserved effect, $c_{i}$, uncorrelated with $\mathbf{x}_{i t}$ for all $t$ ? (2) Is the strict exogeneity assumption (conditional on $c_{i}$ ) reasonable? The following examples illustrate how we might organize our thinking on these two questions.

Example 10.1 (Program Evaluation): A standard model for estimating the effects of job training or other programs on subsequent wages is\\
$\log \left(\right.$ wage $\left._{i t}\right)=\theta_{t}+\mathbf{z}_{i t} \gamma+\delta_{1}$ prog $_{i t}+c_{i}+u_{i t}$,\\
where $i$ indexes individual and $t$ indexes time period. The parameter $\theta_{t}$ denotes a time-varying intercept, and $\mathbf{z}_{i t}$ is a set of observable characteristics that affect wage and may also be correlated with program participation.

Evaluation data sets are often collected at two points in time. At $t=1$, no one has participated in the program, so that prog $_{i 1}=0$ for all $i$. Then, a subgroup is chosen to participate in the program (or the individuals choose to participate), and subsequent wages are observed for the control and treatment groups in $t=2$. Model (10.16) allows for any number of time periods and general patterns of program participation.

The reason for including the individual effect, $c_{i}$, is the usual omitted ability story: if individuals choose whether or not to participate in the program, that choice could be correlated with ability. This possibility is often called the self-selection problem. Alternatively, administrators might assign people based on characteristics that the econometrician cannot observe.

The other issue is the strict exogeneity assumption of the explanatory variables, particularly prog ${ }_{i t}$. Typically, we feel comfortable with assuming that $u_{i t}$ is uncorrelated with prog $_{i t}$. But what about correlation between $u_{i t}$ and, say, prog $_{i, t+1}$ ? Future program participation could depend on $u_{i t}$ if people choose to participate in the future based on shocks to their wage in the past, or if administrators choose people as participants at time $t+1$ who had a low $u_{i t}$. Such feedback might not be very important, since $c_{i}$ is being allowed for, but it could be. See, for example, Bassi (1984) and Ham and Lalonde (1996). Another issue, which is more easily dealt with, is that the training program could have lasting effects. If so, then we should include lags of prog $_{i t}$ in model (10.16). Or, the program itself might last more than one period, in which case prog $_{\text {it }}$ can be replaced by a series of dummy variables for how long unit $i$ at time $t$ has been subject to the program.

Example 10.2 (Distributed Lag Model): Hausman, Hall, and Griliches (1984) estimate nonlinear distributed lag models to study the relationship between patents awarded to a firm and current and past levels of R\&D spending. A linear, five-lag version of their model is\\
patents $_{i t}=\theta_{t}+\mathbf{z}_{i t} \gamma+\delta_{0} R D_{i t}+\delta_{1} R D_{i, t-1}+\cdots+\delta_{5} R D_{i, t-5}+c_{i}+u_{i t}$,\\
where $R D_{i t}$ is spending on $\mathrm{R} \& \mathrm{D}$ for firm $i$ at time $t$ and $\mathbf{z}_{i t}$ contains variables such as firm size (as measured by sales or employees). The variable $c_{i}$ is a firm heterogeneity term that may influence patents ${ }_{i t}$ and that may be correlated with current, past, and future $\mathrm{R} \& \mathrm{D}$ expenditures. Interest lies in the pattern of the $\delta_{j}$ coefficients. As with the other examples, we must decide whether $\mathrm{R} \& \mathrm{D}$ spending is likely to be correlated with $c_{i}$. In addition, if shocks to patents today (changes in $u_{i t}$ ) influence R\&D spending at future dates, then strict exogeneity can fail, and the methods in this chapter will not apply.

The next example presents a case where the strict exogeneity assumption is necessarily false, and the unobserved effect and the explanatory variable must be correlated.

Example 10.3 (Lagged Dependent Variable): A simple dynamic model of wage determination with unobserved heterogeneity is


\begin{equation*}
\log \left(\text { wage }_{i t}\right)=\beta_{1} \log \left(\text { wage }_{i, t-1}\right)+c_{i}+u_{i t}, \quad t=1,2, \ldots, T . \tag{10.18}
\end{equation*}


Often, interest lies in how persistent wages are (as measured by the size of $\beta_{1}$ ) after controlling for unobserved heterogeneity (individual productivity), $c_{i}$. Letting $y_{i t}= \log \left(\right.$ wage $\left._{i t}\right)$, a standard assumption would be\\
$\mathrm{E}\left(u_{i t} \mid y_{i, t-1}, \ldots, y_{i 0}, c_{i}\right)=0$,\\
which means that all of the dynamics are captured by the first lag. Let $x_{i t}=y_{i, t-1}$. Then, under assumption (10.19), $u_{i t}$ is uncorrelated with $\left(x_{i t}, x_{i, t-1}, \ldots, x_{i 1}\right)$, but $u_{i t}$ cannot be uncorrelated with $\left(x_{i, t+1}, \ldots, x_{i T}\right)$, as $x_{i, t+1}=y_{i t}$. In fact,\\
$\mathrm{E}\left(y_{i t} u_{i t}\right)=\beta_{1} \mathrm{E}\left(y_{i, t-1} u_{i t}\right)+\mathrm{E}\left(c_{i} u_{i t}\right)+\mathrm{E}\left(u_{i t}^{2}\right)=\mathrm{E}\left(u_{i t}^{2}\right)>0$,\\
because $\mathrm{E}\left(y_{i, t-1} u_{i t}\right)=0$ and $\mathrm{E}\left(c_{i} u_{i t}\right)=0$ under assumption (10.19). Therefore, the strict exogeneity assumption never holds in unobserved effects models with lagged dependent variables.

In addition, $y_{i, t-1}$ and $c_{i}$ are necessarily correlated (since at time $t-1, y_{i, t-1}$ is the left-hand-side variable). Not only must strict exogeneity fail in this model, but the exogeneity assumption required for pooled OLS estimation of model (10.18) is also violated. We will study estimation of such models in Chapter 11.

\subsection*{10.3 Estimating Unobserved Effects Models by Pooled Ordinary Least Squares}
Under certain assumptions, the pooled OLS estimator can be used to obtain a consistent estimator of $\boldsymbol{\beta}$ in model (10.11). Write the model as\\
$y_{i t}=\mathbf{x}_{i t} \boldsymbol{\beta}+v_{i t}, \quad t=1,2, \ldots, T$,\\
where $v_{i t} \equiv c_{i}+u_{i t}, t=1, \ldots, T$ are the composite errors. For each $t, v_{i t}$ is the sum of the unobserved effect and an idiosyncratic error. From Section 7.8, we know that pooled OLS estimation of this equation is consistent if $\mathrm{E}\left(\mathbf{x}_{i t}^{\prime} v_{i t}\right)=\mathbf{0}, t=1,2, \ldots, T$. Practically speaking, no correlation between $x_{i t}$ and $v_{i t}$ means that we are assuming $\mathrm{E}\left(\mathbf{x}_{i t}^{\prime} u_{i t}\right)=\mathbf{0}$ and\\
$\mathrm{E}\left(\mathbf{x}_{i t}^{\prime} c_{i}\right)=\mathbf{0}, \quad t=1,2, \ldots, T$.

Equation (10.22) is the more restrictive assumption, since $\mathrm{E}\left(\mathbf{x}_{i t}^{\prime} u_{i t}\right)=\mathbf{0}$ holds if we have successfully modeled $\mathrm{E}\left(y_{i t} \mid \mathbf{x}_{i t}, c_{i}\right)$.

In static and finite distributed lag models we are sometimes willing to make the assumption (10.22); in fact, we will do so in the next section on random effects estimation. As seen in Example 10.3, models with lagged dependent variables in $\mathbf{x}_{i t}$ must violate assumption (10.22) because $y_{i, t-1}$ and $c_{i}$ must be correlated.

Even if assumption (10.22) holds, the composite errors will be serially correlated due to the presence of $c_{i}$ in each time period. Therefore, inference using pooled OLS requires the robust variance matrix estimator and robust test statistics from Chapter 7. Because $v_{i t}$ depends on $c_{i}$ for all $t$, the correlation between $v_{i t}$ and $v_{i s}$ does not generally decrease as the distance $|t-s|$ increases; in time series parlance, the $v_{i t}$ are not weakly dependent across time. (We show this property explicitly in the next section when $\left\{u_{i t}: t=1, \ldots, T\right\}$ is homoskedastic and serially uncorrelated.) Therefore, it is important that we be able to rely on large- $N$ and fixed- $T$ asymptotics when applying pooled OLS.

As we discussed in Chapter 7, each $\left(\mathbf{y}_{i}, \mathbf{X}_{i}\right)$ has $T$ rows and should be ordered chronologically, and the $\left(\mathbf{y}_{i}, \mathbf{X}_{i}\right)$ should be stacked from $i=1, \ldots, N$. The order of the cross section observations is, as usual, irrelevant.

\subsection*{10.4 Random Effects Methods}
\subsection*{10.4.1 Estimation and Inference under the Basic Random Effects Assumptions}
As with pooled OLS, a random effects analysis puts $c_{i}$ into the error term. In fact, random effects analysis imposes more assumptions than those needed for pooled

OLS: strict exogeneity in addition to orthogonality between $c_{i}$ and $\mathbf{x}_{i t}$. Stating the assumption in terms of conditional means, we have\\
assumption RE.1:\\
(a) $\mathrm{E}\left(u_{i t} \mid \mathbf{x}_{i}, c_{i}\right)=0, t=1, \ldots, T$; (b) $\mathrm{E}\left(c_{i} \mid \mathbf{x}_{i}\right)=\mathrm{E}\left(c_{i}\right)=0$\\
where $\mathbf{x}_{i} \equiv\left(\mathbf{x}_{i 1}, \mathbf{x}_{i 2}, \ldots, \mathbf{x}_{i T}\right)$.\\
In Section 10.2 we discussed the meaning of the strict exogeneity Assumption RE.1a. Assumption RE.1b is how we will state the orthogonality between $c_{i}$ and each $\mathbf{x}_{i t}$. For obtaining consistency results, we could relax RE.1b to assumption (10.22), but in practice (10.22) affords little more generality, and we will use Assumption RE.1b later to derive the traditional asymptotic variance for the random effects estimator. Assumption RE.1b is always implied by the assumption that the $\mathbf{x}_{i t}$ are nonrandom and $\mathrm{E}\left(c_{i}\right)=0$, or by the assumption that $c_{i}$ is independent of $\mathbf{x}_{i}$. The important part is $\mathrm{E}\left(c_{i} \mid \mathbf{x}_{i}\right)=\mathrm{E}\left(c_{i}\right)$; the assumption $\mathrm{E}\left(c_{i}\right)=0$ is without loss of generality, provided an intercept is included in $\mathbf{x}_{i t}$, as should always be the case.

Why do we maintain Assumption RE. 1 when it is more restrictive than needed for a pooled OLS analysis? The random effects approach exploits the serial correlation in the composite error, $v_{i t}=c_{i}+u_{i t}$, in a generalized least squares (GLS) framework. In order to ensure that feasible GLS is consistent, we need some form of strict exogeneity between the explanatory variables and the composite error. Under Assumption RE. 1 we can write\\
$y_{i t}=\mathbf{x}_{i t} \boldsymbol{\beta}+v_{i t}$,\\
$\mathrm{E}\left(v_{i t} \mid \mathbf{x}_{i}\right)=0, \quad t=1,2, \ldots, T$,\\
where\\
$v_{i t}=c_{i}+u_{i t}$.

Equation (10.24) shows that $\left\{\mathbf{x}_{i t}: t=1, \ldots, T\right\}$ satisfies the strict exogeneity assumption SGLS. 1 (see Chapter 7) in the model (10.23). Therefore, we can apply GLS methods that account for the particular error structure in equation (10.25).

Write the model (10.23) for all $T$ time periods as


\begin{equation*}
\mathbf{y}_{i}=\mathbf{X}_{i} \boldsymbol{\beta}+\mathbf{v}_{i} \tag{10.26}
\end{equation*}


and $\mathbf{v}_{i}$ can be written as $\mathbf{v}_{i}=c_{i} \mathbf{j}_{T}+\mathbf{u}_{i}$, where $\mathbf{j}_{T}$ is the $T \times 1$ vector of ones. Define the (unconditional) variance matrix of $\mathbf{v}_{i}$ as\\
$\boldsymbol{\Omega} \equiv \mathrm{E}\left(\mathbf{v}_{i} \mathbf{v}_{i}^{\prime}\right)$,\\
a $T \times T$ matrix that we assume to be positive definite. Remember, this matrix is necessarily the same for all $i$ because of the random sampling assumption in the cross section.

For consistency of GLS, we need the usual rank condition for GLS:\\
assumption RE.2: rank $\mathrm{E}\left(\mathbf{X}_{i}^{\prime} \boldsymbol{\Omega}^{-1} \mathbf{X}_{i}\right)=K$.\\
Applying the results from Chapter 7, we know that GLS and feasible GLS are consistent under Assumptions RE. 1 and RE.2. A general FGLS analysis, using an unrestricted variance estimator $\boldsymbol{\Omega}$, is consistent and $\sqrt{N}$-asymptotically normal as $N \rightarrow \infty$. But we would not be exploiting the unobserved effects structure of $v_{i t}$. A standard random effects analysis adds assumptions on the idiosyncratic errors that give $\boldsymbol{\Omega}$ a special form. The first assumption is that the idiosyncratic errors $u_{i t}$ have a constant unconditional variance across $t$ :\\
$\mathrm{E}\left(u_{i t}^{2}\right)=\sigma_{u}^{2}, \quad t=1,2, \ldots, T$.

The second assumption is that the idiosyncratic errors are serially uncorrelated:\\
$\mathrm{E}\left(u_{i t} u_{i s}\right)=0, \quad$ all $t \neq s$.

Under these two assumptions, we can derive the variances and covariances of the elements of $\mathbf{v}_{i}$. Under Assumption RE.1a, $\mathrm{E}\left(c_{i} u_{i t}\right)=0, t=1,2, \ldots, T$, and so\\
$\mathrm{E}\left(v_{i t}^{2}\right)=\mathrm{E}\left(c_{i}^{2}\right)+2 \mathrm{E}\left(c_{i} u_{i t}\right)+\mathrm{E}\left(u_{i t}^{2}\right)=\sigma_{c}^{2}+\sigma_{u}^{2}$,\\
where $\sigma_{c}^{2}=\mathrm{E}\left(c_{i}^{2}\right)$. Also, for all $t \neq s$,\\
$\mathrm{E}\left(v_{i t} v_{i s}\right)=\mathrm{E}\left[\left(c_{i}+u_{i t}\right)\left(c_{i}+u_{i s}\right)\right]=\mathrm{E}\left(c_{i}^{2}\right)=\sigma_{c}^{2}$.\\
Therefore, under assumptions RE.1, (10.28), and (10.29), $\boldsymbol{\Omega}$ takes the special form\\
$\boldsymbol{\Omega}=\mathrm{E}\left(\mathbf{v}_{i} \mathbf{v}_{i}^{\prime}\right)=\left(\begin{array}{cccc}\sigma_{c}^{2}+\sigma_{u}^{2} & \sigma_{c}^{2} & \cdots & \sigma_{c}^{2} \\ \sigma_{c}^{2} & \sigma_{c}^{2}+\sigma_{u}^{2} & \cdots & \vdots \\ \vdots & & \ddots & \sigma_{c}^{2} \\ \sigma_{c}^{2} & & & \sigma_{c}^{2}+\sigma_{u}^{2}\end{array}\right)$.

Because $\mathbf{j}_{T} \mathbf{j}_{T}^{\prime}$ is the $T \times T$ matrix with unity in every element, we can write the matrix (10.30) as\\
$\boldsymbol{\Omega}=\sigma_{u}^{2} \mathbf{I}_{T}+\sigma_{c}^{2} \mathbf{j}_{T} \mathbf{j}_{T}^{\prime}$.

When $\boldsymbol{\Omega}$ has the form (10.31), we say it has the random effects structure. Rather than depending on $T(T+1) / 2$ unrestricted variances and covariances, as would be the case in a general GLS analysis, $\boldsymbol{\Omega}$ depends only on two parameters, $\sigma_{c}^{2}$ and $\sigma_{u}^{2}$, regardless of the size of $T$. The correlation between the composite errors $v_{i t}$ and $v_{i s}$ does not depend on the difference between $t$ and $s: \operatorname{Corr}\left(v_{i s}, v_{i t}\right)=\sigma_{c}^{2} /\left(\sigma_{c}^{2}+\sigma_{u}^{2}\right) \geq 0$, $s \neq t$. This correlation is also the ratio of the variance of $c_{i}$ to the variance of the composite error, and it is useful as a measure of the relative importance of the unobserved effect $c_{i}$.

Unlike the stable AR(1) model we briefly discussed in Section 7.8.6, the correlation between the composite errors $v_{i t}$ and $v_{i s}$ does not tend to zero as $t$ and $s$ get far apart under the RE covariance structure. If $v_{i t}=\rho v_{i, t-1}+e_{i t}$, where $|\rho|<1$ and the $\left\{e_{i t}\right\}$ are serially uncorrelated, then $\operatorname{Cov}\left(v_{i t}, v_{i s}\right)$ converges to zero pretty quickly as $|t-s|$ gets large. (Of course, the convergence is faster with smaller $|\rho|$.) Unlike standard models for serial correlation in time series settings, the random effects assumption implies strong persistence in the unobservables over time, due, of course, to the presence of $c_{i}$.

Assumptions (10.28) and (10.29) are special to random effects. For efficiency of feasible GLS, we assume that the variance matrix of $\mathbf{v}_{i}$ conditional on $\mathbf{x}_{i}$ is constant:\\
$\mathrm{E}\left(\mathbf{v}_{i} \mathbf{v}_{i}^{\prime} \mid \mathbf{x}_{i}\right)=\mathrm{E}\left(\mathbf{v}_{i} \mathbf{v}_{i}^{\prime}\right)$.

Assumptions (10.28), (10.29), and (10.32) are implied by our third RE assumption:\\
assumption RE.3: (a) $\mathrm{E}\left(\mathbf{u}_{i} \mathbf{u}_{i}^{\prime} \mid \mathbf{x}_{i}, c_{i}\right)=\sigma_{u}^{2} \mathbf{I}_{T}$; (b) $\mathrm{E}\left(c_{i}^{2} \mid \mathbf{x}_{i}\right)=\sigma_{c}^{2}$.\\
Under Assumption RE.3a, $\mathrm{E}\left(u_{i t}^{2} \mid \mathbf{x}_{i}, c_{i}\right)=\sigma_{u}^{2}, t=1, \ldots, T$, which implies assumption (10.28), and $\mathrm{E}\left(u_{i t} u_{i s} \mid \mathbf{x}_{i}, c_{i}\right)=0, t \neq s, t, s=1, \ldots, T$, which implies assumption (10.29) (both by the usual iterated expectations argument). But Assumption RE.3a is stronger because it assumes that the conditional variances are constant and the conditional covariances are zero. Along with Assumption RE.1b, Assumption RE.3b is the same as $\operatorname{Var}\left(c_{i} \mid \mathbf{x}_{i}\right)=\operatorname{Var}\left(c_{i}\right)$, which is a homoskedasticity assumption on the unobserved effect $c_{i}$. Under Assumption RE.3, assumption (10.32) holds and $\boldsymbol{\Omega}$ has the form (10.30).

To implement an FGLS procedure, define $\sigma_{v}^{2}=\sigma_{c}^{2}+\sigma_{u}^{2}$. For now, assume that we have consistent estimators of $\sigma_{u}^{2}$ and $\sigma_{c}^{2}$. Then we can form\\
$\hat{\boldsymbol{\Omega}} \equiv \hat{\sigma}_{u}^{2} \mathbf{I}_{T}+\hat{\sigma}_{c}^{2} \mathbf{j}_{T} \mathbf{j}_{T}^{\prime}$,\\
a $T \times T$ matrix that we assume to be positive definite. In a panel data context, the FGLS estimator that uses the variance matrix (10.33) is what is known as the random effects estimator.


\begin{equation*}
\hat{\boldsymbol{\beta}}_{R E}=\left(\sum_{i=1}^{N} \mathbf{X}_{i}^{\prime} \hat{\boldsymbol{\Omega}}^{-1} \mathbf{X}_{i}\right)^{-1}\left(\sum_{i=1}^{N} \mathbf{X}_{i}^{\prime} \hat{\boldsymbol{\Omega}}^{-1} \mathbf{y}_{i}\right) . \tag{10.34}
\end{equation*}


Before we discuss obtaining $\hat{\boldsymbol{\Omega}}$ and performing asymptotic inference under the full set of RE assumptions, it is important to know that the RE estimator is generally consistent with or without Assumption RE.3. Clearly, the form of $\hat{\boldsymbol{\Omega}}$ is motivated by Assumption RE.3. Nevertheless, as we discussed in Section 7.5.3, incorrectly imposing restrictions on $\mathrm{E}\left(\mathbf{v}_{i} \mathbf{v}_{i}^{\prime}\right)$ does not cause inconsistency of FGLS provided the appropriate strict exogeneity assumption holds, as it does under Assumption RE.1. Of course, we have to insert $\operatorname{plim}(\hat{\boldsymbol{\Omega}})$ (which necessarily has the RE form) in place of $\boldsymbol{\Omega}$ in Assumption RE.2, but that is a minor change. Practically important is that without Assumption RE.3, we need to make inference fully robust, but we hold off on that until Section 10.4.2.

Under Assumption RE.3, the RE estimator is efficient in the class of estimators consistent under $\mathrm{E}\left(\mathbf{v}_{i} \mid \mathbf{x}_{i}\right)=\mathbf{0}$, including pooled OLS and a variety of weighted least squares estimators, because RE is asymptotically equivalent to GLS under Assumptions RE.1-RE.3. The usual feasible GLS variance matrix-see equation (7.54)-is valid under Assumptions RE.1-RE.3. The only difference from the general analysis is that $\hat{\boldsymbol{\Omega}}$ is chosen as in expression (10.33).

In order to implement the RE procedure, we need to obtain $\hat{\sigma}_{c}^{2}$ and $\hat{\sigma}_{u}^{2}$. Actually, it is easiest to first find $\hat{\sigma}_{v}^{2}=\hat{\sigma}_{c}^{2}+\hat{\sigma}_{u}^{2}$. Under Assumption RE.3a, $\sigma_{v}^{2}=T^{-1} \sum_{t=1}^{T} \mathrm{E}\left(v_{i t}^{2}\right)$ for all $i$; therefore, averaging $v_{i t}^{2}$ across all $i$ and $t$ would give a consistent estimator of $\sigma_{v}^{2}$. But we need to estimate $\boldsymbol{\beta}$ to make this method operational. A convenient initial estimator of $\boldsymbol{\beta}$ is the pooled OLS estimator, denoted here by $\check{\boldsymbol{\beta}}$. Let $\check{v}_{i t}$ denote the pooled OLS residuals. A consistent estimator of $\sigma_{v}^{2}$ is


\begin{equation*}
\hat{\sigma}_{v}^{2}=\frac{1}{(N T-K)} \sum_{i=1}^{N} \sum_{t=1}^{T} \check{v}_{i t}^{2}, \tag{10.35}
\end{equation*}


which is the usual variance estimator from the OLS regression on the pooled data. The degrees-of-freedom correction in equation (10.35)-that is, the use of $N T-K$ rather than $N T$-has no effect asymptotically. Under Assumptions RE.1-RE.3, equation (10.35) is a consistent estimator (actually, a $\sqrt{N}$-consistent estimator) of $\sigma_{v}^{2}$.

To find a consistent estimator of $\sigma_{c}^{2}$, recall that $\sigma_{c}^{2}=\mathrm{E}\left(v_{i t} v_{i s}\right)$, all $t \neq s$. Therefore, for each $i$, there are $T(T-1) / 2$ nonredundant error products that can be used to estimate $\sigma_{c}^{2}$. If we sum all these combinations and take the expectation, we get, for each $i$,


\begin{align*}
\mathrm{E}\left(\sum_{t=1}^{T-1} \sum_{s=t+1}^{T} v_{i t} v_{i s}\right) & =\sum_{t=1}^{T-1} \sum_{s=t+1}^{T} \mathrm{E}\left(v_{i t} v_{i s}\right)=\sum_{t=1}^{T-1} \sum_{s=t+1}^{T} \sigma_{c}^{2}=\sigma_{c}^{2} \sum_{t=1}^{T-1}(T-t) \\
& =\sigma_{c}^{2}((T-1)+(T-2)+\cdots+2+1)=\sigma_{c}^{2} T(T-1) / 2 \tag{10.36}
\end{align*}


where we have used the fact that the sum of the first $T-1$ positive integers is $T(T-1) / 2$. As usual, a consistent estimator is obtained by replacing the expectation with an average (across $i$ ) and replacing $v_{i t}$ with its pooled OLS residual. We also make a degrees-of-freedom adjustment as a small-sample correction:


\begin{equation*}
\hat{\sigma}_{c}^{2}=\frac{1}{[N T(T-1) / 2-K]} \sum_{i=1}^{N} \sum_{t=1}^{T-1} \sum_{s=t+1}^{T} \check{v}_{i t} \check{v}_{i s} \tag{10.37}
\end{equation*}


is a $\sqrt{N}$-consistent estimator of $\sigma_{c}^{2}$ under Assumptions RE.1-RE.3. Equation (10.37), without the degrees-of-freedom adjustment, appears elsewhere-for example, Baltagi (2001, Sect. 2.3). Given $\hat{\sigma}_{v}^{2}$ and $\hat{\sigma}_{c}^{2}$, we can form $\hat{\sigma}_{u}^{2}=\hat{\sigma}_{v}^{2}-\hat{\sigma}_{c}^{2}$. (The idiosyncratic error variance, $\sigma_{u}^{2}$, can also be estimated using the fixed effects method, which we discuss in Section 10.5. Also, there are other methods of estimating $\sigma_{c}^{2}$. A common estimator of $\sigma_{c}^{2}$ is based on the between estimator of $\boldsymbol{\beta}$, which we touch on in Section 10.5 ; see Hsiao (2003, Sect. 3.3) and Baltagi (2001, Sect. 2.3). Because the RE estimator is a feasible GLS estimator under strict exogeneity, all that we need are consistent estimators of $\sigma_{c}^{2}$ and $\sigma_{u}^{2}$ in order to obtain a $\sqrt{N}$-efficient estimator of $\boldsymbol{\beta}$.)

As a practical matter, equation (10.37) is not guaranteed to be positive, although it is in the vast majority of applications. A negative value for $\hat{\sigma}_{c}^{2}$ is indicative of negative serial correlation in $u_{i t}$, probably a substantial amount, which means that Assumption RE.3a is violated. Alternatively, some other assumption in the model can be false. We should make sure that time dummies are included in the model if aggregate effects are important; omitting them can induce serial correlation in the implied $u_{i t}$. When the intercepts are allowed to change freely over time, the effects of other aggregate variables will not be identified. If $\hat{\sigma}_{c}^{2}$ is negative, unrestricted FGLS may be called for; see Section 10.4.3.

Example 10.4 (RE Estimation of the Effects of Job Training Grants): We now use the data in JTRAIN1.RAW to estimate the effect of job training grants on firm scrap rates, using a random effects analysis. There are 54 firms that reported scrap rates for each of the years 1987, 1988, and 1989. Grants were not awarded in 1987. Some firms received grants in 1988, others received grants in 1989, and a firm could not receive a grant twice. Since there are firms in 1989 that received a grant only in 1988, it is important to allow the grant effect to persist one period. The estimated equation is

$$
\begin{aligned}
\widehat{\log \widehat{(s c r a p)}}= & \underset{(.243)}{.415}-\underset{(.109)}{.093} d 88-\underset{(.132)}{.270} d 89+\underset{(.411)}{.548} \text { union } \\
& -\underset{(.148)}{.215} \text { grant }-\underset{(.205)}{.377} \text { grant }_{-1} .
\end{aligned}
$$

The lagged value of grant has the larger impact and is statistically significant at the 5 percent level against a one-sided alternative. You are invited to estimate the equation without grant $_{-1}$ to verify that the estimated grant effect is much smaller (on the order of 6.7 percent) and statistically insignificant.

Multiple hypotheses tests are carried out as in any FGLS analysis; see Section 7.6, where $G=T$. In computing an $F$-type statistic based on weighted sums of squared residuals, $\hat{\boldsymbol{\Omega}}$ in expression (10.33) should be based on the pooled OLS residuals from the unrestricted model. Then, obtain the residuals from the unrestricted random effects estimation as $\hat{\mathbf{v}}_{i} \equiv \mathbf{y}_{i}-\mathbf{X}_{i} \hat{\boldsymbol{\beta}}_{R E}$. Let $\tilde{\boldsymbol{\beta}}_{R E}$ denote the REs estimator with the $Q$ linear restrictions imposed, and define the restricted RE residuals as $\tilde{\mathbf{v}}_{i} \equiv \mathbf{y}_{i}-\mathbf{X}_{i} \tilde{\boldsymbol{\beta}}_{R E}$. Insert these into equation (7.56) in place of $\hat{\mathbf{u}}_{i}$ and $\tilde{\mathbf{u}}_{i}$ for a chi-square statistic or into equation (7.57) for an $F$-type statistic.

In Example 10.4, the Wald test for joint significance of grant and grant $_{-1}$ (against a two-sided alternative) yields a $\chi_{2}^{2}$ statistic equal to 3.66 , with $p$-value $=.16$. (This test comes from Stata.)

\subsection*{10.4.2 Robust Variance Matrix Estimator}
Because failure of Assumption RE. 3 does not cause inconsistency in the RE estimator, it is very useful to be able to conduct statistical inference without this assumption. Assumption RE. 3 can fail for two reasons. First, $\mathrm{E}\left(\mathbf{v}_{i} \mathbf{v}_{i}^{\prime} \mid \mathbf{x}_{i}\right)$ may not be constant, so that $\mathrm{E}\left(\mathbf{v}_{i} \mathbf{v}_{i}^{\prime} \mid \mathbf{x}_{i}\right) \neq \mathrm{E}\left(\mathbf{v}_{i} \mathbf{v}_{i}^{\prime}\right)$. This outcome is always a possibility with GLS analysis. Second, $\mathrm{E}\left(\mathbf{v}_{i} \mathbf{v}_{i}^{\prime}\right)$ may not have the RE structure: the idiosyncratic errors $u_{i t}$ may have variances that change over time, or they could be serially correlated. In either case a robust variance matrix is available from the analysis in Chapter 7. We simply use equation (7.52) with $\hat{\mathbf{u}}_{i}$ replaced by $\hat{\mathbf{v}}_{i}=\mathbf{y}_{i}-\mathbf{X}_{i} \hat{\boldsymbol{\beta}}_{R E}, i=1,2, \ldots, N$, the $T \times 1$ vectors of RE residuals.

Robust standard errors are obtained in the usual way from the robust variance matrix estimator, and robust Wald statistics are obtained by the usual formula $W= (\mathbf{R} \hat{\boldsymbol{\beta}}-\mathbf{r})^{\prime}\left(\mathbf{R} \hat{\mathbf{V}} \mathbf{R}^{\prime}\right)^{-1}(\mathbf{R} \hat{\boldsymbol{\beta}}-\mathbf{r})$, where $\hat{\mathbf{V}}$ is the robust variance matrix estimator. Remember, if Assumption RE. 3 is violated, the sum of squared residuals form of the $F$ statistic is not valid.

The idea behind using a robust variance matrix is the following. Assumptions RE.1-RE. 3 lead to a well-known estimation technique whose properties are understood under these assumptions. But it is always a good idea to make the analysis robust whenever feasible. With fixed $T$ and large $N$ asymptotics, we lose nothing in using the robust standard errors and test statistics even if Assumption RE. 3 holds. In Section 10.7.2, we show how the RE estimator can be obtained from a particular pooled OLS regression, which makes obtaining robust standard errors and $t$ and $F$ statistics especially easy.

\subsection*{10.4.3 General Feasible Generalized Least Squares Analysis}
If the idiosyncratic errors $\left\{u_{i t}: t=1,2, \ldots, T\right\}$ are generally heteroskedastic and serially correlated across $t$, a more general estimator of $\boldsymbol{\Omega}$ can be used in FGLS:\\
$\hat{\boldsymbol{\Omega}}=N^{-1} \sum_{i=1}^{N} \check{\mathbf{v}}_{i} \check{\mathbf{v}}_{i}^{\prime}$,\\
where the $\check{\mathbf{v}}_{i}$ would be the pooled OLS residuals. The FGLS estimator is consistent under Assumptions RE. 1 and RE.2, and, if we assume that $\mathrm{E}\left(\mathbf{v}_{i} \mathbf{v}_{i}^{\prime} \mid \mathbf{x}_{i}\right)=\boldsymbol{\Omega}$, then the FGLS estimator is asymptotically efficient and its asymptotic variance estimator takes the usual form.

Using equation (10.38) is more general than the RE analysis. In fact, with large $N$ asymptotics, the general FGLS estimator is just as efficient as the RE estimator under Assumptions RE.1-RE.3. Using equation (10.38) is asymptotically more efficient if $\mathrm{E}\left(\mathbf{v}_{i} \mathbf{v}_{i}^{\prime} \mid \mathbf{x}_{i}\right)=\boldsymbol{\Omega}$, but $\boldsymbol{\Omega}$ does not have the RE form. So why not always use FGLS with $\hat{\boldsymbol{\Omega}}$ given in equation (10.38)? There are historical reasons for using RE methods rather than a general FGLS analysis. The structure of $\boldsymbol{\Omega}$ in the matrix (10.30) was once synonomous with unobserved effects models: any correlation in the composite errors $\left\{v_{i t}: t=1,2, \ldots, T\right\}$ was assumed to be caused by the presence of $c_{i}$. The idiosyncratic errors, $u_{i t}$, were, by definition, taken to be serially uncorrelated and homoskedastic.

If $N$ is not several times larger than $T$, an unrestricted FGLS analysis can have poor finite sample properties because $\hat{\boldsymbol{\Omega}}$ has $T(T+1) / 2$ estimated elements. Even though estimation of $\boldsymbol{\Omega}$ does not affect the asymptotic distribution of the FGLS estimator, it certainly affects its finite sample properties. Random effects estimation requires estimation of only two variance parameters for any $T$.

With very large $N$, using the general estimate of $\boldsymbol{\Omega}$ is an attractive alternative, especially if the estimate in equation (10.38) appears to have a pattern different from the RE pattern. As a middle ground between a traditional random effects analysis\\
and a full-blown FGLS analysis, we might specify a particular structure for the idiosyncratic error variance matrix $\mathrm{E}\left(\mathbf{u}_{i} \mathbf{u}_{i}^{\prime}\right)$. For example, if $\left\{u_{i t}\right\}$ follows a stable firstorder autoregressive process with autocorrelation coefficient $\rho$ and variance $\sigma_{u}^{2}$, then $\boldsymbol{\Omega}=\mathrm{E}\left(\mathbf{u}_{i} \mathbf{u}_{i}^{\prime}\right)+\sigma_{c}^{2} \mathbf{j}_{T} \mathbf{j}_{T}^{\prime}$ depends in a known way on only three parameters, $\sigma_{u}^{2}, \sigma_{c}^{2}$, and $\rho$. These parameters can be estimated after initial pooled OLS estimation, and then an FGLS procedure using the particular structure of $\boldsymbol{\Omega}$ is easy to implement. We do not cover such possibilities explicitly; see, for example, MaCurdy (1982). Some are preprogrammed in statistical packages, but not necessarily with the option of making inference robust to misspecification of $\boldsymbol{\Omega}$ or system heteroskedasticity.

\subsection*{10.4.4 Testing for the Presence of an Unobserved Effect}
If the standard random effects assumptions RE.1-RE. 3 hold but the model does not actually contain an unobserved effect, pooled OLS is efficient and all associated pooled OLS statistics are asymptotically valid. The absence of an unobserved effect is statistically equivalent to $\mathrm{H}_{0}: \sigma_{c}^{2}=0$.

To test $\mathrm{H}_{0}: \sigma_{c}^{2}=0$, we can use the simple test for $\mathrm{AR}(1)$ serial correlation covered in Chapter 7 (see equation (7.81)). The $\operatorname{AR}(1)$ test is valid because the errors $v_{i t}$ are serially uncorrelated under the null $\mathrm{H}_{0}: \sigma_{c}^{2}=0$ (and we are assuming that $\left\{\mathbf{x}_{i t}\right\}$ is strictly exogenous). However, a better test is based directly on the estimator of $\sigma_{c}^{2}$ in equation (10.37).

Breusch and Pagan (1980) derive a statistic using the Lagrange multiplier principle in a likelihood setting (something we cover in Chapter 13). We will not derive the Breusch and Pagan statistic because we are not assuming any particular distribution for the $v_{i t}$. Instead, we derive a similar test that has the advantage of being valid for any distribution of $\mathbf{v}_{i}$ and only states that the $v_{i t}$ are uncorrelated under the null. (In particular, the statistic is valid for heteroskedasticity in the $v_{i t}$.)

From equation (10.37), we base a test of $\mathrm{H}_{0}: \sigma_{c}^{2}=0$ on the null asymptotic distribution of


\begin{equation*}
N^{-1 / 2} \sum_{i=1}^{N} \sum_{t=1}^{T-1} \sum_{s=t+1}^{T} \hat{v}_{i t} \hat{v}_{i s} \tag{10.39}
\end{equation*}


which is essentially the estimator $\hat{\sigma}_{c}^{2}$ scaled up by $\sqrt{N}$. Because of strict exogeneity, this statistic has the same limiting distribution (as $N \rightarrow \infty$ with fixed $T$ ) when we replace the pooled OLS residuals $\hat{v}_{i t}$ with the errors $v_{i t}$ (see Problem 7.4). For any distribution of the $v_{i t}, N^{-1 / 2} \sum_{i=1}^{N} \sum_{t=1}^{T-1} \sum_{s=t+1}^{T} v_{i t} v_{i s}$ has a limiting normal distribution (under the null that the $v_{i t}$ are serially uncorrelated) with variance $\mathrm{E}\left(\sum_{t=1}^{T-1} \sum_{s=t+1}^{T} v_{i t} v_{i s}\right)^{2}$. We can estimate this variance in the usual way (take away\\
the expectation, average across $i$, and replace $v_{i t}$ with $\hat{v}_{i t}$ ). When we put expression (10.39) over its asymptotic standard error we get the statistic


\begin{equation*}
\frac{\sum_{i=1}^{N} \sum_{t=1}^{T-1} \sum_{s=t+1}^{T} \hat{v}_{i t} \hat{v}_{i s}}{\left[\sum_{i=1}^{N}\left(\sum_{t=1}^{T-1} \sum_{s=t+1}^{T} \hat{v}_{i t} \hat{v}_{i s}\right)^{2}\right]^{1 / 2}} \tag{10.40}
\end{equation*}


Under the null hypothesis that the $v_{i t}$ are serially uncorrelated, this statistic is distributed asymptotically as standard normal. Unlike the Breusch-Pagan statistic, with expression (10.40) we can reject $\mathrm{H}_{0}$ for negative estimates of $\sigma_{c}^{2}$, although negative estimates are rare in practice (unless we have already differenced the data, something we discuss in Section 10.6).

The statistic in expression (10.40) can detect many kinds of serial correlation in the composite error $v_{i t}$, and so a rejection of the null should not be interpreted as implying that the RE error structure must be true. Finding that the $v_{i t}$ are serially uncorrelated is not very surprising in applications, especially since $\mathbf{x}_{i t}$ cannot contain lagged dependent variables for the methods in this chapter.

It is probably more interesting to test for serial correlation in the $\left\{u_{i t}\right\}$, as this is a test of the RE form of $\boldsymbol{\Omega}$. Baltagi and Li (1995) obtain a test under normality of $c_{i}$ and $\left\{u_{i t}\right\}$, based on the Lagrange multiplier principle. In Section 10.7.2, we discuss a simpler test for serial correlation in $\left\{u_{i t}\right\}$ using a pooled OLS regression on transformed data, which does not rely on normality.

\subsection*{10.5 Fixed Effects Methods}
\subsection*{10.5.1 Consistency of the Fixed Effects Estimator}
Again consider the linear unobserved effects model for $T$ time periods:


\begin{equation*}
y_{i t}=\mathbf{x}_{i t} \boldsymbol{\beta}+c_{i}+u_{i t}, \quad t=1, \ldots, T . \tag{10.41}
\end{equation*}


The RE approach to estimating $\boldsymbol{\beta}$ effectively puts $c_{i}$ into the error term, under the assumption that $c_{i}$ is orthogonal to $\mathbf{x}_{i t}$, and then accounts for the implied serial correlation in the composite error $v_{i t}=c_{i}+u_{i t}$ using a GLS analysis. In many applications the whole point of using panel data is to allow for $c_{i}$ to be arbitrarily correlated with the $\mathbf{x}_{i t}$. A fixed effects analysis achieves this purpose explicitly.

The $T$ equations in the model (10.41) can be written as


\begin{equation*}
\mathbf{y}_{i}=\mathbf{X}_{i} \boldsymbol{\beta}+c_{i} \mathbf{j}_{T}+\mathbf{u}_{i}, \tag{10.42}
\end{equation*}


where $\mathbf{j}_{T}$ is still the $T \times 1$ vector of ones. As usual, equation (10.42) represents a single random draw from the cross section.

The first fixed effects (FE) assumption is strict exogeneity of the explanatory variables conditional on $c_{i}$ :\\
assumption FE.1: $\quad \mathrm{E}\left(u_{i t} \mid \mathbf{x}_{i}, c_{i}\right)=0, t=1,2, \ldots, T$.\\
This assumption is identical to the first part of Assumption RE.1. Thus, we maintain strict exogeneity of $\left\{\mathbf{x}_{i t}: t=1, \ldots, T\right\}$ conditional on the unobserved effect. The key difference is that we do not assume RE.1b. In other words, for FE analysis, $\mathrm{E}\left(c_{i} \mid \mathbf{x}_{i}\right)$ is allowed to be any function of $\mathbf{x}_{i}$.

By relaxing RE.1b we can consistently estimate partial effects in the presence of time-constant omitted variables that can be arbitrarily related to the observables $\mathbf{x}_{i t}$. Therefore, FE analysis is more robust than random effects analysis. As we suggested in Section 10.1, this robustness comes at a price: without further assumptions, we cannot include time-constant factors in $\mathbf{x}_{i t}$. The reason is simple: if $c_{i}$ can be arbitrarily correlated with each element of $\mathbf{x}_{i t}$, there is no way to distinguish the effects of time-constant observables from the time-constant unobservable $c_{i}$. When analyzing individuals, factors such as gender or race cannot be included in $\mathbf{x}_{i t}$. For analyzing firms, industry cannot be included in $\mathbf{x}_{i t}$ unless industry designation changes over time for at least some firms. For cities, variables describing fixed city attributes, such as whether or not the city is near a river, cannot be included in $\mathbf{x}_{i t}$.

The fact that $\mathbf{x}_{i t}$ cannot include time-constant explanatory variables is a drawback in certain applications, but when the interest is only on time-varying explanatory variables, it is convenient not to have to worry about modeling time-constant factors that are not of direct interest.

In panel data analysis the term "time-varying explanatory variables" means that each element of $\mathbf{x}_{i t}$ varies over time for some cross section units. Often there are elements of $\mathbf{x}_{i t}$ that are constant across time for a subset of the cross section. For example, if we have a panel of adults and one element of $\mathbf{x}_{i t}$ is education, we can allow education to be constant for some part of the sample. But we must have education changing for some people in the sample.

As a general specification, let $d 2_{t}, \ldots, d T_{t}$ denote time period dummies so that $d s_{t}=1$ if $s=t$, and zero otherwise (often these are defined in terms of specific years, such as $d 88_{t}$, but at this level we call them time period dummies). Let $\mathbf{z}_{i}$ be a vector of time-constant observables, and let $\mathbf{w}_{i t}$ be a vector of time-varying variables. Suppose $y_{i t}$ is determined by


\begin{align*}
& y_{i t}=\theta_{1}+\theta_{2} d 2_{t}+\cdots+\theta_{T} d T_{t}+\mathbf{z}_{i} \gamma_{1}+d 2_{t} \mathbf{z}_{i} \gamma_{2} \\
& \quad+\cdots+d T_{t} \mathbf{z}_{i} \gamma_{T}+\mathbf{w}_{i t} \boldsymbol{\delta}+c_{i}+u_{i t}  \tag{10.43}\\
& \quad \mathrm{E}\left(u_{i t} \mid \mathbf{z}_{i}, \mathbf{w}_{i 1}, \mathbf{w}_{i 2}, \ldots, \mathbf{w}_{i T}, c_{i}\right)=0, \quad t=1,2, \ldots, T . \tag{10.44}
\end{align*}


We hope that this model represents a causal relationship, where the conditioning on $c_{i}$ allows us to control for unobserved factors that are time constant. Without further assumptions, the intercept $\theta_{1}$ cannot be identified and the vector $\gamma_{1}$ on $\mathbf{z}_{i}$ cannot be identified, because $\theta_{1}+\mathbf{z}_{i} \gamma_{1}$ cannot be distinguished from $c_{i}$. Note that $\theta_{1}$ is the intercept for the base time period, $t=1$, and $\gamma_{1}$ measures the effects of $\mathbf{z}_{i}$ on $y_{i t}$ in period $t=1$. Even though we cannot identify the effects of the $\mathbf{z}_{i}$ in any particular time period, $\gamma_{2}, \gamma_{3}, \ldots, \gamma_{T}$ are identified, and therefore we can estimate the differences in the partial effects on time-constant variables relative to a base period. In particular, we can test whether the effects of time-constant variables have changed over time. As a specific example, if $y_{i t}=\log \left(\right.$ wage $\left._{i t}\right)$ and one element of $\mathbf{z}_{i}$ is a female binary variable, then we can estimate how the gender gap has changed over time, even though we cannot estimate the gap in any particular time period.

The idea for estimating $\boldsymbol{\beta}$ under Assumption FE. 1 is to transform the equations to eliminate the unobserved effect $c_{i}$. When at least two time periods are available, there are several transformations that accomplish this purpose. In this section we study the fixed effects transformation, also called the within transformation. The FE transformation is obtained by first averaging equation (10.41) over $t=1, \ldots, T$ to get the cross section equation\\
$\bar{y}_{i}=\overline{\mathbf{x}}_{i} \boldsymbol{\beta}+c_{i}+\bar{u}_{i}$,\\
where $\bar{y}_{i}=T^{-1} \sum_{t=1}^{T} y_{i t}, \quad \overline{\mathbf{x}}_{i}=T^{-1} \sum_{t=1}^{T} \mathbf{x}_{i t}$, and $\bar{u}_{i}=T^{-1} \sum_{t=1}^{T} u_{i t}$. Subtracting equation (10.45) from equation (10.41) for each $t$ gives the FE transformed equation,\\
$y_{i t}-\bar{y}_{i}=\left(\mathbf{x}_{i t}-\overline{\mathbf{x}}_{i}\right) \boldsymbol{\beta}+u_{i t}-\bar{u}_{i}$\\
or\\
$\ddot{y}_{i t}=\ddot{\mathbf{x}}_{i t} \boldsymbol{\beta}+\ddot{u}_{i t}, \quad t=1,2, \ldots, T$,\\
where $\ddot{y}_{i t} \equiv y_{i t}-\bar{y}_{i}, \ddot{\mathbf{x}}_{i t} \equiv \mathbf{x}_{i t}-\overline{\mathbf{x}}_{i}$, and $\ddot{u}_{i t} \equiv u_{i t}-\bar{u}_{i}$. The time demeaning of the original equation has removed the individual specific effect $c_{i}$.

With $c_{i}$ out of the picture, it is natural to think of estimating equation (10.46) by pooled OLS. Before investigating this possibility, we must remember that equation (10.46) is an estimating equation: the interpretation of $\boldsymbol{\beta}$ comes from the (structural) conditional expectation $\mathrm{E}\left(y_{i t} \mid \mathbf{x}_{i}, c_{i}\right)=\mathrm{E}\left(y_{i t} \mid \mathbf{x}_{i t}, c_{i}\right)=\mathbf{x}_{i t} \boldsymbol{\beta}+c_{i}$.

To see whether pooled OLS estimation of equation (10.46) will be consistent, we need to show that the key pooled OLS assumption (Assumption POLS. 1 from Chapter 7) holds in equation (10.46). That is,\\
$\mathrm{E}\left(\ddot{\mathbf{x}}_{i t}^{\prime} \ddot{u}_{i t}\right)=\mathbf{0}, \quad t=1,2, \ldots, T$.

For each $t$, the left-hand side of equation (10.47) can be written as $\mathrm{E}\left[\left(\mathbf{x}_{i t}-\overline{\mathbf{x}}_{i}\right)^{\prime}\left(u_{i t}-\bar{u}_{i}\right)\right]$. Now, under Assumption FE.1, $u_{i t}$ is uncorrelated with $\mathbf{x}_{i s}$, for all $s, t=1,2, \ldots, T$. It follows that $u_{i t}$ and $\bar{u}_{i}$ are uncorrelated with $\mathbf{x}_{i t}$ and $\overline{\mathbf{x}}_{i}$ for $t=1,2, \ldots, T$. Therefore, assumption (10.47) holds under Assumption FE.1, and so pooled OLS applied to equation (10.46) can be expected to produce consistent estimators. We can actually say a lot more than condition (10.47): under Assumption FE.1, $\mathrm{E}\left(\ddot{u}_{i t} \mid \mathbf{x}_{i}\right)=\mathrm{E}\left(u_{i t} \mid \mathbf{x}_{i}\right)-\mathrm{E}\left(\bar{u}_{i} \mid \mathbf{x}_{i}\right)=0$, which in turn implies that $\mathrm{E}\left(\ddot{u}_{i t} \mid \ddot{\mathbf{x}}_{i 1}, \ldots, \ddot{\mathbf{x}}_{i T}\right)=0$, since each $\ddot{\mathbf{x}}_{i t}$ is just a function of $\mathbf{x}_{i}=\left(\mathbf{x}_{i 1}, \ldots, \mathbf{x}_{i T}\right)$. This result shows that the $\ddot{\mathbf{x}}_{i t}$ satisfy the conditional expectation form of the strict exogeneity assumption in the model (10.46). Among other things, this conclusion implies that the FE estimator of $\boldsymbol{\beta}$ that we will derive is actually unbiased under Assumption FE.1.

It is important to see that assumption (10.47) fails if we try to relax the strict exogeneity assumption to something weaker, such as $\mathrm{E}\left(\mathbf{x}_{i t}^{\prime} u_{i t}\right)=\mathbf{0}$, all $t$, because this assumption does not ensure that $\mathbf{x}_{i s}$ is uncorrelated with $u_{i t}, s \neq t$. In Chapter 11 we study violations of the strict exogeneity assumption in more detail.

The FE estimator, denoted by $\hat{\boldsymbol{\beta}}_{F E}$, is the pooled OLS estimator from the regression\\
$\ddot{y}_{i t}$ on $\ddot{\mathbf{x}}_{i t}, \quad t=1,2, \ldots, T ; i=1,2, \ldots, N$.

The FE estimator is simple to compute once the time demeaning has been carried out. Some econometrics packages have special commands to carry out FE estimation (and commands to carry out the time demeaning for all $i$ ). It is also fairly easy to program this estimator in matrix-oriented languages.

To study the FE estimator a little more closely, write equation (10.46) for all time periods as\\
$\ddot{\mathbf{y}}_{i}=\ddot{\mathbf{X}}_{i} \boldsymbol{\beta}+\ddot{\mathbf{u}}_{i}$,\\
where $\ddot{\mathbf{y}}_{i}$ is $T \times 1, \ddot{\mathbf{X}}_{i}$ is $T \times K$, and $\ddot{\mathbf{u}}_{i}$ is $T \times 1$. This set of equations can be obtained by premultiplying equation (10.42) by a time-demeaning matrix. Define $\mathbf{Q}_{T} \equiv \mathbf{I}_{T}- \mathbf{j}_{T}\left(\mathbf{j}_{T}^{\prime} \mathbf{j}_{T}\right)^{-1} \mathbf{j}_{T}^{\prime}$, which is easily seen to be a $T \times T$ symmetric, idempotent matrix with rank $T-1$. Further, $\mathbf{Q}_{T} \mathbf{j}_{T}=\mathbf{0}, \mathbf{Q}_{T} \mathbf{y}_{i}=\ddot{\mathbf{y}}_{i}, \mathbf{Q}_{T} \mathbf{X}_{i}=\ddot{\mathbf{X}}_{i}$, and $\mathbf{Q}_{T} \mathbf{u}_{i}=\ddot{\mathbf{u}_{i}}$, and so premultiplying equation (10.42) by $\mathbf{Q}_{T}$ gives the demeaned equations (10.49).

In order to ensure that the FE estimator is well behaved asymptotically, we need a standard rank condition on the matrix of time-demeaned explanatory variables:\\
assumption FE.2: $\quad \operatorname{rank}\left(\sum_{t=1}^{T} \mathrm{E}\left(\ddot{\mathbf{x}}_{i t}^{\prime} \ddot{\mathbf{x}}_{i t}\right)\right)=\operatorname{rank}\left[\mathrm{E}\left(\ddot{\mathbf{X}}_{i}^{\prime} \ddot{\mathbf{X}}_{i}\right)\right]=K$.\\
If $\mathbf{x}_{i t}$ contains an element that does not vary over time for any $i$, then the corresponding element in $\ddot{\mathbf{x}}_{i t}$ is identically zero for all $t$ and any draw from the cross section. Since $\ddot{\mathbf{X}}_{i}$ would contain a column of zeros for all $i$, Assumption FE. 2 could not be true. Assumption FE. 2 shows explicitly why time-constant variables are not allowed in fixed effects analysis (unless they are interacted with time-varying variables, such as time dummies).

The FE estimator can be expressed as


\begin{equation*}
\hat{\boldsymbol{\beta}}_{F E}=\left(\sum_{i=1}^{N} \ddot{\mathbf{X}}_{i}^{\prime} \ddot{\mathbf{X}}_{i}\right)^{-1}\left(\sum_{i=1}^{N} \ddot{\mathbf{X}}_{i}^{\prime} \ddot{\mathbf{y}}_{i}\right)=\left(\sum_{i=1}^{N} \sum_{t=1}^{T} \ddot{\mathbf{x}}_{i t}^{\prime} \ddot{\mathbf{x}}_{i t}\right)^{-1}\left(\sum_{i=1}^{N} \sum_{t=1}^{T} \ddot{\mathbf{x}}_{i t}^{\prime} \ddot{y}_{i t}\right) . \tag{10.50}
\end{equation*}


It is also called the within estimator because it uses the time variation within each cross section. The between estimator, which uses only variation between the cross section observations, is the OLS estimator applied to the time-averaged equation (10.45). This estimator is not consistent under Assumption FE. 1 because $\mathrm{E}\left(\overline{\mathbf{x}}_{i}^{\prime} c_{i}\right)$ is not necessarily zero. The between estimator is consistent under Assumption RE. 1 and a standard rank condition, but it effectively discards the time series information in the data set. It is more efficient to use the RE estimator.

Under Assumption FE. 1 and the finite sample version of Assumption FE.2, namely, $\operatorname{rank}\left(\ddot{\mathbf{X}^{\prime}} \ddot{\mathbf{X}}\right)=K, \hat{\boldsymbol{\beta}}_{F E}$ can be shown to be unbiased conditional on $\mathbf{X}$.

\subsection*{10.5.2 Asymptotic Inference with Fixed Effects}
Without further assumptions the FE estimator is not necessarily the most efficient estimator based on Assumption FE.1. The next assumption ensures that FE is efficient.\\
assumption FE.3: $\quad \mathrm{E}\left(\mathbf{u}_{i} \mathbf{u}_{i}^{\prime} \mid \mathbf{x}_{i}, c_{i}\right)=\sigma_{u}^{2} \mathbf{I}_{T}$.\\
Assumption FE. 3 is identical to Assumption RE.3a. Since $\mathrm{E}\left(\mathbf{u}_{i} \mid \mathbf{x}_{i}, c_{i}\right)=\mathbf{0}$ by Assumption FE.1, Assumption FE. 3 is the same as saying $\operatorname{Var}\left(\mathbf{u}_{i} \mid \mathbf{x}_{i}, c_{i}\right)=\sigma_{u}^{2} \mathbf{I}_{T}$ if Assumption FE. 1 also holds. As with Assumption RE.3a, it is useful to think of Assumption FE. 3 as having two parts. The first is that $\mathrm{E}\left(\mathbf{u}_{i} \mathbf{u}_{i}^{\prime} \mid \mathbf{x}_{i}, c_{i}\right)=\mathrm{E}\left(\mathbf{u}_{i} \mathbf{u}_{i}^{\prime}\right)$, which is standard in system estimation contexts (see equation (7.53)). The second is that the unconditional variance matrix $\mathrm{E}\left(\mathbf{u}_{i} \mathbf{u}_{i}^{\prime}\right)$ has the special form $\sigma_{u}^{2} \mathbf{I}_{T}$. This implies that the idiosyncratic errors $u_{i t}$ have a constant variance across $t$ and are serially uncorrelated, just as in assumptions (10.28) and (10.29).

Assumption FE.3, along with Assumption FE.1, implies that the unconditional variance matrix of the composite error $\mathbf{v}_{i}=c_{i} \mathbf{j}_{T}+\mathbf{u}_{i}$ has the RE form. However, without Assumption RE.3b, $\mathrm{E}\left(\mathbf{v}_{i} \mathbf{v}_{i}^{\prime} \mid \mathbf{x}_{i}\right) \neq \mathrm{E}\left(\mathbf{v}_{i} \mathbf{v}_{i}^{\prime}\right)$. While this result matters for inference with the RE estimator, it has no bearing on a fixed effects analysis.

It is not obvious that Assumption FE. 3 has the desired consequences of ensuring efficiency of FE and leading to simple computation of standard errors and test statistics. Consider the demeaned equation (10.46). Normally, for pooled OLS to be relatively efficient, we require that the $\left\{\ddot{u}_{i t}: t=1,2, \ldots, T\right\}$ be homoskedastic across $t$ and serially uncorrelated. The variance of $\ddot{u}_{i t}$ can be computed as


\begin{align*}
\mathrm{E}\left(\ddot{u}_{i t}^{2}\right) & =\mathrm{E}\left[\left(u_{i t}-\bar{u}_{i}\right)^{2}\right]=\mathrm{E}\left(u_{i t}^{2}\right)+\mathrm{E}\left(\bar{u}_{i}^{2}\right)-2 \mathrm{E}\left(u_{i t} \bar{u}_{i}\right) \\
& =\sigma_{u}^{2}+\sigma_{u}^{2} / T-2 \sigma_{u}^{2} / T=\sigma_{u}^{2}(1-1 / T) \tag{10.51}
\end{align*}


which verifies (unconditional) homoskedasticity across $t$. However, for $t \neq s$, the covariance between $\ddot{u}_{i t}$ and $\ddot{u}_{i s}$ is

$$
\begin{aligned}
\mathrm{E}\left(\ddot{u}_{i t} \ddot{u}_{i s}\right) & =\mathrm{E}\left[\left(u_{i t}-\bar{u}_{i}\right)\left(u_{i s}-\bar{u}_{i}\right)\right]=\mathrm{E}\left(u_{i t} u_{i s}\right)-\mathrm{E}\left(u_{i t} \bar{u}_{i}\right)-\mathrm{E}\left(u_{i s} \bar{u}_{i}\right)+\mathrm{E}\left(\bar{u}_{i}^{2}\right) \\
& =0-\sigma_{u}^{2} / T-\sigma_{u}^{2} / T+\sigma_{u}^{2} / T=-\sigma_{u}^{2} / T<0
\end{aligned}
$$

Combining this expression with the variance in equation (10.51) gives, for all $t \neq s$,


\begin{equation*}
\operatorname{Corr}\left(\ddot{u}_{i t}, \ddot{u}_{i s}\right)=-1 /(T-1) \tag{10.52}
\end{equation*}


which shows that the time-demeaned errors $\ddot{u}_{t i}$ are negatively serially correlated. (As $T$ gets large, the correlation tends to zero.)

It turns out that, because of the nature of time demeaning, the serial correlation in the $\ddot{u}_{i t}$ under Assumption FE. 3 causes only minor complications. To find the asymptotic variance of $\hat{\boldsymbol{\beta}}_{F E}$, write

$$
\sqrt{N}\left(\hat{\boldsymbol{\beta}}_{F E}-\boldsymbol{\beta}\right)=\left(N^{-1} \sum_{i=1}^{N} \ddot{\mathbf{X}}_{i}^{\prime} \ddot{\mathbf{X}}_{i}\right)^{-1}\left(N^{-1 / 2} \sum_{i=1}^{N} \ddot{\mathbf{X}}_{i}^{\prime} \mathbf{u}_{i}\right)
$$

where we have used the important fact that $\ddot{\mathbf{X}}_{i}^{\prime} \ddot{\mathbf{u}}_{i}=\mathbf{X}_{i}^{\prime} \mathbf{Q}_{T} \mathbf{u}_{i}=\ddot{\mathbf{X}}_{i}^{\prime} \mathbf{u}_{i}$. Under Assumption FE.3, $\mathrm{E}\left(\mathbf{u}_{i} \mathbf{u}_{i}^{\prime} \mid \ddot{\mathbf{X}}_{i}\right)=\sigma_{u}^{2} \mathbf{I}_{T}$. From the system OLS analysis in Chapter 7 it follows that\\
$\sqrt{N}\left(\hat{\boldsymbol{\beta}}_{F E}-\boldsymbol{\beta}\right) \sim \operatorname{Normal}\left(\mathbf{0}, \sigma_{u}^{2}\left[\mathrm{E}\left(\ddot{\mathbf{X}}_{i}^{\prime} \ddot{\mathbf{X}}_{i}\right)\right]^{-1}\right)$,\\
and so\\
$\operatorname{Avar}\left(\hat{\boldsymbol{\beta}}_{F E}\right)=\sigma_{u}^{2}\left[\mathrm{E}\left(\ddot{\mathbf{X}}_{i}^{\prime} \ddot{\mathbf{X}}_{i}\right)\right]^{-1} / N$.

Given a consistent estimator $\hat{\sigma}_{u}^{2}$ of $\sigma_{u}^{2}$, equation (10.53) is easily estimated by also replacing $\mathrm{E}\left(\ddot{\mathbf{X}}_{i}^{\prime} \ddot{\mathbf{X}}_{i}\right)$ with its sample analogue $N^{-1} \sum_{i=1}^{N} \ddot{\mathbf{X}}_{i}^{\prime} \ddot{\mathbf{X}}_{i}$ :\\
$\left.\widehat{\operatorname{Avar}\left(\hat{\boldsymbol{\beta}}_{F E}\right.}\right)=\hat{\sigma}_{u}^{2}\left(\sum_{i=1}^{N} \ddot{\mathbf{X}}_{i}^{\prime} \ddot{\mathbf{X}}_{i}\right)^{-1}=\hat{\sigma}_{u}^{2}\left(\sum_{i=1}^{N} \sum_{t=1}^{T} \ddot{\mathbf{x}}_{i t}^{\prime} \ddot{\mathbf{x}}_{i t}\right)^{-1}$.

The asymptotic standard errors of the FE estimates are obtained as the square roots of the diagonal elements of the matrix (10.54).

Expression (10.54) is very convenient because it looks just like the usual OLS variance matrix estimator that would be reported from the pooled OLS regression (10.48). However, there is one catch, and this comes in obtaining the estimator $\hat{\sigma}_{u}^{2}$ of $\sigma_{u}^{2}$. The errors in the transformed model are $\ddot{u}_{i t}$, and these errors are what the OLS residuals from regression (10.48) estimate. Since $\sigma_{u}^{2}$ is the variance of $u_{i t}$, we must use a little care.

To see how to estimate $\sigma_{u}^{2}$, we use equation (10.51) summed across $t: \sum_{t=1}^{T} \mathrm{E}\left(\ddot{u}_{i t}^{2}\right)= (T-1) \sigma_{u}^{2}$, and so $[N(T-1)]^{-1} \sum_{i=1}^{N} \sum_{t=1}^{T} \mathrm{E}\left(\ddot{u}_{i t}^{2}\right)=\sigma_{u}^{2}$. Now, define the fixed effects residuals as\\
$\hat{\ddot{u}}_{i t}=\ddot{y}_{i t}-\ddot{\mathbf{x}}_{i t} \hat{\boldsymbol{\beta}}_{F E}, \quad t=1,2, \ldots, T ; i=1,2, \ldots, N$,\\
which are simply the OLS residuals from the pooled regression (10.48). Then a consistent estimator of $\sigma_{u}^{2}$ under Assumptions FE.1-FE. 3 is\\
$\hat{\sigma}_{u}^{2}=\mathrm{SSR} /[N(T-1)-K]$,\\
where $\mathrm{SSR}=\sum_{i=1}^{N} \sum_{t=1}^{T} \hat{\ddot{u}}_{i t}^{2}$. The subtraction of $K$ in the denominator of equation (10.56) does not matter asymptotically, but it is standard to make such a correction. In fact, under Assumptions FE.1-FE.3, it can be shown that $\hat{\sigma}_{u}^{2}$ is actually an unbiased estimator of $\sigma_{u}^{2}$ conditional on $\mathbf{X}$ (and therefore unconditionally as well).

Pay careful attention to the denominator in equation (10.56). This is not the degrees of freedom that would be obtained from regression (10.48). In fact, the usual variance estimate from regression (10.48) would be SSR/( NT - K), which has a probability limit less than $\sigma_{u}^{2}$ as $N$ gets large. The difference between SSR/( $N T-K$ ) and equation (10.56) can be substantial when $T$ is small.

The upshot of all this is that the usual standard errors reported from the regression (10.48) will be too small on average because they use the incorrect estimate of $\sigma_{u}^{2}$. Of course, computing equation (10.56) directly is pretty trivial. But, if a standard regression package is used after time demeaning, it is perhaps easiest to adjust the usual standard errors directly. Since $\hat{\sigma}_{u}$ appears in the standard errors, each standard error is simply multiplied by the factor $\{(N T-K) /[N(T-1)-K]\}^{1 / 2}$. As an example, if $N=500, T=3$, and $K=10$, the correction factor is about 1.227.

If an econometrics package has an option for explicitly obtaining fixed effects estimates using panel data, $\sigma_{u}^{2}$ will be properly estimated, and you do not have to worry about adjusting the standard errors. Many software packages also compute an estimate of $\sigma_{c}^{2}$, which is useful to determine how large the variance of the unobserved component is to the variance of the idiosyncratic component. Given $\hat{\boldsymbol{\beta}}_{F E}, \hat{\sigma}_{v}^{2}= (N T-K)^{-1} \sum_{i=1}^{N} \sum_{t=1}^{T}\left(y_{i t}-\mathbf{x}_{i t} \hat{\boldsymbol{\beta}}_{F E}\right)^{2}$ is a consistent estimator of $\sigma_{v}^{2}=\sigma_{c}^{2}+\sigma_{u}^{2}$, and so a consistent estimator of $\sigma_{c}^{2}$ is $\hat{\sigma}_{v}^{2}-\hat{\sigma}_{u}^{2}$. (See Problem 10.14 for a discussion of why the estimated variance of the unobserved effect in an FE analysis is generally larger than that for an RE analysis.)

Example 10.5 (FE Estimation of the Effects of Job Training Grants): Using the data in JTRAIN1.RAW, we estimate the effect of job training grants using the FE estimator. The variable union has been dropped because it does not vary over time for any of the firms in the sample. The estimated equation with standard errors is

$$
\widehat{\log (s c r a p)}=\underset{(.109)}{-.080} d 88-\underset{(.133)}{.247} d 89-\underset{(.151)}{.252} \text { grant }-\underset{(.210)}{.422} \text { grant }_{-1}
$$

Compared with the RE estimations the grant is estimated to have a larger effect, both contemporaneously and lagged one year. The $t$ statistics are also somewhat more significant with FE .

Under Assumptions FE.1-FE.3, multiple restrictions are most easily tested using an $F$ statistic, provided the degrees of freedom are appropriately computed. Let $\mathrm{SSR}_{u r}$ be the unrestricted SSR from regression (10.48), and let $\mathrm{SSR}_{r}$ denote the restricted sum of squared residuals from a similar regression, but with $Q$ restrictions imposed on $\boldsymbol{\beta}$. Then

$$
F=\frac{\left(\mathrm{SSR}_{r}-\mathrm{SSR}_{u r}\right)}{\mathrm{SSR}_{u r}} \cdot \frac{[N(T-1)-K]}{Q}
$$

is approximately $F$ distributed with $Q$ and $N(T-1)-K$ degrees of freedom. (The precise statement is that $Q \cdot F \sim \chi_{Q}^{2}$ as $N \rightarrow \infty$ under $\mathrm{H}_{0}$.) When this equation is applied to Example 10.5, the $F$ statistic for joint significance of grant and grant $_{-1}$ is $F=2.23$, with $p$-value $=.113$.

\subsection*{10.5.3 Dummy Variable Regression}
So far we have viewed the $c_{i}$ as being unobservable random variables, and for most applications this approach gives the appropriate interpretation of $\boldsymbol{\beta}$. Traditional approaches to fixed effects estimation view the $c_{i}$ as parameters to be estimated\\
along with $\boldsymbol{\beta}$. In fact, if Assumption FE. 2 is changed to its finite sample version, $\operatorname{rank}\left(\ddot{\mathbf{X}}^{\prime} \ddot{\mathbf{X}}\right)=K$, then the model under Assumptions FE.1-FE. 3 satisfies the GaussMarkov assumptions conditional on $\mathbf{X}$.

If the $c_{i}$ are parameters to estimate, how would we estimate each $c_{i}$ along with $\boldsymbol{\beta}$ ? One possibility is to define $N$ dummy variables, one for each cross section observation: $d n_{i}=1$ if $n=i, d n_{i}=0$ if $n \neq i$. Then, run the pooled OLS regression\\
$y_{i t}$ on $d 1_{i}, d 2_{i}, \ldots, d N_{i}, \mathbf{x}_{i t}, \quad t=1,2, \ldots, T ; i=1,2, \ldots, N$.

Then, $\hat{c}_{1}$ is the coefficient on $d 1_{i}, \hat{c}_{2}$ is the coefficient on $d 2_{i}$, and so on.\\
It is a nice exercise in least squares mechanics-in particular, partitioned regression (see Davidson and MacKinnon, 1993, Sect. 1.4)-to show that the estimator of $\boldsymbol{\beta}$ obtained from regression (10.57) is, in fact, the FE estimator. This is why $\hat{\boldsymbol{\beta}}_{F E}$ is sometimes referred to as the dummy variable estimator. Also, the residuals from regression (10.57) are identical to the residuals from regression (10.48). One benefit of regression (10.57) is that it produces the appropriate estimate of $\sigma_{u}^{2}$ because it uses $N T-N-K=N(T-1)-K$ as the degrees of freedom. Therefore, if it can be done, regression (10.57) is a convenient way to carry out FE analysis under Assumptions FE.1-FE.3.

There is an important difference between the $\hat{c}_{i}$ and $\hat{\boldsymbol{\beta}}_{F E}$. We already know that $\hat{\boldsymbol{\beta}}_{F E}$ is consistent with fixed $T$ as $N \rightarrow \infty$. This is not the case with the $\hat{c}_{i}$. Each time a new cross section observation is added, another $c_{i}$ is added, and information does not accumulate on the $c_{i}$ as $N \rightarrow \infty$. Each $\hat{c}_{i}$ is an unbiased estimator of $c_{i}$ when the $c_{i}$ are treated as parameters, at least if we maintain Assumption FE. 1 and the finite sample analogue of Assumption FE.2. When we add Assumption FE.3, the GaussMarkov assumptions hold (conditional on $\mathbf{X}$ ), and $\hat{c}_{1}, \hat{c}_{2}, \ldots, \hat{c}_{N}$ are best linear unbiased conditional on $\mathbf{X}$. (The $\hat{c}_{i}$ give practical examples of estimators that are unbiased but not consistent.)

Econometric software that computes fixed effects estimates rarely reports the "estimates" of the $c_{i}$ (at least in part because there are typically very many). Sometimes it is useful to obtain the $\hat{c}_{i}$ even when regression (10.57) is infeasible. Using the OLS first-order conditions for the dummy variable regression, it can be shown that\\
$\hat{c}_{i}=\bar{y}_{i}-\overline{\mathbf{x}}_{i} \hat{\boldsymbol{\beta}}_{F E}, \quad i=1,2, \ldots, N$,\\
which makes sense because $\hat{c}_{i}$ is the intercept for cross section unit $i$. We can then focus on specific units-which is more fruitful with larger $T$-or we can compute the sample average, sample median, or sample quantiles of the $\hat{c}_{i}$ to get some idea of how heterogeneity is distributed in the population. For example, the sample average\\
$\hat{\mu}_{c} \equiv N^{-1} \sum_{i=1}^{N} \hat{c}_{i}=N^{-1} \sum_{i=1}^{N}\left(\bar{y}_{i}-\overline{\mathbf{x}}_{i} \hat{\boldsymbol{\beta}}_{F E}\right)$\\
is easily seen to be a consistent estimator (as $N \rightarrow \infty$ ) of the population average $\mu_{c} \equiv \mathrm{E}\left(c_{i}\right)=\mathrm{E}\left(\bar{y}_{i}\right)-\mathrm{E}\left(\overline{\mathbf{x}}_{i}\right) \boldsymbol{\beta}$. So, although we cannot estimate each $c_{i}$ very well with small $T$, we can often estimate features of the population distribution of $c_{i}$ quite well. In fact, many econometrics software packages report $\hat{\mu}_{c}$ as the "intercept" along with $\hat{\boldsymbol{\beta}}_{F E}$ in fixed effects estimation. In other words, the intercept is simply an estimate of the average heterogeneity. (Seeing an "intercept" with fixed effects output can be a bit confusing because we have seen that the within transformation eliminates any time-constant explanatory variables. But having done that, we can then use the fixed effects estimates of $\boldsymbol{\beta}$ to estimate the mean of the heterogeneity distribution, and this is what is reported as the "intercept.")

As with RE estimation, we can consistently estimate the variance $\sigma_{c}^{2}$ if we add the assumption that $\left\{u_{i t}: t=1,2, \ldots\right\}$ is serially uncorrelated with constant variance, as implied by Assumption FE.3. First, we can consistently estimate the variance of the composite error, $\sigma_{v}^{2}=\sigma_{c}^{2}+\sigma_{u}^{2}$, because $v_{i t}=y_{i t}-\mathbf{x}_{i t} \boldsymbol{\beta}$; therefore,

$$
\hat{\sigma}_{v}^{2}=(N T-K)^{-1} \sum_{i=1}^{N} \sum_{t=1}^{T}\left(y_{i t}-\mathbf{x}_{i t} \hat{\boldsymbol{\beta}}_{F E}\right)^{2}
$$

is consistent for $\sigma_{v}^{2}$ as $N \rightarrow \infty$ (where the subtraction of $K$ is a standard degrees-offreedom adjustment). Then, we can estimate $\sigma_{c}^{2}$ as $\hat{\sigma}_{c}^{2}=\hat{\sigma}_{v}^{2}-\hat{\sigma}_{u}^{2}$, where $\hat{\sigma}_{u}^{2}$ is given by (10.56). We can use $\hat{\sigma}_{c}^{2}$ (assuming it is positive) to assess the variability in heterogeneity about its mean value. (Incidentally, as shown in Problem 10.14, one must be careful in using FE to estimate $\sigma_{c}^{2}$ before applying RE because, with time-constant variables, the FE estimate of the heterogeneity variance will be too large. Pooled OLS should be used instead, as we discussed in Section 10.4.1.)

When the $c_{i}$ are treated as different intercepts, we can compute an exact test of their equality under the classical linear model assumptions (which require, in addition to Assumptions FE. 1 to FE.3, normality of the $u_{i t}$ ). The $F$ statistic has an $F$ distribution with numerator and demoninator degrees of freedom $N-1$ and $N(T-1)-K$, respectively. Interestingly, Orme and Yamagata (2005) show that the $F$ statistic is still justified without normality when $T$ is fixed and $N \rightarrow \infty$, although they assume the idiosycratic errors are homoskedastic and serially uncorrelated.

Generally, we should view the fact that the dummy variable regression (10.57) produces $\hat{\boldsymbol{\beta}}_{F E}$ as the coefficient vector on $\mathbf{x}_{i t}$ as a coincidence. While there are other\\
unobserved effects models where "estimating" the unobserved effects along with the vector $\boldsymbol{\beta}$ results in a consistent estimator of $\boldsymbol{\beta}$, there are many cases where this approach leads to trouble. As we will see in Part IV, many nonlinear panel data models with unobserved effects suffer from an incidental parameters problem, where estimating the incidental parameters, $c_{i}$, along with $\boldsymbol{\beta}$ produces an inconsistent estimator of $\boldsymbol{\beta}$.

\subsection*{10.5.4 Serial Correlation and the Robust Variance Matrix Estimator}
Recall that the FE estimator is consistent and asymptotically normal under Assumptions FE. 1 and FE.2. But without Assumption FE.3, expression (10.54) gives an improper variance matrix estimator. While heteroskedasticity in $u_{i t}$ is always a potential problem, serial correlation is likely to be more important in certain applications. When applying the FE estimator, it is important to remember that nothing rules out serial correlation in $\left\{u_{i t}: t=1, \ldots, T\right\}$. While it is true that the observed serial correlation in the composite errors, $v_{i t}=c_{i}+u_{i t}$, is dominated by the presence of $c_{i}$, there can also be serial correlation that dies out over time. Sometimes, $\left\{u_{i t}\right\}$ can have very strong serial dependence, in which case the usual FE standard errors obtained from expression (10.54) can be very misleading. This possibility tends to be a bigger problem with large $T$. (As we will see, there is no reason to worry about serial correlation in $u_{i t}$ when $T=2$.)

Testing the idiosyncratic errors, $\left\{u_{i t}\right\}$, for serial correlation is somewhat tricky. A key point is that we cannot estimate the $u_{i t}$; because of the time demeaning used in FE, we can only estimate the time-demeaned errors, $\ddot{u}_{i t}$. As shown in equation (10.52), the time-demeaned errors are negatively correlated if the $u_{i t}$ are uncorrelated. When $T=2, \ddot{u}_{i 1}=-\ddot{u}_{i 2}$ for all $i$, and so there is perfect negative correlation. This finding shows that for $T=2$ it is pointless to use the $\ddot{u}_{i t}$ to test for any kind of serial correlation pattern.

When $T \geq 3$, we can use equation (10.52) to determine if there is serial correlation in $\left\{u_{i t}\right\}$. Naturally, we use the fixed effects residuals, $\hat{\ddot{u}}_{i t}$. One simplification is obtained by applying Problem 7.4: we can ignore the estimation error in $\boldsymbol{\beta}$ in obtaining the asymptotic distribution of any test statistic based on sample covariances and variances. In other words, it is as if we are using the $\ddot{u}_{i t}$, rather than the $\hat{\ddot{u}}_{i t}$. The test is complicated by the fact that the $\left\{\ddot{u}_{i t}\right\}$ are serially correlated under the null hypothesis. There are two simple possibilities for dealing with this. First, we can just use any two time periods (say, the last two), to test equation (10.52) using a simple regression. In other words, run the regression\\
$\hat{\ddot{u}}_{i T}$ on $\hat{\ddot{u}}_{i, T-1}, \quad i=1, \ldots, N$,\\
and use $\hat{\delta}$, the coefficient on $\hat{\ddot{u}}_{i, T-1}$, along with its standard error, to test $\mathrm{H}_{0}: \delta= -1 /(T-1)$, where $\delta=\operatorname{Corr}\left(\ddot{u}_{i, T-1}, \ddot{u}_{i T}\right)$. Under Assumptions FE.1-FE.3, the usual $t$ statistic has an asymptotic normal distribution. (It is trivial to make this test robust to heteroskedasticity.)

Alternatively, we can use more time periods if we make the $t$ statistic robust to arbitrary serial correlation. In other words, run the pooled OLS regression\\
$\hat{\ddot{u}}_{i t}$ on $\hat{\ddot{u}}_{i, t-1}, \quad t=3, \ldots, T ; i=1, \ldots, N$,\\
and use the fully robust standard error for pooled OLS; see equation (7.26). It may seem a little odd that we make a test for serial correlation robust to serial correlation, but this need arises because the null hypothesis is that the time-demeaned errors are serially correlated. This approach clearly does not produce an optimal test against, say, AR(1) correlation in the $u_{i t}$, but it is very simple and may be good enough to indicate a problem.

Without Assumption FE.3, the asymptotic variance of $\sqrt{N}\left(\hat{\boldsymbol{\beta}}_{F E}-\boldsymbol{\beta}\right)$ has the sandwich form, $\left[\mathrm{E}\left(\ddot{\mathbf{X}}_{i}^{\prime} \ddot{\mathbf{X}}_{i}\right)\right]^{-1} \mathrm{E}\left(\ddot{\mathbf{X}}_{i}^{\prime} \ddot{\mathbf{u}}_{i}^{\prime} \ddot{\mathbf{u}}_{i}^{\prime} \ddot{\mathbf{X}}_{i}\right)\left[\mathrm{E}\left(\ddot{\mathbf{X}}_{i}^{\prime} \ddot{\mathbf{X}}_{i}\right)\right]^{-1}$. Thefore, if we suspect or find evidence of serial correlation, we should, at a minimum, compute a fully robust variance matrix estimator (and corresponding test statistics) for the FE estimator. But this is just as in Chapter 7 applied to the time-demeaned set of equations. Let $\hat{\ddot{\mathbf{u}}}_{i}=\ddot{\mathbf{y}}_{i}-\ddot{\mathbf{X}}_{i} \hat{\boldsymbol{\beta}}_{F E}, i=1,2, \ldots, N$, denote the $T \times 1$ vectors of FE residuals. Direct application of equation (7.28) gives\\
$\operatorname{Avar}\left(\hat{\boldsymbol{\beta}}_{F E}\right)=(\ddot{\mathbf{X}} \prime \ddot{\mathbf{X}})^{-1}\left(\sum_{i=1}^{N} \ddot{\mathbf{X}}_{i}^{\prime} \hat{\ddot{u}}_{i}^{\prime} \hat{\mathbf{u}}_{i}^{\prime} \ddot{\mathbf{X}}_{i}\right)(\ddot{\mathbf{X}} \prime \ddot{\mathbf{X}})^{-1}$,\\
which was suggested by Arellano (1987) and follows from the general results of White (1984, Chap. 6). The robust variance matrix estimator is valid in the presence of any heteroskedasticity or serial correlation in $\left\{u_{i t}: t=1, \ldots, T\right\}$, provided that $T$ is small relative to $N$. (Remember, equation (7.28) is generally justified for fixed $T$, $N \rightarrow \infty$ asymptotics.) The robust standard errors are obtained as the square roots of the diagonal elements of the matrix (10.59), and matrix (10.59) can be used as the $\hat{\mathbf{V}}$ matrix in constructing Wald statistics. Unfortunately, the sum of squared residuals form of the $F$ statistic is no longer asymptotically valid when Assumption FE. 3 fails.

Example 10.5 (continued): We now report the robust standard errors for the $\log$ (scrap) equation along with the usual FE standard errors:

$$
\widehat{\log \widehat{(s c r a p})}=\underset{\substack{(.109) \\[.096]}}{.080} d 88-\underset{(.133)}{.247} d 89-\underset{(.151)}{.252} \text { grant }-\underset{(.210)}{.422} \text { grant }_{-1} .
$$

The robust standard error on grant is actually smaller than the usual standard error, while the robust standard error on grant $_{-1}$ is larger than the usual one. As a result, the absolute value of the $t$ statistic on grant $_{-1}$ drops from about 2 to just over 1.5.

Remember, with fixed $T$ as $N \rightarrow \infty$, the robust standard errors are just as valid asymptotically as the nonrobust ones when Assumptions FE.1-FE. 3 hold. But the usual standard errors and test statistics may be better behaved under Assumptions FE.1-FE. 3 if $N$ is not very large relative to $T$, especially if $u_{i t}$ is normally distributed.

\subsection*{10.5.5 Fixed Effects Generalized Least Squares}
Recall that Assumption FE. 3 can fail for two reasons. The first is that the conditional variance matrix does not equal the unconditional variance matrix: $\mathrm{E}\left(\mathbf{u}_{i} \mathbf{u}_{i}^{\prime} \mid \mathbf{x}_{i}, c_{i}\right) \neq \mathrm{E}\left(\mathbf{u}_{i} \mathbf{u}_{i}^{\prime}\right)$. Even if $\mathrm{E}\left(\mathbf{u}_{i} \mathbf{u}_{i}^{\prime} \mid \mathbf{x}_{i}, c_{i}\right)=\mathrm{E}\left(\mathbf{u}_{i} \mathbf{u}_{i}^{\prime}\right)$, the unconditional variance matrix may not be scalar: $\mathrm{E}\left(\mathbf{u}_{i} \mathbf{u}_{i}^{\prime}\right) \neq \sigma_{u}^{2} \mathbf{I}_{T}$, which means either that the variance of $u_{i t}$ changes with $t$ or, probably more important, that there is serial correlation in the idiosyncratic errors. The robust variance matrix (10.59) is valid in any case.

Rather than compute a robust variance matrix for the FE estimator, we can instead relax Assumption FE. 3 to allow for an unrestricted, albeit constant, conditional covariance matrix. This is a natural route to follow if the robust standard errors of the FE estimator are too large to be useful and if there is evidence of serial dependence or a time-varying variance in the $u_{i t}$.

ASSUMPTION FEGLS.3: $\quad \mathrm{E}\left(\mathbf{u}_{i} \mathbf{u}_{i}^{\prime} \mid \mathbf{x}_{i}, c_{i}\right)=\boldsymbol{\Lambda}$, a $T \times T$ positive definite matrix.\\
Under Assumption FEGLS.3, $\mathrm{E}\left(\ddot{\mathbf{u}}_{i} \ddot{\mathbf{u}}_{i}^{\prime} \mid \ddot{\mathbf{x}}_{i}\right)=\mathrm{E}\left(\ddot{\mathbf{u}}_{i} \ddot{\mathbf{u}}_{i}^{\prime}\right)$. Further, using $\ddot{\mathbf{u}}_{i}=\mathbf{Q}_{T} \mathbf{u}_{i}$,\\
$\mathrm{E}\left(\ddot{\mathbf{u}}_{i} \ddot{\mathbf{u}}_{i}^{\prime}\right)=\mathbf{Q}_{T} \mathrm{E}\left(\mathbf{u}_{i} \mathbf{u}_{i}^{\prime}\right) \mathbf{Q}_{T}=\mathbf{Q}_{T} \boldsymbol{\Lambda} \mathbf{Q}_{T}$.\\
which has rank $T-1$. The deficient rank in expression (10.60) causes problems for the usual approach to GLS, because the variance matrix cannot be inverted. One way to proceed is to use a generalized inverse. A much easier approach-and one that turns out to be algebraically identical-is to drop one of the time periods from the analysis. It can be shown (see Im, Ahn, Schmidt, and Wooldridge, 1999) that it does not matter which of these time periods is dropped: the resulting GLS estimator is the same.

For concreteness, suppose we drop time period $T$, leaving the equations


\begin{align*}
& \ddot{y}_{i 1}=\ddot{\mathbf{x}}_{i 1} \boldsymbol{\beta}+\ddot{u}_{i 1} \\
& \vdots  \tag{10.61}\\
& \ddot{y}_{i, T-1}=\ddot{\mathbf{x}}_{i, T-1} \boldsymbol{\beta}+\ddot{u}_{i, T-1} .
\end{align*}


So that we do not have to introduce new notation, we write the system (10.61) as equation (10.49), with the understanding that now $\ddot{\mathbf{y}}_{i}$ is $(T-1) \times 1, \ddot{\mathbf{X}}_{i}$ is $(T-1) \times K$, and $\ddot{\mathbf{u}}_{i}$ is $(T-1) \times 1$. Define the $(T-1) \times(T-1)$ positive definite matrix $\boldsymbol{\Omega} \equiv \mathrm{E}\left(\ddot{\mathbf{u}}_{i} \ddot{\mathbf{u}}_{i}^{\prime}\right)$. We do not need to make the dependence of $\boldsymbol{\Omega}$ on $\boldsymbol{\Lambda}$ and $\mathbf{Q}_{T}$ explicit; the key point is that, if no restrictions are made on $\boldsymbol{\Lambda}$, then $\boldsymbol{\Omega}$ is also unrestricted.

To estimate $\boldsymbol{\Omega}$, we estimate $\boldsymbol{\beta}$ by fixed effects in the first stage. After dropping the last time period for each $i$, define the $(T-1) \times 1$ residuals $\ddot{\mathbf{u}}_{i}=\ddot{\mathbf{y}}_{i}-\ddot{\mathbf{X}}_{i} \hat{\boldsymbol{\beta}}_{F E}, i= 1,2, \ldots, N$. A consistent estimator of $\boldsymbol{\Omega}$ is\\
$\hat{\boldsymbol{\Omega}}=N^{-1} \sum_{i=1}^{N} \check{\mathbf{u}}_{i} \check{\mathbf{u}}_{i}^{\prime}$.

The fixed effects GLS (FEGLS) estimator is defined by\\
$\hat{\boldsymbol{\beta}}_{F E G L S}=\left(\sum_{i=1}^{N} \ddot{\mathbf{X}}_{i}^{\prime} \hat{\boldsymbol{\Omega}}^{-1} \ddot{\mathbf{X}}_{i}\right)^{-1}\left(\sum_{i=1}^{N} \ddot{\mathbf{X}}_{i}^{\prime} \hat{\boldsymbol{\Omega}}^{-1} \ddot{\mathbf{y}}_{i}\right)$,\\
where $\ddot{\mathbf{X}}_{i}$ and $\ddot{\mathbf{y}}_{i}$ are defined with the last time period dropped. For consistency of FEGLS, we replace Assumption FE. 2 with a new rank condition:\\
assumption FEGLS.2: $\quad \operatorname{rank} \mathrm{E}\left(\ddot{\mathbf{X}}_{i}^{\prime} \boldsymbol{\Omega}^{-1} \ddot{\mathbf{X}}_{i}\right)=K$.\\
Under Assumptions FE. 1 and FEGLS.2, the FEGLS estimator is consistent. When we add Assumption FEGLS.3, the asymptotic variance is easy to estimate:\\
$\widehat{\operatorname{Avar}\left(\hat{\boldsymbol{\beta}}_{F E G L S}\right)}=\left(\sum_{i=1}^{N} \ddot{\mathbf{X}}_{i}^{\prime} \hat{\boldsymbol{\Omega}}^{-1} \ddot{\mathbf{X}}_{i}\right)^{-1}$.\\
The sum of squared residual statistics from FGLS can be used to test multiple restrictions. Note that $G=T-1$ in the $F$ statistic in equation (7.57).

The FEGLS estimator was proposed by Kiefer (1980) when the $c_{i}$ are treated as parameters. As we just showed, the procedure consistently estimates $\boldsymbol{\beta}$ when we view $c_{i}$ as random and allow it to be arbitrarily correlated with $\mathbf{x}_{i t}$.

The FEGLS estimator is asymptotically no less efficient than the FE estimator under Assumption FEGLS.3, even when $\boldsymbol{\Lambda}=\sigma_{u}^{2} \mathbf{I}_{T}$. Generally, if $\boldsymbol{\Lambda} \neq \sigma_{u}^{2} \mathbf{I}_{T}$, FEGLS is more efficient than FE, but this conclusion relies on the large- $N$, fixed- $T$ asymptotics. Unfortunately, because FEGLS still uses the fixed effects transformation to remove $c_{i}$, it can have large asymptotic standard errors if the matrices $\ddot{\mathbf{X}}_{i}$ have columns close to zero.

Rather than allowing $\boldsymbol{\Omega}$ to be an unrestricted matrix, we can impose restrictions on $\boldsymbol{\Lambda}$ that imply $\boldsymbol{\Omega}$ has a restricted form. For example, Bhargava, Franzini, and

Narendranathan (1982) (BFN) assume that $\left\{u_{i t}\right\}$ follows a stable, homoskedastic AR(1) model. This assumption implies that $\boldsymbol{\Omega}$ depends on only three parameters, $\sigma_{c}^{2}$, $\sigma_{u}^{2}$, and the AR coefficient, $\rho$, no matter how large $T$ is. BFN obtain a transformation that eliminates the unobserved effect, $c_{i}$, and removes the serial correlation in $u_{i t}$. They also propose estimators of $\rho$, so that feasible GLS is possible.

Modeling $\left\{u_{i t}\right\}$ as a specific time series process is attractive when $N$ is not very large relative to $T$, as estimating an unrestricted covariance matrix for $\ddot{\mathbf{u}}_{i}$ (the $(T-1) \times 1$ vector of time-demeaned errors) without large $N$ can lead to poor finitesample performance of the FGLS estimator. However, the only general statements we can make concern fixed- $T, N \rightarrow \infty$ asymptotics. In this scenario, the FGLS estimator that uses unrestricted $\boldsymbol{\Omega}$ is no less asymptotically efficient than an FGLS estimator that puts restrictions on $\boldsymbol{\Omega}$. And, if the restrictions on $\boldsymbol{\Omega}$ are incorrect, the estimator that imposes the restrictions is less asymptotically efficient. Therefore, on purely theoretical grounds, we prefer an estimator of the type in equation (10.62).

As with any FGLS estimator, it is always a good idea to compute a fully robust variance matrix estimator for the FEGLS estimator. The robust variance matrix estimator is still given by equation (7.52), but now we insert the time-demeaned regressors and FEGLS residuals in place of $\mathbf{X}_{i}$ and $\hat{\mathbf{u}}_{i}$, respectively. Therefore, the fully robust variance matrix estimator for the FEGLS estimator is\\
$\widehat{\operatorname{Avar}\left(\hat{\boldsymbol{\beta}}_{F E G L S}\right)}=\left(\sum_{i=1}^{N} \ddot{\mathbf{X}}_{i}^{\prime} \hat{\boldsymbol{\Omega}}^{-1} \ddot{\mathbf{X}}_{i}\right)^{-1}\left(\sum_{i=1}^{N} \ddot{\mathbf{X}}_{i}^{\prime} \hat{\boldsymbol{\Omega}}^{-1} \hat{\ddot{\mathbf{u}}}_{i} \hat{\ddot{\mathbf{u}}}_{i}^{\prime} \hat{\boldsymbol{\Omega}}^{-1} \ddot{\mathbf{X}}_{i}^{\prime}\right)\left(\sum_{i=1}^{N} \ddot{\mathbf{X}}_{i}^{\prime} \hat{\boldsymbol{\Omega}}^{-1} \ddot{\mathbf{X}}_{i}\right)^{-1}$,\\
where $\hat{\ddot{\mathbf{u}}}_{i}=\ddot{\mathbf{y}}_{i}-\ddot{\mathbf{X}}_{i} \hat{\boldsymbol{\beta}}_{\text {FEGLS }}$ are the $(T-1) \times 1$ vectors of FEGLS residuals. If $\hat{\boldsymbol{\Omega}}$ is unrestricted the robust variance matrix estimator is robust to system heteroskedasticity, which is generally present if $\mathrm{E}\left(\ddot{\mathbf{u}}_{i} \ddot{\mathbf{u}}_{i}^{\prime} \mid \ddot{\mathbf{X}}_{i}\right) \neq \mathrm{E}\left(\ddot{\mathbf{u}}_{i} \ddot{\mathbf{u}}_{i}^{\prime}\right)$. Remember, even if we allow the unconditional variance matrix of $\ddot{\mathbf{u}}_{i}$ to be unrestricted, we can never guarantee that the conditional variance matrix does not depend on the regressors. When $\hat{\boldsymbol{\Omega}}$ is restricted, such as in the $\operatorname{AR}(1)$ model for $\left\{u_{i t}\right\}$, the robust estimator is also robust to the retrictions imposed being incorrect. Using an AR(1) model for the idiosyncratic errors might be substantially more efficient than just using the usual FE estimator, and we can guard against incorrect asymptotic inference by using a fully robust variance matrix estimator. As we discussed earlier, if $T$ is somewhat large, for better small-sample properties we may wish to impose restrictions on $\boldsymbol{\Omega}$ rather than use an unrestricted estimator.

Baltagi and Li (1991) consider FEGLS estimation when $\left\{u_{i t}\right\}$ follows an $\operatorname{AR}(1)$ process. But their estimator of $\rho$ is based on the FE residual-say, from the regression $\hat{\ddot{u}}_{i t}$ on $\hat{\ddot{u}}_{i, t-1}, t=2, \ldots, T$-and this estimator is inconsistent for fixed $T$ with\\
$N \rightarrow \infty$. In fact, if $\left\{u_{i t}\right\}$ is serially uncorrelated, we know from equation (10.52) that $\operatorname{plim}(\hat{\rho})=-1 /(T-1)$. With large $T$ the inconsistency may be small, but what are the consequences generally of using an inconsistent estimator of $\rho$ ? We know by now that using an inconsistent estimator of a variance matrix does not lead to inconsistency of a feasible GLS estimator, and this is no exception. But, of course, inference should be made fully robust, because if $\operatorname{plim}(\hat{\rho}) \neq \rho$, GLS does not fully eliminate the serial correlation in $\left\{u_{i t}\right\}$. Even if $\operatorname{plim}(\hat{\rho})$ is "close" to $\rho$, we may want to guard against more general forms of serial correlation or violation of system homoskedasticity.

\subsection*{10.5.6 Using Fixed Effects Estimation for Policy Analysis}
There are other ways to interpret the FE transformation to illustrate why fixed effects is useful for policy analysis and program evaluation. Consider the model\\
$y_{i t}=\mathbf{x}_{i t} \boldsymbol{\beta}+v_{i t}=\mathbf{z}_{i t} \boldsymbol{\gamma}+\delta w_{i t}+v_{i t}$,\\
where $v_{i t}$ may or may not contain an unobserved effect. Let $w_{i t}$ be the policy variable of interest; it could be continuous or discrete. The vector $\mathbf{z}_{i t}$ contains other controls that might be correlated with $w_{i t}$, including time-period dummy variables.

As an exercise, you can show that sufficient for consistency of fixed effects, along with the rank condition FE.2, is\\
$\mathrm{E}\left[\mathbf{x}_{i t}^{\prime}\left(v_{i t}-\bar{v}_{i}\right)\right]=\mathbf{0}, \quad t=1,2, \ldots, T$.\\
This assumption shows that each element of $\mathbf{x}_{i t}$, and in particular the policy variable $w_{i t}$, can be correlated with $\bar{v}_{i}$. What FE requires for consistency is that $w_{i t}$ be uncorrelated with deviations of $v_{i t}$ from the average over the time period. So a policy variable, such as program participation, can be systematically related to the persistent component in the error $v_{i t}$ as measured by $\bar{v}_{i}$. It is for this reason that FE is often superior to pooled OLS or random effects for applications where participation in a program is determined by preprogram attributes that also affect $y_{i t}$.

\subsection*{10.6 First Differencing Methods}
\subsection*{10.6.1 Inference}
In Section 10.1 we used differencing to eliminate the unobserved effect $c_{i}$ with $T=2$. We now study the differencing transformation in the general case of model (10.41). For completeness, we state the first assumption as follows.\\
assumption FD.1: Same as Assumption FE.1.\\
We emphasize that the model and the interpretation of $\boldsymbol{\beta}$ are exactly as in Section 10.5 . What differs is our method for estimating $\boldsymbol{\beta}$.

Lagging the model (10.41) one period and subtracting gives\\
$\Delta y_{i t}=\Delta \mathbf{x}_{i t} \boldsymbol{\beta}+\Delta u_{i t}, \quad t=2,3, \ldots, T$,\\
where $\Delta y_{i t}=y_{i t}-y_{i, t-1}, \Delta \mathbf{x}_{i t}=\mathbf{x}_{i t}-\mathbf{x}_{i, t-1}$, and $\Delta u_{i t}=u_{i t}-u_{i, t-1}$. As with the FE transformation, this first-differencing transformation eliminates the unobserved effect $c_{i}$. In differencing we lose the first time period for each cross section: we now have $T-1$ time periods for each $i$, rather than $T$. If we start with $T=2$, then, after differencing, we arrive at one time period for each cross section: $\Delta y_{i 2}=\Delta \mathbf{x}_{i 2} \boldsymbol{\beta}+\Delta u_{i 2}$.

Equation (10.63) makes it clear that the elements of $\mathbf{x}_{i t}$ must be time varying (for at least some cross section units); otherwise $\Delta \mathbf{x}_{i t}$ has elements that are identically zero for all $i$ and $t$. Also, while the intercept in the original equation gets differenced away, equation (10.63) contains changes in time dummies if $\mathbf{x}_{i t}$ contains time dummies. In the $T=2$ case, the coefficient on the second-period time dummy becomes the intercept in the differenced equation. If we difference the general equation (10.43) we get


\begin{align*}
\Delta y_{i t}= & \theta_{2}\left(\Delta d 2_{t}\right)+\cdots+\theta_{T}\left(\Delta d T_{t}\right)+\left(\Delta d 2_{t}\right) \mathbf{z}_{i} \gamma_{2} \\
& +\cdots+\left(\Delta d T_{t}\right) \mathbf{z}_{i} \gamma_{T}+\Delta \mathbf{w}_{i t} \boldsymbol{\delta}+\Delta u_{i t} \tag{10.64}
\end{align*}


The parameters $\theta_{1}$ and $\gamma_{1}$ are not identified because they disappear from the transformed equation, just as with fixed effects.

The first-difference (FD) estimator, $\hat{\boldsymbol{\beta}}_{F D}$, is the pooled OLS estimator from the regression\\
$\Delta y_{i t}$ on $\Delta \mathbf{x}_{i t}, \quad t=2, \ldots, T ; i=1,2, \ldots, N$.

Under Assumption FD.1, pooled OLS estimation of the first-differenced equations will be consistent because\\
$\mathrm{E}\left(\Delta \mathbf{x}_{i t}^{\prime} \Delta u_{i t}\right)=0, \quad t=2,3, \ldots, T$.

Therefore, Assumption POLS. 1 from Section 7.8 holds. In fact, strict exogeneity holds in the FD equation:\\
$\mathrm{E}\left(\Delta u_{i t} \mid \Delta \mathbf{x}_{i 2}, \Delta \mathbf{x}_{i 3}, \ldots, \Delta \mathbf{x}_{i T}\right)=0, \quad t=2,3, \ldots, T$,\\
which means the FD estimator is actually unbiased conditional on $\mathbf{X}$.\\
To arrive at assumption (10.66) we clearly can get by with an assumption weaker than Assumption FD.1. The key point is that assumption (10.66) fails if $u_{i t}$ is corre-\\
lated with $\mathbf{x}_{i, t-1}, \mathbf{x}_{i t}$, or $\mathbf{x}_{i, t+1}$, and so we just assume that $\mathbf{x}_{i s}$ is uncorrelated with $u_{i t}$ for all $t$ and $s$.

For completeness, we state the rank condition for the FD estimator:\\
ASSUMPTION FD.2: $\quad \operatorname{rank}\left(\sum_{t=2}^{T} \mathrm{E}\left(\Delta \mathbf{x}_{i t}^{\prime} \Delta \mathbf{x}_{i t}\right)\right)=K$.\\
Assumption FD. 2 clearly rules out explanatory variables in $\mathbf{x}_{i t}$ that are fixed across time for all $i$. It also excludes perfect collinearity among the time-varying variables after differencing. There are subtle ways in which perfect collinearity can arise in $\Delta \mathbf{x}_{i t}$. For example, suppose that in a panel of individuals working in every year, we collect information on labor earnings, workforce experience ( $\operatorname{exper}_{i t}$ ), and other variables. Then, because experience increases by one each year for every person in the sample, $\Delta$ exper $_{i t}=1$ for all $i$ and $t$. If $\mathbf{x}_{i t}$ also contains a full set of year dummies, then $\Delta$ exper $_{i t}$ is perfectly collinear with the changes in the year dummies. This is easiest to see in the $T=2$ case, where $\Delta d 2_{t}=1$ for $t=2$, and so $\Delta \operatorname{exper}_{i 2}=\Delta d 2_{2}$ (remember, we only use the second time period in estimation). In the general case, it can be seen that for all $t=2, \ldots, T, \Delta d 2_{t}+2 \Delta d 3_{t}+\cdots+(T-1) \Delta d T_{t}=1$, and so $\Delta \operatorname{exper}_{i t}$ is perfectly collinear with $\Delta d 2_{t}, \Delta d 3_{t}, \ldots, \Delta d T_{t}$.

Assuming the data have been ordered as we discussed earlier, first differencing is easy to implement provided we keep track of which transformed observations are valid and which are not. Differences for observation numbers $1, T+1,2 T+1$, $3 T+1, \ldots$, and $(N-1) T+1$ should be set to missing. These observations correspond to the first time period for every cross section unit in the original data set; by definition, there is no first difference for the $t=1$ observations. A little care is needed so that differences between the first time period for unit $i+1$ and the last time period for unit $i$ are not treated as valid observations. Making sure these are set to missing is easy when a year variable or time period dummies have been included in the data set.

One reason to prefer the FD estimator to the FE estimator is that FD is easier to implement without special software. Are there statistical reasons to prefer FD to FE? Recall that, under Assumptions FE.1-FE.3, the FE estimator is asymptotically efficient in the class of estimators using the strict exogeneity assumption FE.1. Therefore, the FD estimator is less efficient than FE under Assumptions FE.1-FE.3. Assumption FE. 3 is key to the efficiency of FE. It assumes homoskedasticity and no serial correlation in $u_{i t}$. Assuming that the $\left\{u_{i t}: t=1,2, \ldots T\right\}$ are serially uncorrelated may be too strong. An alternative assumption is that the first differences of the idiosyncratic errors, $\left\{e_{i t} \equiv \Delta u_{i t}, t=2, \ldots, T\right\}$, are serially uncorrelated (and have constant variance).

ASSUMPTION FD.3: $\mathrm{E}\left(\mathbf{e}_{i} \mathbf{e}_{i}^{\prime} \mid \mathbf{x}_{i 1}, \ldots, \mathbf{x}_{i T}, c_{i}\right)=\sigma_{e}^{2} \mathbf{I}_{T-1}$, where $\mathbf{e}_{i}$ is the $(T-1) \times 1$ vector containing $e_{i t}, t=2, \ldots, T$.

Under Assumption FD. 3 we can write $u_{i t}=u_{i, t-1}+e_{i t}$, so that no serial correlation in $e_{i t}$ implies that $u_{i t}$ is a random walk. A random walk has substantial serial dependence, and so Assumption FD. 3 represents an opposite extreme from Assumption FE.3.

Under Assumptions FD.1-FD.3, it can be shown that the FD estimator is most efficient in the class of estimators using the strict exogeneity assumption FE.1. Further, from the pooled OLS analysis in Section 7.8,\\
$\widehat{\operatorname{Avar}\left(\hat{\boldsymbol{\beta}}_{F D}\right)}=\hat{\sigma}_{e}^{2}\left(\Delta \mathbf{X}^{\prime} \Delta \mathbf{X}\right)^{-1}$,\\
where $\hat{\sigma}_{e}^{2}$ is a consistent estimator of $\sigma_{e}^{2}$. The simplest estimator is obtained by computing the OLS residuals


\begin{equation*}
\hat{e}_{i t}=\Delta y_{i t}-\Delta \mathbf{x}_{i t} \hat{\boldsymbol{\beta}}_{F D} \tag{10.68}
\end{equation*}


from the pooled regression (10.65). A consistent estimator of $\sigma_{e}^{2}$ is\\
$\hat{\sigma}_{e}^{2}=[N(T-1)-K]^{-1} \sum_{i=1}^{N} \sum_{t=2}^{T} \hat{e}_{i t}^{2}$,\\
which is the usual error variance estimator from regression (10.65). These equations show that, under Assumptions FD.1-FD.3, the usual OLS standard errors from the FD regression (10.65) are asymptotically valid.

Unlike in the FE regression (10.48), the denominator in equation (10.69) is correctly obtained from regression (10.65). Dropping the first time period appropriately captures the lost degrees of freedom ( $N$ of them).

Under Assumption FD.3, all statistics reported from the pooled regression on the first-differenced data are asymptotically valid, including $F$ statistics based on sums of squared residuals.

\subsection*{10.6.2 Robust Variance Matrix}
If Assumption FD. 3 is violated, then, as usual, we can compute a robust variance matrix. The estimator in equation (7.26) applied in this context is\\
$\left.\widehat{\operatorname{Avar}\left(\hat{\boldsymbol{\beta}}_{F D}\right.}\right)=\left(\Delta \mathbf{X}^{\prime} \Delta \mathbf{X}\right)^{-1}\left(\sum_{i=1}^{N} \Delta \mathbf{X}_{i}^{\prime} \hat{\mathbf{e}}_{i} \hat{\mathbf{e}}_{i}^{\prime} \Delta \mathbf{X}_{i}\right)\left(\Delta \mathbf{X}^{\prime} \Delta \mathbf{X}\right)^{-1}$,\\
where $\Delta \mathbf{X}$ denotes the $N(T-1) \times K$ matrix of stacked first differences of $\mathbf{x}_{i t}$.

Example 10.6 (FD Estimation of the Effects of Job Training Grants): We now estimate the effect of job training grants on $\log$ (scrap) using first differencing. Specifically, we use pooled OLS on\\
$\Delta \log \left(\operatorname{scrap}_{i t}\right)=\delta_{1}+\delta_{2} d 89_{t}+\beta_{1} \Delta \operatorname{grant}_{i t}+\beta_{2} \Delta \operatorname{grant}_{i, t-1}+\Delta u_{i t}$.\\
Rather than difference the year dummies and omit the intercept, we simply include an intercept and a dummy variable for 1989 to capture the aggregate time effects. If we were specifically interested in the year effects from the structural model (in levels), then we should difference those as well.

The estimated equation is

$\Delta \widehat{\log (\text { scrap })}=\underset{\substack{(.091) \\[.088]}}{\underset{(.091}{.}} \underset{(.111]}{.096} d 89-\underset{(.131)}{.223} \Delta$ grant $-\underset{(.235)}{.351} \Delta$ grant $_{-1}$\\
$R^{2}=.037$,\\
where the usual standard errors are in parentheses and the robust standard errors are in brackets. We report $R^{2}$ here because it has a useful interpretation: it measures the amount of variation in the growth in the scrap rate that is explained by Δgrant and $\Delta$ grant $_{-1}$ (and d89). The estimates on grant and grant $_{-1}$ are fairly similar to the FE estimates, although grant is now statistically more significant than grant ${ }_{-1}$. The usual $F$ test for joint significance of $\Delta$ grant and $\Delta$ grant $_{-1}$ is 1.53 , with $p$-value $=.222$.

In the previous example we used a device that is common when applying FD estimation. Namely, rather than drop an overall intercept and include the differenced time dummies, we estimated an intercept and then included time dummies for $T-2$ of the remaining periods-in Example 10.6, just the third period. Generally, rather than using the $T-1$ regressors $\left(\Delta d 2_{t}, \Delta d 3_{t}, \ldots, \Delta d T_{t}\right)$, it is often more convenient to use, say $\left(1, d 3_{t}, \ldots, d T_{t}\right)$. Because these sets of regressors involving the time dummies are nonsingular linear transformations of each other, the estimated coefficients on the other variables do not change, nor do their standard errors or any test statistics. In most regression packages, including an overall intercept makes it easier to obtain the appropriate $R$-squared for the FD equation. Of course, if one is interested in the coefficients on the original time dummies, then it is easiest to simply include all dummies in FD form and omit an overall intercept.

\subsection*{10.6.3 Testing for Serial Correlation}
Under Assumption FD.3, the errors $e_{i t} \equiv \Delta u_{i t}$ should be serially uncorrelated. We can easily test this assumption given the pooled OLS residuals from regression\\
(10.65). Since the strict exogeneity assumption holds, we can apply the simple form of the test in Section 7.8. The regression is based on $T-2$ time periods:\\
$\hat{e}_{i t}=\hat{\rho}_{1} \hat{e}_{i, t-1}+$ error $_{i t}, \quad t=3,4, \ldots, T ; i=1,2, \ldots, N$.

The test statistic is the usual $t$ statistic on $\hat{\rho}_{1}$. With $T=2$ this test is not available, nor is it necessary. With $T=3$, regression (10.71) is just a cross section regression because we lose the $t=1$ and $t=2$ time periods.

If the idiosyncratic errors $\left\{u_{i t}: t=1,2, \ldots, T\right\}$ are uncorrelated to begin with, $\left\{e_{i t}: t=2,3, \ldots, T\right\}$ will be autocorrelated. In fact, under Assumption FE. 3 it is easily shown that $\operatorname{Corr}\left(e_{i t}, e_{i, t-1}\right)=-.5$. In any case, a finding of significant serial correlation in the $e_{i t}$ warrants computing the robust variance matrix for the FD estimator.

Example 10.6 (continued): We test for $\mathrm{AR}(1)$ serial correlation in the first-differenced equation by regressing $\hat{e}_{i t}$ on $\hat{e}_{i, t-1}$ using the year 1989. We get $\hat{\rho}_{1}=.237$ with $t$ statistic $=1.76$. There is marginal evidence of positive serial correlation in the first differences $\Delta u_{i t}$. Further, $\hat{\rho}_{1}=.237$ is very different from $\rho_{1}=-.5$, which is implied by the standard random and fixed effects assumption that the $u_{i t}$ are serially uncorrelated.

An alternative to computing robust standard errors and test statistics is to use an FDGLS analysis under the assumption that $\mathrm{E}\left(\mathbf{e}_{i} \mathbf{e}_{i}^{\prime} \mid \mathbf{x}_{i}\right)$ is a constant $(T-1) \times (T-1)$ matrix. We omit the details, as they are similar to the FEGLS case in Section 10.5.5. As with FEGLS, we could impose structure on $\mathrm{E}\left(\mathbf{u}_{i} \mathbf{u}_{i}^{\prime}\right)$, such as a stable, homoskedastic $\mathrm{AR}(1)$ model, and then derive $\mathrm{E}\left(\mathbf{e}_{i} \mathbf{e}_{i}^{\prime}\right)$ in terms of a small set of parameters.

\subsection*{10.6.4 Policy Analysis Using First Differencing}
First differencing a structural equation with an unobserved effect is a simple yet powerful method of program evaluation. Many questions can be addressed by having a two-year panel data set with control and treatment groups available at two points in time.

In applying first differencing, we should difference all variables appearing in the structural equation to obtain the estimating equation, including any binary indicators indicating participation in the program. The estimates should be interpreted in the original equation because it allows us to think of comparing different units in the cross section at any point in time, where one unit receives the treatment and the other does not.

In one special case it does not matter whether the policy variable is differenced. Assume that $T=2$, and let $p r o g_{i t}$ denote a binary indicator set to one if person $i$ was in the program at time $t$. For many programs, $\operatorname{prog}_{i 1}=0$ for all $i$ : no one participated in the program in the initial time period. In the second time period, prog $_{i 2}$ is unity for\\
those who participate in the program and zero for those who do not. In this one case, $\Delta$ prog $_{i}=$ prog $_{i 2}$, and the FD equation can be written as\\
$\Delta y_{i 2}=\theta_{2}+\Delta \mathbf{z}_{i 2} \gamma+\delta_{1}$ prog $_{i 2}+\Delta u_{i 2}$.

The effect of the policy can be obtained by regressing the change in $y$ on the change in $\mathbf{z}$ and the policy indicator. When $\Delta \mathbf{z}_{i 2}$ is omitted, the estimate of $\delta_{1}$ from equation (10.72) is the difference-in-differences (DD) estimator (see Problem 10.4): $\hat{\delta}_{1}=\overline{\Delta y}_{\text {treat }}-\overline{\Delta y}_{\text {control }}$. This is similar to the DD estimator from Section 6.3-see equation (6.53)-but there is an important difference: with panel data, the differences over time are for the same cross section units.

If some people participated in the program in the first time period, or if more than two periods are involved, equation (10.72) can give misleading answers. In general, the equation that should be estimated is\\
$\Delta y_{i t}=\xi_{t}+\Delta \mathbf{z}_{i t} \gamma+\delta_{1} \Delta$ prog $_{i t}+\Delta u_{i t}$,\\
where the program participation indicator is differenced along with everything else, and the $\xi_{t}$ are new period intercepts. Example 10.6 is one such case. Extensions of the model, where prog $_{\text {it }}$ appears in other forms, are discussed in Chapter 11.

\subsection*{10.7 Comparison of Estimators}
\subsection*{10.7.1 Fixed Effects versus First Differencing}
When we have only two time periods, FE estimation and FD produce identical estimates and inference, as you are asked to show in Problem 10.3. First differencing is easier to implement, and all procedures that can be applied to a single cross sectionsuch as heteroskedasticity-robust inference-can be applied directly.

When $T>2$ and we are confident the strict exogeneity assumption holds, the choice between FD and FE hinges on the assumptions about the idiosyncratic errors, $u_{i t}$. In particular, the FE estimator is more efficient under Assumption FE.3-the $u_{i t}$ are serially uncorrelated - while the FD estimator is more efficient when $u_{i t}$ follows a random walk. In many cases, the truth is likely to lie somewhere in between.

If FE and FD estimates differ in ways that cannot be attributed to sampling error, we should worry about violations of the strict exogeneity assumption. If $u_{i t}$ is correlated with $\mathbf{x}_{i s}$ for any $t$ and $s$, FE and FD generally have different probability limits. Any of the standard endogeneity problems, including measurement error, timevarying omitted variables, and simultaneity, generally cause correlation between $\mathbf{x}_{i t}$ and $u_{i t}$-that is, contemporaneous correlation - which then causes both FD and FE\\
to be inconsistent and to have different probability limits. (We explicitly consider these problems in Chapter 11.) In addition, correlation between $u_{i t}$ and $\mathbf{x}_{i s}$ for $s \neq t$ causes FD and FE to be inconsistent. When lagged $\mathbf{x}_{i t}$ is correlated with $u_{i t}$, we can solve lack of strict exogeneity by including lags and interpreting the equation as a distributed lag model. More problematical is when $u_{i t}$ is correlated with future $\mathbf{x}_{i t}$ : only rarely does putting future values of explanatory variables in an equation lead to an interesting economic model. In Chapter 11 we show how to estimate the parameters consistently when there is feedback from $u_{i t}$ to $\mathbf{x}_{i s}, s>t$.

If we maintain contemporaneous exogeneity, that is


\begin{equation*}
\mathrm{E}\left(\mathbf{x}_{i t}^{\prime} u_{i t}\right)=\mathbf{0}, \tag{10.74}
\end{equation*}


then we can show that the FE estimator generally has an inconsistency that shrinks to zero at the rate $1 / T$, while the inconsistency of the FD estimator is essentially independent of $T$. More precisely, we can find the probability limits of the FE and FD estimators (with fixed $T$ and $N \rightarrow \infty$ ) and then study the plims as a function of $T$. The inconsistency in these estimators is sometimes, rather loosely, called the "asymptotic bias."

Generally, under Assumption FE.1, we can write


\begin{equation*}
\operatorname{plim}_{N \rightarrow \infty}\left(\hat{\boldsymbol{\beta}}_{F E}\right)=\boldsymbol{\beta}+\left[T^{-1} \sum_{t=1}^{T} \mathrm{E}\left(\ddot{\mathbf{x}}_{i t}^{\prime} \ddot{\mathbf{x}}_{i t}\right)\right]^{-1}\left[T^{-1} \sum_{t=1}^{T} \mathrm{E}\left(\ddot{\mathbf{x}}_{i t}^{\prime} u_{i t}\right)\right]^{-1}, \tag{10.75}
\end{equation*}


where $\ddot{\mathbf{x}}_{i t}=\mathbf{x}_{i t}-\overline{\mathbf{x}}_{i}$, as always, and we emphasize that we are taking the plim for fixed $T$ and $N \rightarrow \infty$. Under (10.74), $\left.\mathrm{E}\left(\ddot{\mathbf{x}}_{i t}^{\prime} u_{i t}\right)=\mathrm{E}\left[\left(\mathbf{x}_{i t}-\overline{\mathbf{x}}_{i}\right)^{\prime} u_{i t}\right)\right]=-\mathrm{E}\left(\overline{\mathbf{x}}_{i}^{\prime} u_{i t}\right)$, and so $T^{-1} \sum_{t=1}^{T} \mathrm{E}\left(\ddot{\mathbf{x}}_{i t}^{\prime} u_{i t}\right)=-T^{-1} \sum_{t=1}^{T} \mathrm{E}\left(\overline{\mathbf{x}}_{i}^{\prime} u_{i t}\right)=-\mathrm{E}\left(\overline{\mathbf{x}}_{i}^{\prime} \bar{u}_{i}\right)$. Now, we can easily bound both average moments in (10.75) if we assume that the process $\left\{\left(\mathbf{x}_{i t}, u_{i t}\right): t=\right. 1,2, \ldots\}$, considered as a time series, is stable and weakly dependent. Actually, as will become clear, what we really should assume is that the time-demeaned sequence, $\left\{\ddot{\mathbf{x}}_{i t}\right\}$, is weakly dependent. This allows for, say, covariate processes such as $\mathbf{x}_{i t}= \mathbf{h}_{i}+\mathbf{r}_{i t}$, where $\mathbf{h}_{i}$ is time-constant heterogeneity and $\left\{\mathbf{r}_{i t}\right\}$ is weakly dependent, because then $\ddot{\mathbf{x}}_{i t}=\ddot{\mathbf{r}}_{i t}$; the persistent component, $\mathbf{h}_{i}$, has been removed. For notational convenience, we just assume that $\left\{\mathbf{x}_{i t}\right\}$ is weakly dependent. (See Hamilton (1994) and Wooldridge (1994a) for general discussions of weak dependence for time series processes. Weakly dependent processes are commonly, if somewhat misleadingly, referred to as "stationary" processes.)

Under weak dependence of the covariate process, $T^{-1} \sum_{t=1}^{T} \mathrm{E}\left(\ddot{\mathbf{x}}_{i t}^{\prime} \ddot{\mathbf{x}}_{i t}\right)$ is bounded as a function of $T$, and so is its inverse if we strengthen Assumption FE. 2 to hold uniformly in $T$-a mild restriction that holds under the rank condition if we impose\\
stationarity on $\left\{\mathbf{x}_{i t}: t=1,2, \ldots\right\}$. Further, $\operatorname{Var}\left(\overline{\mathbf{x}}_{i}\right)$ and $\operatorname{Var}\left(\bar{u}_{i}\right)$ are $O\left(T^{-1}\right)$ under weak dependence. By the Cauchy-Schwartz inequality (for example, Davidson, 1994, Chap. 9) we have, for each $j=1, \ldots, K,\left|\mathrm{E}\left(\bar{x}_{i j} \bar{u}_{i}\right)\right| \leq\left[\operatorname{Var}\left(\bar{x}_{i j}\right) \operatorname{Var}\left(\bar{u}_{i}\right)\right]^{1 / 2}=O\left(T^{-1}\right)$. It follows that


\begin{equation*}
\operatorname{plim}_{N \rightarrow \infty}\left(\hat{\boldsymbol{\beta}}_{F E}\right)=\boldsymbol{\beta}+O(1) \cdot O\left(T^{-1}\right)=\boldsymbol{\beta}+O\left(T^{-1}\right) \equiv \boldsymbol{\beta}+\mathbf{r}_{F E}(T), \tag{10.76}
\end{equation*}


where $\mathbf{r}_{F E}(T)=O\left(T^{-1}\right)$ is the inconsistency as a function of $T$. It follows that the inconsistency in the FE estimator can be small if we have a reasonably large $T$, although the exact size of $\mathbf{r}_{F E}(T)$ depends on features of the stochastic process $\left\{\left(\mathbf{x}_{i t}, u_{i t}\right): t=1,2, \ldots\right\}$.

Hsiao (2003, Sect. 4.2) derives the exact form of $r_{F E}(T)$ (a scalar) for the stable autoregressive model, that is, with $x_{i t}=y_{i, t-1}$ and $|\beta|<1$. The term $r_{F E}(T)$ is negative and, for fixed $T$, increases in absolute value as $\beta$ approaches unity. Unfortunately, finding $\mathbf{r}_{F E}(T)$ in general means modeling $\left\{\left(\mathbf{x}_{i t}, u_{i t}\right): t=1,2, \ldots\right\}$, and so we cannot know how close $\mathbf{r}_{F E}(T)$ is to zero for a given $T$ without making specific assumptions. (In the AR(1) model, we effectively modeled the stochastic process of the regressor and error because the former is the lagged dependent variable and the latter is assumed to be serially uncorrelated.)

For the FD estimator, the general probability limit is


\begin{equation*}
\operatorname{plim}_{N \rightarrow \infty}\left(\hat{\boldsymbol{\beta}}_{F D}\right)=\boldsymbol{\beta}+\left[(T-1)^{-1} \sum_{t=2}^{T} \mathrm{E}\left(\Delta \mathbf{x}_{i t}^{\prime} \Delta \mathbf{x}_{i t}\right)\right]^{-1}\left[(T-1)^{-1} \sum_{t=2}^{T} \mathrm{E}\left(\Delta \mathbf{x}_{i t}^{\prime} \Delta u_{i t}\right)\right]^{-1} . \tag{10.77}
\end{equation*}


If $\left\{\mathbf{x}_{i t}: t=1,2, \ldots\right\}$ is weakly dependent, so is $\Delta \mathbf{x}_{i t}$, and so the first average in (10.77) is generally bounded. (In fact, under stationarity this average does not depend on $T$.) Under (10.74),\\
$\mathrm{E}\left(\Delta \mathbf{x}_{i t}^{\prime} \Delta u_{i t}\right)=-\left[\mathrm{E}\left(\mathbf{x}_{i t}^{\prime} u_{i, t-1}\right)+\mathrm{E}\left(\mathbf{x}_{i, t-1}^{\prime} u_{i t}\right)\right]$,\\
which is generally nonzero. Under stationarity, $\mathrm{E}\left(\Delta \mathbf{x}_{i t}^{\prime} \Delta u_{i t}\right)$ does not depend on $t$, and so the second average in (10.77) does not depend on $T$. Even if we assume $u_{i t}$ is uncorrelated with past covariates-a sequential exogeneity assumption conditional on $c_{i}$-so that $\mathrm{E}\left(\mathbf{x}_{i, t-1}^{\prime} u_{i t}\right)=\mathbf{0}, \mathrm{E}\left(\mathbf{x}_{i t}^{\prime} u_{i, t-1}\right)$ does not equal zero if there is feedback from current idiosyncratic errors to future values of the covariates. Therefore, with sequential exogeneity, the second term is $-(T-1)^{-1} \sum_{t=2}^{T} \mathrm{E}\left(\mathbf{x}_{i t}^{\prime} u_{i, t-1}\right)$, which equals $-\mathrm{E}\left(\mathbf{x}_{i 2}^{\prime} u_{i 1}\right)$ under stationarity, and is generally $O(1)$ even without stationarity.

The previous analysis shows that under contemporaneous exogeneity and weak dependence of the regressors and idiosyncratic errors, the FE estimator has an\\
advantage over the FD estimator when $T$ is large. Interestingly, a more detailed analysis shows that the same order of magnitude holds for $\mathbf{r}_{F E}(T)$ if some regressors have unit roots in their time series representation, that is, they are "integrated of order one," or I(1). (See Hamilton (1994) or Wooldridge (1994a) for more on I(1) processes.) The error process must still be assumed $\mathrm{I}(0)$ (weakly dependent). For example, in the scalar case, if $\left\{x_{i t}: t=1,2, \ldots\right\}$ is $\mathrm{I}(1)$ without a drift term, then $T^{-2} \sum_{t=1}^{T} \mathrm{E}\left[\left(x_{i t}-\bar{x}_{i}\right)^{2}\right]=O(1)$ and $\operatorname{Var}\left(\bar{x}_{i}\right)=O(T)$. If we maintain that $\left\{u_{i t}: t=\right. 1,2, \ldots\}$ is $\mathrm{I}(0)$, then $\operatorname{Var}\left(\bar{u}_{i}\right)=O\left(T^{-1}\right)$, and simple algebra shows that $r_{F E}(T)= O\left(T^{-1}\right)$. The orders of magnitude for moments of averages of $\mathrm{I}(1)$ series can be inferred from Hamilton (1994, Prop. 17.1), for example, or shown directly. Harris and Tzavalis (1999) established the $O\left(T^{-1}\right)$ rate for the inconsistency for the autoregressive model with $\beta=1$-in fact, they show $r_{F E}(T)=-3 /(T+1)$-although the process must be without linear trend when $\beta=1$.

One way to summarize the $O\left(T^{-1}\right)$ inconsistency for $\hat{\beta}_{F E}$ when $-1<\beta \leq 1$ is that it applies to the model $y_{i t}=\beta y_{i, t-1}+(1-\beta) a_{i}+u_{i t}$, so that when $\beta=1,\left\{y_{i t}: t=\right. 0,1, \ldots\}$ is $\mathrm{I}(1)$ without drift (and therefore its mean is $\mathrm{E}\left(y_{i 0}\right)$ for all $t$ ). The FD estimator of $\beta$ has inconsistency $O(1)$ in such models, and more generally with $\mathrm{I}(1)$ regressors if strict exogeneity fails.

The previous analysis certainly favors FE estimation when contemporaneous exogeneity holds but strict exogeneity fails, even if $\left\{y_{i t}\right\}$ and some elements of $\left\{\mathbf{x}_{i t}\right\}$ have unit roots. Unfortunately, there is a catch: the finding that $\mathbf{r}_{F E}(T)=O\left(T^{-1}\right)$ depends critically on the idiosyncratic errors $\left\{u_{i t}\right\}$ being an $\mathrm{I}(0)$ sequence. In the terminology of time series econometrics, $y_{i t}$ and $\mathbf{x}_{i t}$ must be "cointegrated" (see Hamilton, 1994, Chap. 19). If, in a time series sense, our model represents a spurious regression-that is, there is no value of $\boldsymbol{\beta}$ such that $y_{i t}-\mathbf{x}_{i t} \boldsymbol{\beta}-c_{i}$ is $\mathrm{I}(0)$-then the FE approach is no longer superior to FD when comparing inconsistencies. In fact, for fixed $N$, the spurious regression problem for FE becomes more acute as $T$ gets large. By contrast, FD removes any unit roots in $y_{i t}$ and $\mathbf{x}_{i t}$, and so spurious regression is not an issue (but lack of strict exogeneity might be).

When $T$ and $N$ are similar in magnitude, a more realistic scenario is to let $T$ and $N$ grow at the same time (and perhaps the same rate). In this scenario, convergence results for partial sums of $\mathrm{I}(1)$ processes and functions of them are needed to obtain limiting distribution results. Considering large $T$ asymptotics is beyond the scope of this text. Phillips and Moon $(1999,2000)$ discuss unit roots, spurious regression, cointegration, and a variety of estimation and testing procedures for panel data. See Baltagi (2001, Chap. 12) for a summary.

Because strict exogeneity plays such an important role in FE and FD estimation with small $T$, it is important to have a way of formally detecting its violation. One\\
possibility is to directly compare $\hat{\boldsymbol{\beta}}_{F E}$ and $\hat{\boldsymbol{\beta}}_{F D}$ via a Hausman test. It could be important to use a robust form of the test that maintains neither Assumption FE. 3 nor Assumption FD.3, so that neither estimator is assumed efficient under the null. The Hausman test has no systematic power for detecting violation of the second moment assumptions (either FE. 3 or FD.3); it is consistent only against alternatives where $\mathrm{E}\left(\mathbf{x}_{i s}^{\prime} u_{i t}\right) \neq \mathbf{0}$ for some ( $s, t$ ) pairs.

Even if we assume that FE or FD is asymptotically efficient under the null, a drawback to the traditional form of the Hausman test is that, with aggregate time effects, the asymptotic variance of $\sqrt{N}\left(\hat{\boldsymbol{\beta}}_{F E}-\hat{\boldsymbol{\beta}}_{F D}\right)$ is singular. In fact, in a model with only aggregate time effects, it can be shown that the FD and FE estimators are identical. Problem 10.6 asks you to work through the statistic in the case without aggregate time effects. Here, we focus on regression-based tests, which are easy to compute even with time-period dummies and easy to make fully robust.

If $T=2$, it is easy to test for strict exogeneity. In the equation $\Delta y_{i}=\Delta \mathbf{x}_{i} \boldsymbol{\beta}+\Delta u_{i}$, neither $\mathbf{x}_{i 1}$ nor $\mathbf{x}_{i 2}$ should be significant as additional explanatory variables in the FD equation. We simply add, say, $\mathbf{x}_{i 2}$ to the FD equation and carry out an $F$ test for significance of $\mathbf{x}_{i 2}$. With more than two time periods, a test of strict exogeneity is a test of $\mathrm{H}_{0}: \boldsymbol{\gamma}=\mathbf{0}$ in the expanded equation\\
$\Delta y_{i t}=\Delta \mathbf{x}_{i t} \boldsymbol{\beta}+\mathbf{w}_{i t} \boldsymbol{\gamma}+\Delta u_{i t}, \quad t=2, \ldots, T$,\\
where $\mathbf{w}_{i t}$ is a subset of $\mathbf{x}_{i t}$ (that would exclude time dummies). Using the Wald approach, this test can be made robust to arbitrary serial correlation or heteroskedasticity; under Assumptions FD.1-FD. 3 the usual $F$ statistic is asymptotically valid.

A test of strict exogeneity using fixed effects, when $T>2$, is obtained by specifying the equation\\
$y_{i t}=\mathbf{x}_{i t} \boldsymbol{\beta}+\mathbf{w}_{i, t+1} \boldsymbol{\delta}+c_{i}+u_{i t}, \quad t=1,2, \ldots, T-1$,\\
where $\mathbf{w}_{i, t+1}$ is again a subset of $\mathbf{x}_{i, t+1}$. Under strict exogeneity, $\boldsymbol{\delta}=\mathbf{0}$, and we can carry out the test using FE estimation. (We lose the last time period by leading $\mathbf{w}_{i t}$.) An example is given in Problem 10.12.

Under strict exogeneity, we can use a GLS procedure on either the time-demeaned equation or the FD equation. If the variance matrix of $\mathbf{u}_{i}$ is unrestricted, it does not matter which transformation we use. Intuitively, this point is pretty clear, since allowing $\mathrm{E}\left(\mathbf{u}_{i} \mathbf{u}_{i}^{\prime}\right)$ to be unrestricted places no restrictions on $\mathrm{E}\left(\ddot{\mathbf{u}}_{i} \ddot{\mathbf{u}}_{i}^{\prime}\right)$ or $\mathrm{E}\left(\Delta \mathbf{u}_{i} \Delta \mathbf{u}_{i}^{\prime}\right)$. $\operatorname{Im}$, Ahn, Schmidt, and Wooldridge (1999) show formally that the FEGLS and FDGLS estimators are asymptotically equivalent under Assumptions FE. 1 and FEGLS. 3 and\\
the appropriate rank conditions. Of course the system homoskedasticity assumption $\mathrm{E}\left(\mathbf{u}_{i} \mathbf{u}_{i}^{\prime} \mid \mathbf{x}_{i}, c_{i}\right)=\mathrm{E}\left(\mathbf{u}_{i} \mathbf{u}_{i}^{\prime}\right)$ is maintained.

\subsection*{10.7.2 Relationship between the Random Effects and Fixed Effects Estimators}
In cases where the key variables in $\mathbf{x}_{t}$ do not vary much over time, FE and FD methods can lead to imprecise estimates. We may be forced to use random effects estimation in order to learn anything about the population parameters. If an RE analysis is appropriate-that is, if $c_{i}$ is orthogonal to $\mathbf{x}_{i t}$-then the RE estimators can have much smaller variances than the FE or FD estimators. We now obtain an expression for the RE estimator that allows us to compare it with the FE estimator. Using the fact that $\mathbf{j}_{T}^{\prime} \mathbf{j}_{T}=T$, we can write $\boldsymbol{\Omega}$ under the random effects structure as

$$
\begin{aligned}
\boldsymbol{\Omega} & =\sigma_{u}^{2} \mathbf{I}_{T}+\sigma_{c}^{2} \mathbf{j}_{T} \mathbf{j}_{T}^{\prime}=\sigma_{u}^{2} \mathbf{I}_{T}+T \sigma_{c}^{2} \mathbf{j}_{T}\left(\mathbf{j}_{T}^{\prime} \mathbf{j}_{T}\right)^{-1} \mathbf{j}_{T}^{\prime} \\
& =\sigma_{u}^{2} \mathbf{I}_{T}+T \sigma_{c}^{2} \mathbf{P}_{T}=\left(\sigma_{u}^{2}+T \sigma_{c}^{2}\right)\left(\mathbf{P}_{T}+\eta \mathbf{Q}_{T}\right),
\end{aligned}
$$

where $\mathbf{P}_{T} \equiv \mathbf{I}_{T}-\mathbf{Q}_{T}=\mathbf{j}_{T}\left(\mathbf{j}_{T}^{\prime} \mathbf{j}_{T}\right)^{-1} \mathbf{j}_{T}^{\prime}$ and $\eta \equiv \sigma_{u}^{2} /\left(\sigma_{u}^{2}+T \sigma_{c}^{2}\right)$. Next, define $\mathbf{S}_{T} \equiv \mathbf{P}_{T}+\eta \mathbf{Q}_{T}$. Then $\mathbf{S}_{T}^{-1}=\mathbf{P}_{T}+(1 / \eta) \mathbf{Q}_{T}$, as can be seen by direct matrix multiplication. Further, $\mathbf{S}_{T}^{-1 / 2}=\mathbf{P}_{T}+(1 / \sqrt{\eta}) \mathbf{Q}_{T}$, because multiplying this matrix by itself gives $\mathbf{S}_{T}^{-1}$ (the matrix is clearly symmetric, since $\mathbf{P}_{T}$ and $\mathbf{Q}_{T}$ are symmetric). After simple algebra, it can be shown that $\mathbf{S}_{T}^{-1 / 2}=(1-\lambda)^{-1}\left[\mathbf{I}_{T}-\lambda \mathbf{P}_{T}\right]$, where $\lambda=1-\sqrt{\eta}$. Therefore,\\
$\boldsymbol{\Omega}^{-1 / 2}=\left(\sigma_{u}^{2}+T \sigma_{c}^{2}\right)^{-1 / 2}(1-\lambda)^{-1}\left[\mathbf{I}_{T}-\lambda \mathbf{P}_{T}\right]=\left(1 / \sigma_{u}\right)\left[\mathbf{I}_{T}-\lambda \mathbf{P}_{T}\right]$,\\
where $\lambda=1-\left[\sigma_{u}^{2} /\left(\sigma_{u}^{2}+T \sigma_{c}^{2}\right)\right]^{1 / 2}$. Assume for the moment that we know $\lambda$. Then the RE estimator is obtained by estimating the transformed equation $\mathbf{C}_{T} \mathbf{y}_{i}=\mathbf{C}_{T} \mathbf{X}_{i} \boldsymbol{\beta}+ \mathbf{C}_{T} \mathbf{v}_{i}$ by system OLS, where $\mathbf{C}_{T} \equiv\left[\mathbf{I}_{T}-\lambda \mathbf{P}_{T}\right]$. Write the transformed equation as\\
$\breve{\mathbf{y}}_{i}=\breve{\mathbf{X}}_{i} \boldsymbol{\beta}+\breve{\mathbf{v}}_{i}$.

The variance matrix of $\breve{\mathbf{v}}_{i}$ is $\mathrm{E}\left(\breve{\mathbf{v}}_{i} \breve{\mathbf{v}}_{i}^{\prime}\right)=\mathbf{C}_{T} \boldsymbol{\Omega} \mathbf{C}_{T}=\sigma_{u}^{2} \mathbf{I}_{T}$, which verifies that $\breve{\mathbf{v}}_{i}$ has variance matrix ideal for system OLS estimation.

The $t$ th element of $\breve{\mathbf{y}}_{i}$ is easily seen to be $y_{i t}-\lambda \bar{y}_{i}$, and similarly for $\breve{\mathbf{X}}_{i}$. Therefore, system OLS estimation of equation (10.78) is just pooled OLS estimation of\\
$y_{i t}-\lambda \bar{y}_{i}=\left(\mathbf{x}_{i t}-\lambda \overline{\mathbf{x}}_{i}\right) \boldsymbol{\beta}+\left(v_{i t}-\lambda \bar{v}_{i}\right)$\\
over all $t$ and $i$. The errors in this equation are serially uncorrelated and homoskedastic under Assumption RE.3; therefore, they satisfy the key conditions for pooled OLS analysis. The feasible RE estimator replaces the unknown $\lambda$ with its estimator, $\hat{\lambda}$, so that $\hat{\boldsymbol{\beta}}_{R E}$ can be computed from the pooled OLS regression\\
$\breve{y}_{i t}$ on $\breve{\mathbf{x}}_{i t}, \quad t=1, \ldots, T ; i=1, \ldots, N$,\\
where now $\breve{\mathbf{x}}_{i t}=\mathbf{x}_{i t}-\hat{\lambda} \overline{\mathbf{x}}_{i}$ and $\breve{y}_{i t}=y_{i t}-\hat{\lambda} \bar{y}_{i}$, all $t$ and $i$. Therefore, we can write\\
$\hat{\boldsymbol{\beta}}_{R E}=\left(\sum_{i=1}^{N} \sum_{t=1}^{T} \breve{\mathbf{x}}_{i t}^{\prime} \breve{\mathbf{x}}_{i t}\right)^{-1}\left(\sum_{i=1}^{N} \sum_{t=1}^{T} \breve{\mathbf{x}}_{i t}^{\prime} \breve{y}_{i t}\right)$.

The usual variance estimate from the pooled OLS regression (10.79), SSR/( NT - K), is a consistent estimator of $\sigma_{u}^{2}$. The usual $t$ statistics and $F$ statistics from the pooled regression are asymptotically valid under Assumptions RE.1-RE.3. For $F$ tests, we obtain $\hat{\lambda}$ from the unrestricted model.

Equation (10.80) shows that the RE estimator is obtained by a quasi-time demeaning: rather than removing the time average from the explanatory and dependent variables at each $t$, RE estimation removes a fraction of the time average. If $\hat{\lambda}$ is close to unity, the RE and FE estimates tend to be close. To see when this result occurs, write $\hat{\lambda}$ as\\
$\hat{\lambda}=1-\left\{1 /\left[1+T\left(\hat{\sigma}_{c}^{2} / \hat{\sigma}_{u}^{2}\right)\right]\right\}^{1 / 2}$,\\
where $\hat{\sigma}_{u}^{2}$ and $\hat{\sigma}_{c}^{2}$ are consistent estimators of $\sigma_{u}^{2}$ and $\sigma_{c}^{2}$ (see Section 10.4). When $T\left(\hat{\sigma}_{c}^{2} / \hat{\sigma}_{u}^{2}\right)$ is large, the second term in $\hat{\lambda}$ is small, in which case $\hat{\lambda}$ is close to unity. In fact, $\hat{\lambda} \rightarrow 1$ as $T \rightarrow \infty$ or as $\hat{\sigma}_{c}^{2} / \hat{\sigma}_{u}^{2} \rightarrow \infty$. For large $T$, it is not surprising to find similar estimates from FE and RE . Even with small $T, \mathrm{RE}$ can be close to FE if the estimated variance of $c_{i}$ is large relative to the estimated variance of $u_{i t}$, a case often relevant for applications. (As $\lambda$ approaches unity, the precision of the RE estimator approaches that of the FE estimator, and the effects of time-constant explanatory variables become harder to estimate.)

Example 10.7 (Job Training Grants): In Example 10.4, $T=3, \hat{\sigma}_{u}^{2} \approx .248$, and $\hat{\sigma}_{c}^{2} \approx 1.932$, which gives $\hat{\lambda} \approx .797$. This helps explain why the RE and FE estimates are reasonably close.

Equations (10.80) and (10.81) also show how RE and pooled OLS are related. Pooled OLS is obtained by setting $\hat{\lambda}=0$, which is never exactly true but could be close. In practice, $\hat{\lambda}$ is not usually close to zero because small values require $\hat{\sigma}_{u}^{2}$ to be large relative to $\hat{\sigma}_{c}^{2}$.

In Section 10.4 we emphasized that consistency of RE estimation hinges on the orthogonality between $c_{i}$ and $\mathbf{x}_{i t}$. In fact, Assumption POLS. 1 is weaker than Assumption RE.1. We now see, because of the particular transformation used by the

RE estimator, that its inconsistency when Assumption RE.1b is violated can be small relative to pooled OLS if $\sigma_{c}^{2}$ is large relative to $\sigma_{u}^{2}$ or if $T$ is large.

If we are primarily interested in the effect of a time-constant variable in a panel data study, the robustness of the FE estimator to correlation between the unobserved effect and the $\mathbf{x}_{i t}$ is practically useless. Without using an instrumental variables approach-something we take up in Chapter 11-RE estimation is probably our only choice. Sometimes, applications of the RE estimator attempt to control for the part of $c_{i}$ correlated with $\mathbf{x}_{i t}$ by including dummy variables for various groups, assuming that we have many observations within each group. For example, if we have panel data on a group of working people, we might include city dummy variables in a wage equation. Or, if we have panel data at the student level, we might include school dummy variables. Including dummy variables for groups controls for a certain amount of heterogeneity that might be correlated with the (time-constant) elements of $\mathbf{x}_{i t}$. By using RE, we can efficiently account for any remaining serial correlation due to unobserved time-constant factors. (Unfortunately, the language used in empirical work can be confusing. It is not uncommon to see school dummy variables referred to as "school fixed effects," even though they appear in an RE analysis at the individual level.)

Regression (10.79) using the quasi-time-demeaned data has several other practical uses. Since it is just a pooled OLS regression that is asymptotically the same as using $\lambda$ in place of $\hat{\lambda}$, we can easily obtain standard errors that are robust to arbitrary heteroskedasticity in $c_{i}$ and $u_{i t}$, as well as arbitrary serial correlation in the $\left\{u_{i t}\right\}$. All that is required is an econometrics package that computes robust standard errors, $t$, and $F$ statistics for pooled OLS regression, such as Stata. Further, we can use the residuals from regression (10.79), say $\hat{r}_{i t}$, to test for serial correlation in $r_{i t} \equiv v_{i t}-\lambda \bar{v}_{i}$, which are serially uncorrelated under Assumption RE.3a. If we detect serial correlation in $\left\{r_{i t}\right\}$, we conclude that Assumption RE.3a is false, which means that the $u_{i t}$ are serially correlated. Although the arguments are tedious, it can be shown that estimation of $\lambda$ and $\boldsymbol{\beta}$ has no effect on the null limiting distribution of the usual (or heteroskedasticity-robust) $t$ statistic from the pooled OLS regression $\hat{r}_{i t}$ on $\hat{r}_{i, t-1}$, $t=2, \ldots, T, i=1, \ldots, N$.

\subsection*{10.7.3 Hausman Test Comparing Random Effects and Fixed Effects Estimators}
Because the key consideration in choosing between an RE and an FE approach is whether $c_{i}$ and $\mathbf{x}_{i t}$ are correlated, it is important to have a method for testing this assumption. Hausman (1978) proposed a test based on the difference between the RE and FE estimates. Since FE is consistent when $c_{i}$ and $\mathbf{x}_{i t}$ are correlated, but RE is\\
inconsistent, a statistically significant difference is interpreted as evidence against the RE assumption RE.1b. This interpretation implicitly assumes that $\left\{\mathbf{x}_{i t}\right\}$ is strictly exogenous with respect to $\left\{u_{i t}\right\}$, that is, Assumption RE.1a (FE.1) holds.

Before we obtain the Hausman test, there are three caveats. First, strict exogeneity, Assumption RE.1a, is maintained under the null and the alternative. Correlation between $\mathbf{x}_{i s}$ and $u_{i t}$ for any $s$ and $t$ causes both FE and RE to be inconsistent, and generally their plims will differ. In order to view the Hausman test as a test of $\operatorname{Cov}\left(\mathbf{x}_{i t}, c_{i}\right)=0$, we must maintain Assumption RE.1a.

A second caveat is that the test is usually implemented assuming that Assumption RE. 3 holds under the null. As we will see, this setup implies that the RE estimator is more efficient than the FE estimator, and it simplifies computation of the test statistic. But we must emphasize that Assumption RE. 3 is an auxiliary assumption, and it is not being tested by the Hausman statistic: the Hausman test has no systematic power against the alternative that Assumption RE. 1 is true but Assumption RE. 3 is false. Failure of Assumption RE. 3 causes the usual Hausman test to have a nonstandard limiting distribution, an issue we return to in a later discussion.

A third caveat concerns the set of parameters that we can compare. Because the FE approach only identifies coefficients on time-varying explanatory variables, we clearly cannot compare FE and RE coefficients on time-constant variables. But there is a more subtle issue: we cannot include in our comparison coefficients on aggregate time effects-that is, variables that change only across $t$. As with the case of comparing FE and FD estimates, the problem with comparing coefficients on aggregate time effects is not one of identification; we know RE and FE both allow inclusion of a full set of time period dummies. The problem is one of singularity in the asymptotic variance matrix of the difference between $\hat{\boldsymbol{\beta}}_{F E}$ and $\hat{\boldsymbol{\beta}}_{R E}$. Problem 10.17 asks you to show that in the model $y_{i t}=\alpha+\mathbf{d}_{t} \boldsymbol{\eta}+\mathbf{w}_{i t} \boldsymbol{\delta}+c_{i}+u_{i t}$, where $\mathbf{d}_{t}$ is a $1 \times R$ vector of aggregate time effects and $\mathbf{w}_{i t}$ is a $1 \times M$ vector of regressors varying across $i$ and $t$, the asymptotic variance of the difference in the FE and RE estimators of $\left(\boldsymbol{\eta}^{\prime}, \boldsymbol{\delta}^{\prime}\right)^{\prime}$ has rank $M$, not $R+M$. (In fact, without $\mathbf{w}_{i t}$ in the model, the FE and RE estimates of $\boldsymbol{\eta}$ are identical. Problem 10.18 asks you to investigate this algebraic result and some related claims for a particular data set.)

Rather than consider the general case-which is complicated by the singularity of the asymptotic covariance matrix-we assume that there are no aggregate time effects in the model in deriving the traditional form of the Hausman test. Therefore, we write the equation as $y_{i t}=\mathbf{x}_{i t} \boldsymbol{\beta}+c_{i}+u_{i t}=\mathbf{z}_{i} \boldsymbol{\gamma}+\mathbf{w}_{i t} \boldsymbol{\delta}+c_{i}+u_{i t}$, where $\mathbf{z}_{i}$ (say, $1 \times J)$ includes at least an intercept and often other time-constant variables. The elements of the $1 \times M$ vector $\mathbf{w}_{i t}$ vary across $i$ and $t$ (as usual, at least for some units\\
and some time periods), and it is the FE and RE estimators of $\boldsymbol{\delta}$ that we compare. For deriving the traditional form of the test, we maintain Assumptions RE. 1 to RE. 3 under the null, as well as the rank condition for fixed effects, Assumption FE.2.

A key component of the traditional Hausman test is showing that the asymptotic variance of the FE estimator is never smaller, and is usually strictly larger, than the asymptotic variance of the RE estimator. We already know the asymptotic variance of the FE estimator:\\
$\operatorname{Avar}\left(\hat{\boldsymbol{\delta}}_{F E}\right)=\sigma_{u}^{2}\left[\mathrm{E}\left(\ddot{\mathbf{W}}_{i}^{\prime} \ddot{\mathbf{W}}_{i}\right)\right]^{-1} / N$.

In the presence of $\mathbf{z}_{i}$ the derivation of $\operatorname{Avar}\left(\hat{\boldsymbol{\delta}}_{R E}\right)$ is a little more tedious, but simplified by the pooled OLS characterization of RE using the quasi-time-demeaned data. Define $\breve{\mathbf{w}}_{i t}=\mathbf{w}_{i t}-\lambda \overline{\mathbf{w}}_{i}$ as the quasi-time demeaned time-varying regressors. To get $\operatorname{Avar}\left(\hat{\boldsymbol{\delta}}_{R E}\right)$, we must obtain the population residuals from the pooled regression $\breve{\mathbf{w}}_{i t}$ on $(1-\lambda) \mathbf{z}_{i}$, which, of course, is the same as dropping the $(1-\lambda)$. Call these population residuals $\tilde{\mathbf{w}}_{i t}$. Then\\
$\operatorname{Avar}\left(\hat{\boldsymbol{\delta}}_{R E}\right)=\sigma_{u}^{2}\left[\mathrm{E}\left(\tilde{\mathbf{W}}_{i}^{\prime} \tilde{\mathbf{W}}_{i}\right)\right]^{-1} / N$.

Next, we want to show that $\operatorname{Avar}\left(\hat{\boldsymbol{\delta}}_{F E}\right)-\operatorname{Avar}\left(\hat{\boldsymbol{\delta}}_{R E}\right)$ is positive definite. But this holds if $\left[\operatorname{Avar}\left(\hat{\boldsymbol{\delta}}_{R E}\right)\right]^{-1}-\left[\operatorname{Avar}\left(\hat{\boldsymbol{\delta}}_{F E}\right)\right]^{-1}$ is positive definite, or $\mathrm{E}\left(\tilde{\mathbf{W}}_{i}^{\prime} \tilde{\mathbf{W}}_{i}\right)-\mathrm{E}\left(\ddot{\mathbf{W}}_{i}^{\prime} \ddot{\mathbf{W}}_{i}\right)$ is positive definite. To demonstrate the latter result, we need to more carefully characterize $\tilde{\mathbf{w}}_{i t}=\breve{\mathbf{w}}_{i t}-\mathbf{z}_{i} \boldsymbol{\Pi}$, where $\boldsymbol{\Pi}=\left[T \cdot \mathrm{E}\left(\mathbf{z}_{i}^{\prime} \mathbf{z}_{i}\right)\right]^{-1} \mathrm{E}\left[\mathbf{z}_{i}^{\prime}\left(\sum_{t=1}^{T} \breve{\mathbf{w}}_{i t}\right)\right]=(1-\lambda)\left[\mathrm{E}\left(\mathbf{z}_{i}^{\prime} \mathbf{z}_{i}\right)\right]^{-1}$. $\mathrm{E}\left(\mathbf{z}_{i}^{\prime} \overline{\mathbf{w}}_{i}\right)$. Straightforward algebra gives $\tilde{\mathbf{w}}_{i t}=\breve{\mathbf{w}}_{i t}-(1-\lambda) \overline{\mathbf{w}}_{i}^{*}$, where $\overline{\mathbf{w}}_{i}^{*}=\mathrm{L}\left(\overline{\mathbf{w}}_{i} \mid \mathbf{z}_{i}\right)$ is the linear projection of $\overline{\mathbf{w}}_{i}$ on $\mathbf{z}_{i}$. We can also write\\
$\tilde{\mathbf{w}}_{i t}=\ddot{\mathbf{w}}_{i t}+(1-\lambda)\left(\overline{\mathbf{w}}_{i}-\overline{\mathbf{w}}_{i}^{*}\right)$,\\
from which it follows immediately


\begin{align*}
\mathrm{E}\left(\tilde{\mathbf{W}}_{i}^{\prime} \tilde{\mathbf{W}}_{i}\right)-\mathrm{E}\left(\ddot{\mathbf{W}}_{i}^{\prime} \ddot{\mathbf{W}}_{i}\right)= & (1-\lambda)^{2} \mathrm{E}\left[\left(\overline{\mathbf{w}}_{i}-\overline{\mathbf{w}}_{i}^{*}\right)^{\prime}\left(\overline{\mathbf{w}}_{i}-\overline{\mathbf{w}}_{i}^{*}\right)\right] \\
& +(1-\lambda) \sum_{t=1}^{T} \ddot{\mathbf{w}}_{i t}^{\prime}\left(\overline{\mathbf{w}}_{i}-\overline{\mathbf{w}}_{i}^{*}\right)+(1-\lambda) \sum_{t=1}^{T}\left(\overline{\mathbf{w}}_{i}-\overline{\mathbf{w}}_{i}^{*}\right)^{\prime} \ddot{\mathbf{w}}_{i t} \\
= & (1-\lambda)^{2} \mathrm{E}\left[\left(\overline{\mathbf{w}}_{i}-\overline{\mathbf{w}}_{i}^{*}\right)^{\prime}\left(\overline{\mathbf{w}}_{i}-\overline{\mathbf{w}}_{i}^{*}\right)\right] \tag{10.85}
\end{align*}


because $\sum_{t=1}^{T} \ddot{\mathbf{w}}_{i t}=\mathbf{0}$ for all $i$. For $\lambda<1$, and provided $\overline{\mathbf{w}}_{i}-\mathrm{L}\left(\overline{\mathbf{w}}_{i} \mid \mathbf{z}_{i}\right)$ has variance with full rank-which means the time averages of the time-varying regressors are not perfectly collinear with the time-constant regressors-we have shown that $\mathrm{E}\left(\tilde{\mathbf{W}}_{i}^{\prime} \tilde{\mathbf{W}}_{i}\right)-\mathrm{E}\left(\ddot{\mathbf{W}}_{i}^{\prime} \ddot{\mathbf{W}}_{i}\right)$ is positive definite.

It can also be shown that, under RE. 1 to RE.3, $\operatorname{Avar}\left(\hat{\boldsymbol{\delta}}_{F E}-\hat{\boldsymbol{\delta}}_{R E}\right)=\operatorname{Avar}\left(\hat{\boldsymbol{\delta}}_{F E}\right)- \operatorname{Avar}\left(\hat{\boldsymbol{\delta}}_{R E}\right)$; Newey and McFadden (1994, Sect. 5.3) provide general sufficient conditions, which are met by the FE and RE estimators under Assumptions RE.1-RE.3. (We cover these conditions in Chapter 14 in our discussion of general efficiency issues; see Lemma 14.1 and the surrounding discussion.) Therefore, we can compute the Hausman statistic as\\
$H=\left(\hat{\boldsymbol{\delta}}_{F E}-\hat{\boldsymbol{\delta}}_{R E}\right)^{\prime}\left[\operatorname{Avar(\hat {\boldsymbol {\delta }}_{FE})}-\widehat{\operatorname{Avar}\left(\hat{\boldsymbol{\delta}}_{R E}\right)}\right]^{-1}\left(\hat{\boldsymbol{\delta}}_{F E}-\hat{\boldsymbol{\delta}}_{R E}\right)$,\\
which has a $\chi_{M}^{2}$ distribution under the null hypothesis. The usual estimates of $\operatorname{Avar}\left(\hat{\boldsymbol{\delta}}_{F E}\right)$ and $\operatorname{Avar}\left(\hat{\boldsymbol{\delta}}_{R E}\right)$ can be used in equation (10.86), but if different estimates of $\sigma_{u}^{2}$ are used, the matrix in the middle need not be positive definite, possibly leading to a negative value of $H$. It is best to use the same estimate of $\sigma_{u}^{2}$ (based on either FE or RE) in both places.

If we are interested in testing the difference in estimators for a single coefficient, we can use a $t$ statistic form of the Hausman test. In particular, if $\delta$ now denotes a scalar on the time-varying variable of interest, we can use $\left(\hat{\delta}_{F E}-\hat{\delta}_{R E}\right) /\left\{\left[\operatorname{se}\left(\hat{\delta}_{F E}\right)\right]^{2}-\right. \left.\left[\operatorname{se}\left(\hat{\delta}_{R E}\right)\right]^{2}\right\}^{1 / 2}$, provided that we use versions of the standard errors that ensure $\operatorname{se}\left(\hat{\delta}_{F E}\right)>\operatorname{se}\left(\hat{\delta}_{R E}\right)$. Under Assumptions RE.1-RE.3, the $t$ statistic has an asymptotic standard normal distribution.

So far we have stated that the null hypothesis is RE.1-RE.3, but expression (10.84) allows us to characterize the implicit null hypothesis. From (10.84), it is seen that deviations between the RE and FE estimates of $\boldsymbol{\delta}$ are due to correlation between $\overline{\mathbf{w}}_{i}-\overline{\mathbf{w}}_{i}^{*}$ and $c_{i}$ (where we maintain Assumption RE.1a, which is strict exogeneity conditional on $c_{i}$ ). In other words, the Hausman test is a test of\\
$\mathrm{H}_{0}: \mathrm{E}\left[\left(\overline{\mathbf{w}}_{i}-\overline{\mathbf{w}}_{i}^{*}\right)^{\prime} c_{i}\right]=\mathbf{0}$.

Equation (10.87) is interesting for several reasons. First, if there are no time-constant variables (except an overall intercept) in the RE estimation, the null hypothesis is $\operatorname{Cov}\left(\overline{\mathbf{w}}_{i}, c_{i}\right)=\mathbf{0}$, which means we are really testing whether the time-average of the $\mathbf{w}_{i t}$ is correlated with the unobserved effect. With time-constant explanatory variables $\mathbf{z}_{i}$, we first remove the correlation between $\overline{\mathbf{w}}_{i}$ and $\mathbf{z}_{i}$ to form the population residuals, $\overline{\mathbf{w}}_{i}-\overline{\mathbf{w}}_{i}^{*}$, before testing for correlation with $c_{i}$. An immediate consequence is that, with a rich set of controls in $\mathbf{z}_{i}$, it is possible for $\overline{\mathbf{w}}_{i}-\overline{\mathbf{w}}_{i}^{*}$ to be uncorrelated with $c_{i}$ even though $\overline{\mathbf{w}}_{i}$ is correlated with $c_{i}$. Not surprisingly, an RE analysis with good controls in $\mathbf{z}_{i}$ and an RE analysis that omits such controls can yield very different estimates of $\boldsymbol{\delta}$, and the RE estimate of $\boldsymbol{\delta}$ with $\mathbf{z}_{i}$ included might be much closer to the FE estimate than if $\mathbf{z}_{i}$ is excluded. Often, discussions of the Hausman test\\
comparing FE and RE assume that only time-varying explanatory variables are included in the RE estimation, but that is too restrictive for applications with timeconstant controls.

Equation (10.87) also suggests a simple regression-based approach to computing a Hausman statistic. In fact, we are led to a regression-based method if we use a particular correlated RE assumption due to Mundlak (1978): $c_{i}=\psi+\overline{\mathbf{w}}_{i} \boldsymbol{\xi}+a_{i}$, where $a_{i}$ has zero mean and is assumed to be uncorrelated with $\mathbf{w}_{i}=\left(\mathbf{w}_{i 1}, \ldots, \mathbf{w}_{i T}\right)$ and $\mathbf{z}_{i}$. Plugging in for $c_{i}$ gives the expanded equation\\
$y_{i t}=\mathbf{x}_{i t} \boldsymbol{\beta}+\overline{\mathbf{w}}_{i} \boldsymbol{\xi}+a_{i}+u_{i t}$,\\
where we absorb $\psi$ into $\boldsymbol{\beta}$ because we assume $\mathbf{x}_{i t}$ includes an intercept. In fact, in addition to an intercept, time-constant variables $\mathbf{z}_{i}$, and $\mathbf{w}_{i t}, \mathbf{x}_{i t}$ can (and usually should) contain a full set of time period dummies. Mundlak (1978) suggested testing $\mathrm{H}_{0}: \boldsymbol{\xi}=\mathbf{0}$ to determine if the heterogeneity was correlated with the time averages of the $\mathbf{w}_{i t}$. (Incidentally, the formulation in equation (10.88) makes it clear that we cannot include the time average of aggregate time variables in $\overline{\mathbf{w}}_{i}$ because the time averages simply would be constant across all $i$.)

How should we estimate the parameters in equation (10.88)? We could estimate the equation by pooled OLS or we could estimate the equation by random effects, which is asymptotically efficient under RE.1-RE.3. As it turns out, the pooled OLS and RE estimates of $\boldsymbol{\xi}$ are identical, and $\hat{\boldsymbol{\xi}}=\hat{\boldsymbol{\delta}}_{B}-\hat{\boldsymbol{\delta}}_{F E}$, where $\hat{\boldsymbol{\delta}}_{B}$ (the between estimator) is the coefficient vector on $\overline{\mathbf{w}}_{i}$ from the cross section regression $\bar{y}_{i}$ on $\mathbf{z}_{i}, \overline{\mathbf{w}}_{i}, i=1, \ldots, N$. Further, the coefficient vector on $\mathbf{w}_{i t}$ is simply the FE estimate; see Mundlak (1978) and Hsiao (2003, Sect. 3.2) for verification. In other words, the regression-based version of the test explicitly compares the between estimate and the FE estimate of $\boldsymbol{\delta}$. Hausman and Taylor (1981) show that this is the same as basing the test on the RE and FE estimate because the RE estimator is a matrix-weighted average of the between and FE estimators; see also Baltagi (2001, Sect. 2.3).

That the FE estimator of $\boldsymbol{\delta}$-the coefficient on $\mathbf{w}_{i t}$-is obtained as the RE estimator in equation (10.88) sheds further light on the source of efficiency of RE over FE. In effect, the RE estimator of $\boldsymbol{\delta}$ sets $\boldsymbol{\xi}$ to zero. Because $\overline{\mathbf{w}}_{i}$ is correlated with $\mathbf{w}_{i t}$-often highly correlated-dropping it from the equation when it is legimate to do so increases efficiency of the remaining parameter estimates by reducing multicollinearity. Of course, inappropriately setting $\boldsymbol{\xi}$ to zero results in inconsistent estimation of $\boldsymbol{\delta}$, and that is the danger of the RE approach.

If we use pooled OLS to test $H_{0}: \boldsymbol{\xi}=\mathbf{0}$, the usual pooled OLS test statistic will be inappropriate-at a minimum because of the serial correlation induced by $a_{i}$. As we know, it is easy to obtain a fully robust variance matrix estimator, and therefore a\\
fully robust Wald statistic, for pooled OLS. Such a statistic will be fully robust to violations of Assumption RE.3. Alternatively, we can estimate (10.88) by RE, and, if we maintain Assumption RE.3, we can use a standard Wald test computed using the usual RE variance matrix estimator. We also know it is easy to obtain a fully robust Wald statistic, that is, a statistic that does not maintain Assumption RE. 3 under the null. The robust Wald statistic based on the RE estimator is necessarily asymptotically equivalent to the robust Wald statistic for pooled OLS.

Earlier we discussed how the traditional Hausman test maintains Assumption RE. 3 under the null. Why should we make the Hausman robust to violations to Assumption RE.3? There is some confusion on this point in the methodolgical and empirical literatures. Hausman (1978) originally presented his testing principle as applying to situations where one estimator is efficient under the null but inconsistent under the alternative-the RE estimator in this case-and the other estimator is consistent under the null and the alternative but inefficient under the null-the FE estimator. While this scenario simplifies calculation of the test statistic-see equation (10.86)-it by no means is required for the test to make sense. In the current context, whether or not the covariates are correlated with $c_{i}$ (through the time average) has nothing to do with conditional second moment assumptions on $c_{i}$ or $\mathbf{u}_{i}=\left(u_{i 1}, \ldots, u_{i T}\right)^{\prime}$-that is, on whether Assumption RE. 3 holds. In fact, including Assumption RE. 3 in the null hypothesis serves to mask the kinds of misspecification the traditional Hausman test can detect. Certainly the test has power against violations of (10.87)-this is what it is intended to have power against. But the nonrobust Hausman statistic, whether computed in (10.86) or from (10.88)-they are asymptotically equivalent under RE.3-has no systematic power for detecting violation of Assumption RE.3. Specifically, if (10.87) holds and the standard rank conditions hold for FE and RE, the statistic converges in distribution (rather than diverging) whether or not Assumption RE. 3 holds. In other words, the test is inconsistent for testing RE.3. This is easily seen from (10.86). The quadratic form will converge in distribution to a quadratic form in a multivariate normal. If Assumption RE. 3 holds, that quadratic form has a chi-square distribution with $M$ degrees of freedom, but not otherwise. Without Assumption RE.3, using the $\chi_{M}^{2}$ distribution to obtain critical values, or $p$-values, will result in a test that will be undersized or oversized under (10.87)-and we cannot generally tell which is the case. This feature carries over to any Hausman statistic where an auxiliary assumption is maintained that means one estimator is asymptotically efficient under the null. See Wooldridge (1991b) for further discussion in the context of testing conditional mean specifications.

To summarize, we can estimate models that include aggregate time effects, timeconstant variables, and regressors that change across both $i$ and $t$, by RE and FE\\
estimation. But no matter how we compute a test statistic, we can only compare the coefficients on the regressors that change across both $i$ and $t$. The regression-based version of the test in equation (10.88) makes it easy to obtain a statistic with a nondegenerate asymptotic distribution, and it is also easy to make the regression-based test fully robust to violations of Assumption RE.3.

As in any other context that uses statistical inference, it is possible to get a statistical rejection of RE.1b (say, at the 5 percent level) with the differences between the RE and FE estimates being practically small. The opposite case is also possible: there can be seemingly large differences between the RE and FE estimates but, due to large standard errors, the Hausman statistic fails to reject. What should be done in this case? A typical response is to conclude that the random effects assumptions hold and to focus on the RE estimates. Unfortunately, we may be committing a Type II error: failing to reject Assumption RE.1b when it is false.

\section*{Problems}
10.1. Consider a model for new capital investment in a particular industry (say, manufacturing), where the cross section observations are at the county level and there are $T$ years of data for each county:\\
$\log \left(\right.$ invest $\left._{i t}\right)=\theta_{t}+\mathbf{z}_{i t} \gamma+\delta_{1} \operatorname{tax}_{i t}+\delta_{2} \operatorname{disaster}_{i t}+c_{i}+u_{i t}$.\\
The variable $\operatorname{tax}_{i t}$ is a measure of the marginal tax rate on capital in the county, and disaster $_{i t}$ is a dummy indicator equal to one if there was a significant natural disaster in county $i$ at time period $t$ (for example, a major flood, a hurricane, or an earthquake). The variables in $\mathbf{z}_{i t}$ are other factors affecting capital investment, and the $\theta_{t}$ represent different time intercepts.\\
a. Why is allowing for aggregate time effects in the equation important?\\
b. What kinds of variables are captured in $c_{i}$ ?\\
c. Interpreting the equation in a causal fashion, what sign does economic reasoning suggest for $\delta_{1}$ ?\\
d. Explain in detail how you would estimate this model; be specific about the assumptions you are making.\\
e. Discuss whether strict exogeneity is reasonable for the two variables $\operatorname{tax}_{i t}$ and disaster $i t$; assume that neither of these variables has a lagged effect on capital investment.\\
10.2. Suppose you have $T=2$ years of data on the same group of $N$ working individuals. Consider the following model of wage determination:\\
$\log \left(\right.$ wage $\left._{i t}\right)=\theta_{1}+\theta_{2} d 2_{t}+\mathbf{z}_{i t} \gamma+\delta_{1}$ female $_{i}+\delta_{2} d 2_{t} \cdot$ female $_{i}+c_{i}+u_{i t}$.\\
The unobserved effect $c_{i}$ is allowed to be correlated with $\mathbf{z}_{i t}$ and female ${ }_{i}$. The variable $d 2_{t}$ is a time period indicator, where $d 2_{t}=1$ if $t=2$ and $d 2_{t}=0$ if $t=1$. In what follows, assume that\\
$\mathrm{E}\left(u_{i t} \mid\right.$ female $\left._{i}, \mathbf{z}_{i 1}, \mathbf{z}_{i 2}, c_{i}\right)=0, \quad t=1,2$.\\
a. Without further assumptions, what parameters in the log wage equation can be consistently estimated?\\
b. Interpret the coefficients $\theta_{2}$ and $\delta_{2}$.\\
c. Write the log wage equation explicitly for the two time periods. Show that the differenced equation can be written as\\
$\Delta \log \left(\right.$ wage $\left._{i}\right)=\theta_{2}+\Delta \mathbf{z}_{i} \gamma+\delta_{2}$ female $_{i}+\Delta u_{i}$, where $\Delta \log \left(\right.$ wage $\left._{i}\right)=\log \left(\right.$ wage $\left._{i 2}\right)-\log \left(\right.$ wage $\left._{i 1}\right)$, and so on.\\
d. How would you test $\mathrm{H}_{0}: \delta_{2}=0$ if $\operatorname{Var}\left(\Delta u_{i} \mid \Delta \mathbf{z}_{i}\right.$, female $\left.e_{i}\right)$ is not constant?\\
10.3. For $T=2$ consider the standard unoberved effects model\\
$y_{i t}=\mathbf{x}_{i t} \boldsymbol{\beta}+c_{i}+u_{i t}, \quad t=1,2$.\\
Let $\hat{\boldsymbol{\beta}}_{F E}$ and $\hat{\boldsymbol{\beta}}_{F D}$ denote the fixed effects and first difference estimators, respectively.\\
a. Show that the FE and FD estimates are numerically identical.\\
b. Show that the error variance estimates from the FE and FD methods are numerically identical.\\
10.4. A common setup for program evaluation with two periods of panel data is the following. Let $y_{i t}$ denote the outcome of interest for unit $i$ in period $t$. At $t=1$, no one is in the program; at $t=2$, some units are in the control group and others are in the experimental group. Let prog $_{\text {it }}$ be a binary indicator equal to one if unit $i$ is in the program in period $t$; by the program design, $\operatorname{prog}_{i 1}=0$ for all $i$. An unobserved effects model without additional covariates is\\
$y_{i t}=\theta_{1}+\theta_{2} d 2_{t}+\delta_{1}$ prog $_{i t}+c_{i}+u_{i t}, \quad \mathrm{E}\left(u_{i t} \mid\right.$ prog $\left._{i 2}, c_{i}\right)=0$,\\
where $d 2_{t}$ is a dummy variable equal to unity if $t=2$, and zero if $t=1$, and $c_{i}$ is the unobserved effect.\\
a. Explain why including $d 2_{t}$ is important in these contexts. In particular, what problems might be caused by leaving it out?\\
b. Why is it important to include $c_{i}$ in the equation?\\
c. Using the first differencing method, show that $\hat{\theta}_{2}=\overline{\Delta y}_{\text {control }}$ and $\hat{\delta}_{1}=\overline{\Delta y}_{\text {treat }}- \overline{\Delta y}_{\text {control }}$, where $\overline{\Delta y}_{\text {control }}$ is the average change in $y$ over the two periods for the group with prog $_{i 2}=0$, and $\overline{\Delta y}_{\text {treat }}$ is the average change in $y$ for the group where prog $_{i 2}=$ 1. This formula shows that $\hat{\delta}_{1}$, the difference-in-differences estimator, arises out of an unobserved effects panel data model.\\
d. Write down the extension of the model for $T$ time periods.\\
e. A common way to obtain the DD estimator for two years of panel data is from the model\\
$y_{i t}=\alpha_{1}+\alpha_{2}$ start $_{t}+\alpha_{3}$ prog $_{i}+\delta_{1}$ start $_{t}$ prog $_{i}+u_{i t}$,\\
where $\mathrm{E}\left(u_{i t} \mid\right.$ start $_{t}$, prog $\left._{i}\right)=0$, prog $_{i}$ denotes whether unit $i$ is in the program in the second period, and start ${ }_{t}$ is a binary variable indicating when the program starts. In the two-period setup, start $=d 2_{t}$ and prog ${ }_{\text {it }}=$ start $_{t}$ prog $_{i}$. The pooled OLS estimator of $\delta_{1}$ is the DD estimator from part c. With $T>2$, the unobserved effects model from part d and pooled estimation of equation (10.89) no longer generally give the same estimate of the program effect. Which approach do you prefer, and why?\\
10.5. Assume that Assumptions RE. 1 and RE.3a hold, but $\operatorname{Var}\left(c_{i} \mid \mathbf{x}_{i}\right) \neq \operatorname{Var}\left(c_{i}\right)$.\\
a. Describe the general nature of $\mathrm{E}\left(\mathbf{v}_{i} \mathbf{v}_{i}^{\prime} \mid \mathbf{x}_{i}\right)$.\\
b. What are the asymptotic properties of the random effects estimator and the associated test statistics?\\
c. How should the random effects statistics be modified?\\
10.6. For a model where $\mathbf{x}_{i t}$ varies across $i$ and $t$, define the $K \times K$ symmetric matrices $\mathbf{A}_{1} \equiv \mathrm{E}\left(\Delta \mathbf{X}_{i}^{\prime} \Delta \mathbf{X}_{i}\right)$ and $\mathbf{A}_{2} \equiv \mathrm{E}\left(\ddot{\mathbf{X}}_{i}^{\prime} \ddot{\mathbf{X}}_{i}\right)$, and assume both are positive definite. Define $\hat{\boldsymbol{\theta}} \equiv\left(\hat{\boldsymbol{\beta}}_{F D}^{\prime}, \hat{\boldsymbol{\beta}}_{F E}^{\prime}\right)^{\prime}$ and $\boldsymbol{\theta} \equiv\left(\boldsymbol{\beta}^{\prime}, \boldsymbol{\beta}^{\prime}\right)^{\prime}$, both $2 K \times 1$ vectors.\\
a. Under Assumption FE. 1 (and the rank conditions we have given), find $\sqrt{N}(\hat{\boldsymbol{\theta}}-\boldsymbol{\theta})$ in terms of $\mathbf{A}_{1}, \mathbf{A}_{2}, N^{-1 / 2} \sum_{i=1}^{N} \Delta \mathbf{X}_{i}^{\prime} \Delta \mathbf{u}_{i}$, and $N^{-1 / 2} \sum_{i=1}^{N} \ddot{\mathbf{X}}_{i}^{\prime} \ddot{\mathbf{u}}_{i}$ [with a $\mathrm{o}_{p}(1)$ remainder].\\
b. Explain how to consistently estimate $\mathrm{A} \operatorname{var} \sqrt{N}(\hat{\boldsymbol{\theta}}-\boldsymbol{\theta})$ without further assumptions.\\
c. Use parts a and b to obtain a robust Hausman statistic comparing the FD and FE estimators. What is the limiting distribution of your statistic under $\mathrm{H}_{0}$ ?\\
d. If $\mathbf{x}_{i t}=\left(\mathbf{d}_{t}, \mathbf{w}_{i t}\right)$, where $\mathbf{d}_{t}$ is a $\mathbf{1} \times \boldsymbol{R}$ vector of aggregate time variables, can you compare all of the FD and FD estimates of $\beta$ ? Explain.\\
10.7. Use the two terms of data in GPA.RAW to estimate an unobserved effects version of the model in Example 7.8. You should drop the variable cumgpa (since this variable violates strict exogeneity).\\
a. Estimate the model by RE, and interpret the coefficient on the in-season variable.\\
b. Estimate the model by FE; informally compare the estimates to the RE estimates, in particular that on the in-season effect.\\
c. Construct the nonrobust Hausman test comparing RE and FE. Include all variables in $\mathbf{w}_{i t}$ that have some variation across $i$ and $t$, except for the term dummy.\\
10.8. Use the data in NORWAY.RAW for the years 1972 and 1978 for a two-year panel data analysis. The model is a simple distributed lag model:\\
$\log \left(\right.$ crime $\left._{i t}\right)=\theta_{0}+\theta_{1} d 78_{t}+\beta_{1}$ clrprc $_{i, t-1}+\beta_{2}$ clrprc $_{i, t-2}+c_{i}+u_{i t}$.\\
The variable clrprc is the clear-up percentage (the percentage of crimes solved). The data are stored for two years, with the needed lags given as variables for each year.\\
a. First estimate this equation using a pooled OLS analysis. Comment on the deterrent effect of the clear-up percentage, including interpreting the size of the coefficients. Test for serial correlation in the composite error $v_{i t}$ assuming strict exogeneity (see Section 7.8).\\
b. Estimate the equation by FE, and compare the estimates with the pooled OLS estimates. Is there any reason to test for serial correlation? Obtain heteroskedasticityrobust standard errors for the FE estimates.\\
c. Using FE analysis, test the hypothesis $\mathrm{H}_{0}: \beta_{1}=\beta_{2}$. What do you conclude? If the hypothesis is not rejected, what would be a more parsimonious model? Estimate this model.\\
10.9. Use the data in CORNWELL.RAW for this problem.\\
a. Estimate both an RE and an FE version of the model in Problem 7.11a. Compute the regression-based version of the Hausman test comparing RE and FE.\\
b. Add the wage variables (in logarithmic form), and test for joint significance after estimation by FE.\\
c. Estimate the equation by FD, and comment on any notable changes. Do the standard errors change much between FE and FD ?\\
d. Test the FD equation for $\mathrm{AR}(1)$ serial correlation.\\
10.10. An unobserved effects model explaining current murder rates in terms of the number of executions in the last three years is\\
mrdrte $_{i t}=\theta_{t}+\beta_{1}$ exec $_{i t}+\beta_{2}$ unem $_{i t}+c_{i}+u_{i t}$,\\
where $m r d r t e_{i t}$ is the number of murders in state $i$ during year $t$, per 10,000 people; $\operatorname{exec}_{i t}$ is the total number of executions for the current and prior two years; and $\operatorname{unem}_{i t}$ is the current unemployment rate, included as a control.\\
a. Using the data in MURDER.RAW, estimate this model by FD. Notice that you should allow different year intercepts. Test the errors in the FD equation for serial correlation.\\
b. Estimate the model by FE. Are there any important differences from the FD estimates?\\
c. Under what circumstances would exec $_{\text {it }}$ not be strictly exogenous (conditional on $c_{i}$ )?\\
10.11. Use the data in LOWBIRTH.RAW for this question.\\
a. For 1987 and 1990, consider the state-level equation

$$
\begin{aligned}
\text { lowbrth }_{i t}= & \theta_{1}+\theta_{2} \text { d90 }_{t}+\beta_{1} \text { afdcprc }_{i t}+\beta_{2} \log \left(\text { phypc }_{i t}\right) \\
& +\beta_{3} \log \left(\text { bedspc }_{i t}\right)+\beta_{4} \log \left(\text { pcinc }_{i t}\right)+\beta_{5} \log \left(\text { popul }_{i t}\right)+c_{i}+u_{i t},
\end{aligned}
$$

where the dependent variable is percentage of births that are classified as low birth weight and the key explanatory variable is afdcprc, the percentage of the population in the welfare program, Aid to Families with Dependent Children (AFDC). The other variables, which act as controls for quality of health care and income levels, are physicians per capita, hospital beds per capita, per capita income, and population. Interpretating the equation causally, what sign should each $\beta_{j}$ have? (Note: Participation in AFDC makes poor women eligible for nutritional programs and prenatal care.)\\
b. Estimate the preceding equation by pooled OLS, and discuss the results. You should report the usual standard errors and serial correlation-robust standard errors.\\
c. Difference the equation to eliminate the state $\mathrm{FE}, c_{i}$, and reestimate the equation. Interpret the estimate of $\beta_{1}$ and compare it to the estimate from part b . What do you make of $\hat{\beta}_{2}$ ?\\
d. Add $a f d c p r c^{2}$ to the model, and estimate it by FD. Are the estimates on afdcprc and $a f d c p r c^{2}$ sensible? What is the estimated turning point in the quadratic?\\
10.12. The data in WAGEPAN.RAW are from Vella and Verbeek (1998) for 545 men who worked every year from 1980 to 1987. Consider the wage equation

$$
\begin{aligned}
\log \left(\text { wage }_{i t}\right)= & \theta_{t}+\beta_{1} \text { educ }_{i}+\beta_{2} \text { black }_{i}+\beta_{3} \text { hisp }_{i}+\beta_{4} \text { exper }_{i t} \\
& +\beta_{5} \text { exper }_{i t}^{2}+\beta_{6} \text { married }_{i t}+\beta_{7} \text { union }_{i t}+c_{i}+u_{i t} .
\end{aligned}
$$

The variables are described in the data set. Notice that education does not change over time.\\
a. Estimate this equation by pooled OLS, and report the results in standard form. Are the usual OLS standard errors reliable, even if $c_{i}$ is uncorrelated with all explanatory variables? Explain. Compute appropriate standard errors.\\
b. Estimate the wage equation by RE. Compare your estimates with the pooled OLS estimates.\\
c. Now estimate the equation by FE. Why is exper $_{i t}$ redundant in the model even though it changes over time? What happens to the marriage and union premiums as compared with the RE estimates?\\
d. Now add interactions of the form d81•educ, d82•educ, . . , d87•educ and estimate the equation by FE. Has the return to education increased over time?\\
e. Return to the original model estimated by FE in part c. Add a lead of the union variable, union $_{i, t+1}$ to the equation, and estimate the model by FE (note that you lose the data for 1987). Is union $_{i, t+1}$ significant? What does your finding say about strict exogeneity of union membership?\\
f. Return to the original model, but add the interactions black $_{i} \cdot$ union $_{i t}$ and $h_{i s p} p_{i} \cdot$ union $_{i t}$. Do the union wage premiums differ by race? Obtain the usual FE statistics and those fully robust to heteroskedasticity and serial correlation.\\
g. Add union $_{i, t+1}$ to the equation from part f , and obtain a fully robust test of the hypothesis that $\left\{\right.$ union $\left._{i t}: t=1, \ldots, T\right\}$ is strictly exogenous. What do you conclude?\\
10.13. Consider the standard linear unobserved effects model (10.11), under the assumptions\\
$\mathrm{E}\left(u_{i t} \mid \mathbf{x}_{i}, \mathbf{h}_{i}, c_{i}\right)=0, \operatorname{Var}\left(u_{i t} \mid \mathbf{x}_{i}, \mathbf{h}_{i}, c_{i}\right)=\sigma_{u}^{2} h_{i t}, \quad t=1, \ldots, T$,\\
where $\mathbf{h}_{i}=\left(h_{i 1}, \ldots, h_{i T}\right)$. In other words, the errors display heteroskedasticity that depends on $h_{i t}$. (In the leading case, $h_{i t}$ is a function of $\mathbf{x}_{i t}$.) Suppose you estimate $\boldsymbol{\beta}$ by minimizing the weighted sum of squared residuals\\
$\sum_{i=1}^{N} \sum_{t=1}^{T}\left(y_{i t}-a_{1} d 1_{i}-\cdots-a_{N} d N_{i}-\mathbf{x}_{i t} \mathbf{b}\right)^{2} / h_{i t}$\\
with respect to the $a_{i}, i=1, \ldots, N$ and $\mathbf{b}$, where $d n_{i}=1$ if $i=n$. (This would seem to be the natural analogue of the dummy variable regression, modified for known heteroskedasticity. We might call this a fixed effects weighted least squares estimator.)\\
a. Show that the FEWLS estimator is generally consistent for fixed $T$ as $N \rightarrow \infty$. Do you need the variance function to be correctly specified in the sense that $\operatorname{Var}\left(u_{i t} \mid \mathbf{x}_{i}, \mathbf{h}_{i}, c_{i}\right)=\sigma_{u}^{2} h_{i t}, t=1, \ldots, T$ ? Explain.\\
b. Suppose the variance function is correctly specified and $\operatorname{Cov}\left(u_{i t}, u_{i s} \mid \mathbf{x}_{i}, \mathbf{h}_{i}, c_{i}\right)=0$, $t \neq s$. Find the asymptotic variance of $\sqrt{N}\left(\hat{\boldsymbol{\beta}}_{\text {FEWLS }}-\boldsymbol{\beta}\right)$.\\
c. Under the assumptions of part b , how would you estimate $\sigma_{u}^{2}$ and $\mathrm{Avar}\left[\sqrt{N}\left(\hat{\boldsymbol{\beta}}_{\text {FEWLS }}-\boldsymbol{\beta}\right)\right]$ ?\\
d. If the variance function is misspecified, or there is serial correlation in $u_{i t}$, or both, how would you estimate $\operatorname{Avar}\left[\sqrt{N}\left(\hat{\boldsymbol{\beta}}_{\text {FEWLS }}-\boldsymbol{\beta}\right)\right]$ ?\\
10.14. Suppose that we have the unobserved effects model\\
$y_{i t}=\alpha+\mathbf{x}_{i t} \boldsymbol{\beta}+\mathbf{z}_{i} \gamma+h_{i}+u_{i t}$\\
where the $\mathbf{x}_{i t}(1 \times K)$ are time-varying, the $\mathbf{z}_{i}(1 \times M)$ are time-constant, $\mathrm{E}\left(u_{i t} \mid \mathbf{x}_{i}, \mathbf{z}_{i}, h_{i}\right)=0, t=1, \ldots, T$, and $\mathrm{E}\left(h_{i} \mid \mathbf{x}_{i}, \mathbf{z}_{i}\right)=0$. Let $\sigma_{h}^{2}=\operatorname{Var}\left(h_{i}\right)$ and $\sigma_{u}^{2}= \operatorname{Var}\left(u_{i t}\right)$. If we estimate $\boldsymbol{\beta}$ by fixed effects, we are estimating the equation $y_{i t}=\mathbf{x}_{i t} \boldsymbol{\beta}+c_{i}+u_{i t}$, where $c_{i}=\alpha+\mathbf{z}_{i} \gamma+h_{i}$.\\
a. Find $\sigma_{c}^{2} \equiv \operatorname{Var}\left(c_{i}\right)$. Show that $\sigma_{c}^{2}$ is at least as large as $\sigma_{h}^{2}$, and usually strictly larger.\\
b. Explain why estimation of the model by fixed effects will lead to a larger estimated variance of the unobserved effect than if we estimate the model by random effects. Does this result make intuitive sense?\\
c. If $\lambda_{c}$ is the quasi-time-demeaning parameter without $\mathbf{z}_{i}$ in the model and $\lambda_{h}$ is the quasi-time-demeaning parameter with $\mathbf{z}_{i}$ in the model, show that $\lambda_{c} \geq \lambda_{h}$, with strict inequality if $\boldsymbol{\gamma} \neq \mathbf{0}$.\\
d. What does part c imply about using pooled OLS versus FE as the first step estimator for estimating the variance of the unobserved heterogeneity in RE estimation?\\
e. Suppose that, in addition to RE.1-RE. 3 holding in the orginal model, $\mathrm{E}\left(\mathbf{z}_{i} \mid \mathbf{x}_{i}\right)=\mathbf{0}, t=1, \ldots, T$ and $\operatorname{Var}\left(\mathbf{z}_{i} \mid \mathbf{x}_{i}\right)=\operatorname{Var}\left(\mathbf{z}_{i}\right)$. Show directly-that is, by comparing the two asymptotic variances-that the RE estimator that includes $\mathbf{z}_{i}$ is asymptotically more efficient than the RE estimator that excludes $\mathbf{z}_{i}$. (The result also follows from Problem 7.15 without the assumption $\operatorname{Var}\left(c_{i} \mid \mathbf{x}_{i}\right)=\operatorname{Var}\left(c_{i}\right)$; in fact, we only need to assume $\mathbf{z}_{i}$ is uncorrelated with $\mathbf{x}_{i}$.)\\
10.15. Consider the standard unobserved effects model, first under a stronger version of the RE assumptions. Let $v_{i t}=c_{i}+u_{i t}, t=1, \ldots, T$, be the composite errors,\\
as usual. Then, in addition to RE.1-RE.3, assume that the $T \times 1$ vector $\mathbf{v}_{i}$ is independent of $\mathbf{x}_{i}$ and that all conditional expectations involving $\mathbf{v}_{i}$ are linear.\\
a . Let $\bar{v}_{i}$ be the time average of the $v_{i t}$. Argue that $E\left(v_{i t} \mid \mathbf{x}_{i}, \bar{v}_{i}\right)=E\left(v_{i t} \mid \bar{v}_{i}\right)=\bar{v}_{i}$. (Hint: To show the second equality, recall that, because the expectation is assumed to be linear, the coefficient on $\bar{v}_{i}$ is $\operatorname{Cov}\left(v_{i t}, \bar{v}_{i}\right) / \operatorname{Var}\left(\bar{v}_{i}\right)$.)\\
b. Use part a to show that\\
$\mathrm{E}\left(y_{i t} \mid \mathbf{x}_{i}, \bar{y}_{i}\right)=\mathbf{x}_{i t} \boldsymbol{\beta}+\left(\bar{y}_{i}-\overline{\mathbf{x}} \boldsymbol{\beta}\right)$.\\
c. Argue that estimation of $\boldsymbol{\beta}$ based on part b leads to the FE estimator.\\
d. How can you reconcile part d with the fact that the full RE assumptions have been assumed? (Hint: What is the conditional expectation underlying RE estimation?)\\
e. Show that the arguments leading up to part c carry over using linear projections under RE.1-RE. 3 alone.\\
10.16. Assume that through time periods $T+1$, the assumptions in Problem 10.15 hold. Suppose you want to forecast $y_{i, T+1}$ at time $T$, where you know $\mathbf{x}_{i, T+1}$ in addition to all past values on $\mathbf{x}$ and $y$.\\
a. It can be shown, under the given assumptions, that $E\left(v_{i, T+1} \mid \mathbf{x}_{i}, \mathbf{x}_{i, T+1}, v_{i 1}, \ldots, v_{i T}\right) =E\left(v_{i, T+1} \mid \bar{v}_{i}\right)$. Use this to show that $E\left(v_{i, T+1} \mid \bar{v}_{i}\right)=\left[\sigma_{c}^{2} /\left(\sigma_{c}^{2}+\sigma_{u}^{2} / T\right)\right] \bar{v}_{i}$.\\
b. Use part a to derive $E\left(y_{i, T+1} \mid \mathbf{x}_{i 1}, \ldots, \mathbf{x}_{i T}, \mathbf{x}_{i, T+1}, y_{i 1}, \ldots, y_{i T}\right)$.\\
c. Compare the expectation in part b with $E\left(y_{i, T+1} \mid \mathbf{x}_{i 1}, \ldots, \mathbf{x}_{i T}, \mathbf{x}_{i, T+1}\right)$.\\
d. Which of the two expectations in parts $b$ and $c$ leads to a better forecast as defined by smallest variance of the forecast error?\\
e. How would you forecast $y_{i, T+1}$ if you have to estimate all unknown population parameters?\\
10.17. Consider a standard unobserved effects model but where we explicitly separate out aggregate time effects, say $\mathbf{d}_{t}$, a $1 \times R$ vector, where $R \leq T-1$. (These are usually a full set of time period dummies, but they could be other aggregate time variables, such as specific functions of time.) Therefore, the model is\\
$y_{i t}=\alpha+\mathbf{d}_{t} \boldsymbol{\eta}+\mathbf{w}_{i t} \boldsymbol{\delta}+c_{i}+u_{i t}, \quad t=1, \ldots, T$,\\
where $\mathbf{w}_{i t}$ is the $1 \times M$ vector of explanatory variables that vary across $i$ and $t$. Because the $\mathbf{d}_{t}$ do not change across $i$, we take them to be nonrandom. Because we have included an intercept in the model, we can assume that $E\left(c_{i}\right)=0$. Let\\
$\lambda=1-\left\{1 /\left[1+T\left(\sigma_{c}^{2} / \sigma_{u}^{2}\right)\right\}^{1 / 2}\right.$ be the usual quasi-time-demeaning parameter for RE estimation. In what follows, we take $\lambda$ as known because estimating it does not affect the asymptotic distribution results.\\
a. Show that we can write the quasi-time-demeaned equation for RE estimation as\\
$y_{i t}-\lambda \bar{y}_{i}=\mu+\left(\mathbf{d}_{t}-\overline{\mathbf{d}}\right) \boldsymbol{\eta}+\left(\mathbf{w}_{i t}-\lambda \overline{\mathbf{w}}_{i}\right) \boldsymbol{\delta}+\left(v_{i t}-\lambda \bar{v}_{i}\right)$,\\
where $\mu=(1-\lambda) \alpha+(1-\lambda) \overline{\mathbf{d}} \boldsymbol{\eta}, v_{i t}=c_{i}+u_{i t}$, and $\overline{\mathbf{d}}=T^{-1} \sum_{t=1}^{T} \mathbf{d}_{t}$ is nonrandom.\\
b. To simplify the algebra without changing the substance of the findings, assume that $\mu=0$ and that we exclude an intercept in estimating the quasi-time-demeaned equation. Write $\mathbf{g}_{i t}=\left(\mathbf{d}_{t}-\overline{\mathbf{d}}, \mathbf{w}_{i t}-\lambda \overline{\mathbf{w}}_{i}\right)$ and $\boldsymbol{\beta}=\left(\boldsymbol{\eta}^{\prime}, \boldsymbol{\delta}^{\prime}\right)^{\prime}$. We will study the asymptotic distribution of the RE estimator by using the pooled OLS estimator from $y_{i t}-\lambda \bar{y}_{i}$ on $\mathbf{g}_{i t}, t=1, \ldots, T ; i=1, \ldots, N$. Show that under Assumptions RE. 1 and RE.2,

$$
\sqrt{N}\left(\hat{\boldsymbol{\beta}}_{R E}-\boldsymbol{\beta}\right)=\mathbf{A}_{1}^{-1} N^{-1 / 2} \sum_{i=1}^{N} \sum_{t=1}^{T} \mathbf{g}_{i t}^{\prime}\left(v_{i t}-\lambda \bar{v}_{i}\right)+o_{p}(1),
$$

where $\mathbf{A}_{1} \equiv \sum_{t=1}^{T} \mathrm{E}\left(\mathbf{g}_{i t}^{\prime} \mathbf{g}_{i t}\right)$. Further, verify that for any $i$,

$$
\sum_{t=1}^{T}\left(\mathbf{d}_{t}-\overline{\mathbf{d}}\right)\left(v_{i t}-\lambda \bar{v}_{i}\right)=\sum_{t=1}^{T}\left(\mathbf{d}_{t}-\overline{\mathbf{d}}\right) u_{i t}
$$

c. Show that under FE. 1 and FE.2,

$$
\sqrt{N}\left(\hat{\boldsymbol{\beta}}_{F E}-\boldsymbol{\beta}\right)=\mathbf{A}_{2}^{-1} N^{-1 / 2} \sum_{i=1}^{N} \sum_{t=1}^{T} \mathbf{h}_{i t}^{\prime} u_{i t}+o_{p}(1),
$$

where $\mathbf{h}_{i t} \equiv\left(\mathbf{d}_{t}-\overline{\mathbf{d}}, \mathbf{w}_{i t}-\overline{\mathbf{w}}_{i}\right)$ and $\mathbf{A}_{2} \equiv \sum_{t=1}^{T} \mathrm{E}\left(\mathbf{h}_{i t}^{\prime} \mathbf{h}_{i t}\right)$.\\
d. Under RE.1, FE.1, and FE.2, show that $\mathbf{A}_{1} \sqrt{N}\left(\hat{\boldsymbol{\beta}}_{R E}-\boldsymbol{\beta}\right)-\mathbf{A}_{2} \sqrt{N}\left(\hat{\boldsymbol{\beta}}_{F E}-\boldsymbol{\beta}\right)$ has an asymptotic variance matrix of rank $M$ rather than $R+M$.\\
e. What implications does part d have for a Hausman test that compares FE and RE when the model contains aggregate time variables of any sort? Does it matter whether Assumption RE. 3 holds under the null?\\
10.18. Use the data in WAGEPAN.RAW to answer this question.\\
a. Using lwage as the dependent variable, estimate a model that contains an intercept and the year dummies $d 81$ through $d 87$. Used pooled OLS, RE, FE, and FD (where\\
in the latter case you difference the year dummies, along with lwage, and omit an overall constant in the FD regression). What do you conclude about the coefficients on the dummy variables?\\
b. Add the time-constant variables educ, black, and hisp to the model, and estimate it by POLS and RE. How do the coefficients compare? What happens if you estimate the equation by FE?\\
c. Are the POLS and RE standard errors from part b the same? Which ones are probably more reliable?\\
d. Obtain the fully robust standard errors for POLS. Do you prefer these or the usual RE standard errors?\\
e. Obtain the fully robust RE standard errors. How do these compare with the fully robust POLS standard errors, and why?

\section*{11 More Topics in Linear Unobserved Effects Models}

\end{document}
